% Shared source selected by the CCL-1.tex and CCL-2.tex entry points.
\documentclass[10pt]{article}

\usepackage[letterpaper,margin=1in]{geometry}
\usepackage[T1]{fontenc}
\usepackage{lmodern}
\usepackage{microtype}
\usepackage{amsmath,amssymb,amsthm,mathtools}
\usepackage{booktabs,array}
\usepackage{enumitem}
\usepackage[hidelinks]{hyperref}
\ifCCLone
\hypersetup{
  pdftitle={CCL-1: Finite Encodings and Shrinking Cylinders over the Gaussian Rationals},
  pdfauthor={Anonymous}
}
\else
\hypersetup{
  pdftitle={CCL-2: Arithmetic and Finite Calculus with Complex Continued Logarithms},
  pdfauthor={Anonymous}
}
\fi

\setlist{nosep,leftmargin=2em}
\allowdisplaybreaks

\newtheorem{theorem}{Theorem}[section]
\newtheorem{lemma}[theorem]{Lemma}
\newtheorem{proposition}[theorem]{Proposition}
\newtheorem{corollary}[theorem]{Corollary}
\theoremstyle{definition}
\newtheorem{definition}[theorem]{Definition}
\newtheorem{algorithm}[theorem]{Algorithm}
\newtheorem{example}[theorem]{Example}
\theoremstyle{remark}
\newtheorem{remark}[theorem]{Remark}

\newcommand{\Q}{\mathbb{Q}}
\newcommand{\Z}{\mathbb{Z}}
\newcommand{\N}{\mathbb{N}}
\newcommand{\Qi}{\mathbb{Q}(i)}
\newcommand{\Zi}{\mathbb{Z}[i]}
\newcommand{\Norm}{\operatorname{N}}
\newcommand{\Ree}{\operatorname{Re}}
\newcommand{\Imm}{\operatorname{Im}}
\newcommand{\sqfloor}[1]{\left\lfloor #1\right\rfloor_{\!\square}}
\newcommand{\bits}{\{0,1\}}
\newcommand{\pref}{\preccurlyeq}
\newcommand{\Cyl}{\mathsf{Cyl}}
\newcommand{\code}{\beta}
\newcommand{\leanref}[1]{\par\smallskip\noindent
  \emph{Formal verification:}\par\nobreak
  {\raggedright\nolinkurl{#1}.\par}}
\newcommand{\term}{\dagger}
\newcommand{\Proj}[2]{[#1:#2]}
\newcommand{\mat}[4]{\begin{pmatrix}#1&#2\\#3&#4\end{pmatrix}}

\ifCCLone
\title{CCL-1: Finite Encodings and Shrinking Cylinders over the Gaussian Rationals}
\else
\title{CCL-2: Arithmetic and Finite Calculus with Complex Continued Logarithms}
\fi
\author{Anonymous}
\date{August 14, 2026}

\begin{document}
\maketitle

\begin{abstract}
\ifCCLone
Continued fractions and Gosper's continued logarithms repeatedly remove a coarse integral or dyadic part and then apply a homographic inverse map.  We give two analogous finite algorithms over the Gaussian rationals.  The first is a coordinate-floor complex continued fraction
\[
 z_n=t_n+\frac{1+i}{z_{n+1}},\qquad t_n\in\Zi.
\]
Its termination proof is reduced to an explicit finite eight-digit contraction certificate; no unrecorded computer-algebra calculation is used.  The second algorithm, the complex continued logarithm (CCL), emits raw binary words.  It normalizes the Gaussian norm by powers of $1+i$, emits three bits selecting one of the eight representatives
\[
 \pm1,\quad \pm i,\quad \pm(1+i),\quad \pm(1-i),
\]
and then inverts the contracted remainder.  A strict post-initial normalization convention gives direct denominator descent after every non-terminal sector step.

All choices are bounded searches through explicitly finite lists.  Geometric regions are used only as rational predicates, not as named infinite sets.  Complete CCL blocks may be abbreviated by $[k,s]$, but an arbitrary prefix can end inside the unary scale run or inside the three-bit sector code, and those unresolved bits remain part of its cylinder predicate and bit cost.  The canonical truncation uses the formal projective coordinate $[1:0]$; no completed real or complex point and no actual infinite stream is required.  A finite two-witness theorem gives the corresponding bit-prefix minimax statement without suprema over infinite cylinders.

The three angular bits are also parsed incrementally.  Their first two bits are a quarter-turn frame and the last bit chooses the axial or diagonal representative; unit folds multiply only the first decoder column and preserve the represented projective point.  A two-tail identity for Hermitian projective quadrance then shows where genuine numerical contraction occurs.  A completed block of radial norm $L$ contracts the Hermitian projective quadrance between every pair of compatible projective tails by the rational factor
\[
 q(L)=\frac{2L}{(1+L)^2}\le\frac12.
\]
Consequently a prefix containing $n$ completed blocks has the finite pairwise certificate $\chi^2\le2^{-n}$.  Individual raw bits still refine the rational predicate, but strict contraction at every bit is neither required nor true.

We also give an operational semantics for values not returned by a finite rational encoding.  A boxed prefix certificate consists only of a finite word, a rational axis-aligned box, a rational witness, and finite proofs that the prefix predicate is satisfiable and that every rational realization of it lies in the box.  The completed-block contraction theorem converts any fixed affine norm bound and rational witness into an explicit shrinking rational box.  A progressive name answers each requested accuracy by extending the prefix and returning a nested box of arbitrarily small rational squared diagonal.  Thus every observation remains finite, while successive prefixes explicitly produce better enclosures.  For a finite rational name, its initial substrings already form a finite tower of prefix-minimax approximations ending at the exact value.

These finite cylinders and boxed names are the representation-level input to
CCL-2.  That companion paper supplies the redundant-output machinery needed
to couple them persistently to RFAP zero/pole tracks and then develops
arithmetic and finite calculus.  Keeping that layer separate leaves the
canonical encoder, parser, and contraction results here independent of any
particular operation on names.
\else
CCL-1 constructs finite complex continued-logarithm words, ordinary
raw-prefix parser predicates, shrinking projective cylinders, and progressive
rational boxed names over the Gaussian rationals.  This second paper develops
the arithmetic layer on that representation.  Its redundant output sectors
overlap, so a whole rational enclosure can justify a next digit without
deciding an unstable canonical boundary.  A rational squared-error sidecar
transports addition, subtraction, multiplication, and apartness-certified
division.  Projective unary padding handles exact results that cannot be
recognized from opaque non-terminal inputs, and a persistent finite-request
cursor returns every requested number of ordinary CCL bits.

The same interface computes exponential, positive-real logarithm, and complex
logarithm in an explicit branch chart.  All local formulae are finite rational
Horner recurrences with displayed error bounds.  Global range preparation is
a bounded dyadic calculation followed by finitely many arithmetic operations.
Operational equivalence proves the field laws, exponential addition,
exponential--logarithm inverse laws on certified charts, and congruence under
independent refinement schedules.

Finally, finite divided differences give limit-free differential laws,
first- and higher-order jets, Taylor error bounds, coefficient-defined and
progressive path integration, trigonometric compositions, projective
rational-function evaluation, finite families of fractional powers, and an
explicit coarse bit-complexity account.  Every result is a terminating
request over displayed finite data.  No completed stream, completed scalar
field, infinitesimal, limit, or function represented as an infinite graph is
formed.
\fi
\end{abstract}

\ifCCLone
\section{Introduction and constructive conventions}

The ordinary continued fraction is the Euclidean algorithm presented as a homographic digit expansion.  Gosper's continued logarithm replaces unbounded positive partial quotients by unary binary scale runs and a fixed subtraction branch; its arithmetic interpretation goes back to HAKMEM~\cite{hakmem}, and its dynamics and complexity have since been studied in several forms~\cite{borwein-cl,borwein-generalized,collins,shallit}.  Over the Gaussian rationals, Hurwitz's nearest-Gaussian-integer construction is the classical point of comparison~\cite{hurwitz}; modern accounts develop its geometry and dynamics~\cite{hockman,kirsebom,gonzalez}.

The coordinate-floor CCF below is not proposed as a replacement for Hurwitz's theory.  It is the direct Cartesian analogue of the real floor construction and provides a useful contrast with the radial--angular CCL.  This revision is organized around three constructive constraints: every search is bounded by an explicit finite list, geometric information is carried by rational predicates rather than by names for point collections, and every direct observation consists of finite bit words and finite rational certificates.  A value need not be returned as a completed point: it may instead be presented by a procedure that extends its current finite prefix and enclosure to any requested rational accuracy.

Two different approximation statements are kept separate.  Prefix minimax optimality is an information theorem about finite tags.  Numerical enclosure is supplied by a contraction theorem for Hermitian projective quadrance: although a partial unary or angular bit may only change coordinates and refine a predicate, each completed CCL block contracts the Hermitian projective quadrance between compatible tails to at most one half.  On a fixed bounded affine region this gives explicit rational boxes.  Every construction is consequently a finite algorithm over displayed rational data.

Every numerical object in the algorithms is rational.  A Gaussian rational is a pair of rationals $x,y$, written $x+iy$; a Gaussian integer is a pair of integers $a,b$, written $a+bi$.  We use $\Qi$ and $\Zi$ for these two ambient types and put
\[
 \Norm(x+iy)=x^2+y^2.
\]
Whenever a square, annulus, sector, or cylinder is mentioned below, it is introduced as a decidable predicate on finite rational data.  Every named collection used in an argument is an explicitly finite list, such as the eight sector representatives or a bounded list of natural exponents.

No completed real or complex numbers are needed.  In particular, the paper does not assign a completed value to an actual infinite stream.  A finite stream is a finite list of bits, and a finite rational input has a finite output.  Section~7 also permits a non-terminating refinement process, but each invocation of that process returns only a finite extension, a rational box, and finite correctness evidence.  The name is this finite-precision interface, not a pre-existing completed point or a named infinite intersection.  The symbol traditionally written as infinity appears only through the formal projective coordinate
\[
 \omega=\Proj{1}{0}.
\]
It is a pair of rational coordinates, not a completed numerical value.

Natural numbers include zero.  We put
\[
 \gamma=1+i,\qquad \Norm(\gamma)=2.
\]
Direct expansion gives the multiplicative identity $\Norm(zw)=\Norm(z)\Norm(w)$ for Gaussian rationals.  In particular,
\[
 \Norm(\gamma^kz)=2^k\Norm(z)
\]
for every $k\in\N$.

All floor operations, comparisons, finite-list searches, and matrix products used below are computable.  Every phrase of the form ``obtain the exponent'' refers to the output of the displayed bounded scan together with its proof.  The sector table is indexed by three bits, the exceptional CCF table by two finite eight-valued indices, and both termination theorems recurse on an explicitly displayed natural-number norm.

\section{Projective homographic bookkeeping}

A projective rational coordinate is written $\Proj{p}{q}$, where $p,q\in\Qi$ are not both zero; non-zero scalar multiples are identified.  A matrix
\[
 M=\mat{a}{b}{c}{d}
\]
acts by the finite formula
\[
 M\Proj{p}{q}=\Proj{ap+bq}{cp+dq}
\]
whenever the output coordinates are not both zero.  Every elementary decoder matrix below has non-zero determinant, so this condition is automatic for its finite products.  For an affine input $z$ we use $\Proj{z}{1}$.  If $M$ is a finite decoder product, its first column is non-zero and its canonical projective truncation is
\[
 M\omega=\Proj{a}{c}.
\]
This formula remains meaningful when $c=0$ and therefore avoids a division convention.

\begin{lemma}[continued-fraction substitution]\label{lem:cf-substitution}
Let $K$ be a commutative field and let $t,\mu,y\in K$.  Assume that $ty+\mu$ and $y$ are not both zero, and put
\[
 A(t,\mu)=\mat{t}{\mu}{1}{0}.
\]
Then
\[
 A(t,\mu)\Proj{y}{1}=\Proj{ty+\mu}{y}.
\]
If $y\ne0$, this is the affine identity $x=t+\mu/y$.  Moreover,
\[
 \mat{a}{b}{c}{d}A(t,\mu)
 =\mat{at+b}{a\mu}{ct+d}{c\mu}.
\]
\end{lemma}

\begin{proof}
Both statements are direct multiplication.
\end{proof}

We also use the finite matrices
\[
 S_\lambda=\mat{\lambda}{0}{0}{1},\qquad
 T_s=\mat{1}{s}{0}{1},\qquad
 I_\lambda=\mat{0}{1}{\lambda}{0},
\]
representing $z\mapsto\lambda z$, $z\mapsto z+s$, and $z\mapsto1/(\lambda z)$ whenever the affine expressions are defined.

\section{Bounded normalization by powers of two}

The algorithms repeatedly choose a power of two surrounding a positive rational.  It is important that this choice be a bounded finite search rather than an appeal to a logarithm or to unbounded minimization.

\begin{lemma}[finite binary normalization]\label{lem:finite-normalization}
Let $A,B\in\N$ be positive.
\begin{enumerate}[label=\textup{(\alph*)}]
\item If $2B\le A$, scanning the finite list $1,2,\ldots,A$ and retaining the last $k$ satisfying $2^kB\le A$ returns the unique $k$ with
\[
 2^kB\le A<2^{k+1}B.
\]
\item If $A<2B$, scanning the finite list $0,1,\ldots,B$ and returning the first $m$ satisfying $B\le2^mA$ returns the unique $m$ with
\[
 B\le2^mA<2B.
\]
\item If $B<A$, scanning the finite list $0,1,\ldots,A$ and retaining the last $k$ satisfying $2^kB<A$ returns the unique $k$ with
\[
 2^kB<A\le2^{k+1}B.
\]
\end{enumerate}
\end{lemma}

\begin{proof}
The elementary induction $2^n\ge n+1$ supplies all required bounds.  In (a), $k=1$ is admissible and $k=A$ is not, since $2^AB>A$.  Maximality gives the upper inequality.  In (b), $2^BA\ge B$, so the finite scan succeeds.  If the first successful index is positive, the previous index fails and doubling gives the strict upper inequality; if it is zero, the assumed inequality $A<2B$ gives it.  In (c), $k=0$ is admissible and $k=A$ is not.  Maximality again gives the upper inequality.  Strict monotonicity of powers of two gives uniqueness in all three cases.
\end{proof}

For $q=A/B>0$, parts (a) and (b) construct the exponent placing $q$ in $[1,2)$, while part (c) constructs the exponent placing a rational $q>1$ in $(1,2]$.

\section{Two real baseline algorithms}

\subsection{Classical continued fractions}

For $x\in\Q$, put $a=\lfloor x\rfloor$.  If $x=a$, stop; otherwise replace $x$ by $1/(x-a)$.  Thus
\[
 x=a+\frac1{x^+},
\]
and the decoder matrix update is
\[
 M\longmapsto M\mat{a}{1}{1}{0}.
\]

\begin{proposition}\label{prop:real-cf-finite}
The classical algorithm terminates on every rational input.  If it emits $a_0,\ldots,a_n$ and $M=A(a_0,1)\cdots A(a_n,1)$, then $M\omega=\Proj{x}{1}$ for the original input $x$.
\end{proposition}

\begin{proof}
Write a non-integral complete quotient as $p/q$ with $q>0$.  The next denominator is the Euclidean remainder $p-aq$, a non-negative integer strictly smaller than $q$.  This gives a well-founded recursion on $q$.  The reconstruction identity follows by composing the inverse-branch equations and using the terminal equality $x_n=a_n$, equivalently $A(a_n,1)\omega=\Proj{a_n}{1}$.
\end{proof}

\subsection{Real continued logarithms}

For this subsection $x\in\Q$ and $x\ge1$.  If $x=1$, stop.  If $x\ge2$, emit $1$ and replace $x$ by $x/2$.  If $1<x<2$, emit $0$ and replace $x$ by $1/(x-1)$.  The inverse branches give
\[
 1:\quad M\longmapsto M\mat{2}{0}{0}{1},
 \qquad
 0:\quad M\longmapsto M\mat{1}{1}{1}{0}.
\]

\begin{proposition}\label{prop:cl-finite}
For $p,q\in\N$ with $p\ge q>0$, the continued-logarithm algorithm terminates on $p/q$.  Composing the inverse-branch matrices of its emitted word and applying the product to the terminal coordinate $\Proj{1}{1}$ reconstructs $\Proj{p}{q}$.
\end{proposition}

\begin{proof}
Use the finite normalization of \autoref{lem:finite-normalization} to consolidate a run of $k$ ones.  Then
\[
 1\le\frac{p}{2^kq}<2.
\]
If equality holds on the left, the algorithm stops after the run.  Otherwise the following zero bit sends
\[
 \frac pq\longmapsto\frac{2^kq}{p-2^kq}.
\]
The new numerator $2^kq$ is a positive integer strictly smaller than $p$.  Grouping one unary run and its following zero into a macro-step gives well-founded recursion on the numerator.  The reconstruction statement follows by composing the displayed inverse branches back from the terminal state $1$.
\end{proof}

\leanref{CCL.BaselineAlgorithms.cfEncode_reconstruct;
CCL.BaselineAlgorithms.clEncode_reconstruct}

\section{A floor-rotated complex continued fraction}

For $z=x+iy\in\Qi$, define the coordinate floor
\[
 \sqfloor{z}=\lfloor x\rfloor+i\lfloor y\rfloor.
\]

\begin{algorithm}[CCF expansion]\label{alg:ccf}
Given $z_0\in\Qi$, put $t_n=\sqfloor{z_n}$.  If $z_n=t_n$, stop.  Otherwise put
\[
 z_{n+1}=\frac{\gamma}{z_n-t_n}.
\]
\end{algorithm}

Equivalently,
\[
 z_n=t_n+\frac{\gamma}{z_{n+1}},
\]
so \autoref{lem:cf-substitution} gives
\[
 \mat{a}{b}{c}{d}
 \longmapsto
 \mat{at_n+b}{a\gamma}{ct_n+d}{c\gamma}.
\]

Define the square predicate
\[
 \operatorname{Sq}(u)\;:\Longleftrightarrow\;
 0\le\Ree u<1\ \text{ and }\ 0\le\Imm u<1.
\]
For non-zero $u$ satisfying $\operatorname{Sq}(u)$, put
\[
 F(u)=\frac{\gamma}{u}-\sqfloor{\frac{\gamma}{u}}.
\]
Then $\operatorname{Sq}(F(u))$.

\begin{lemma}[digit cone]\label{lem:digit-cone}
Let $u\ne0$ satisfy $\operatorname{Sq}(u)$ and write
\[
 \sqfloor{\frac{\gamma}{u}}=a+bi.
\]
Then
\[
 a\ge1,
 \qquad
 -a-1\le b\le a.
\]
\end{lemma}

\begin{proof}
Write $u=x+iy$.  Since
\[
 \frac{\gamma}{u}
 =\frac{x+y}{x^2+y^2}
  +i\frac{x-y}{x^2+y^2},
\]
we have $\Ree(\gamma/u)\ge|\Imm(\gamma/u)|$.  Also
\[
 x+y-(x^2+y^2)=x(1-x)+y(1-y)>0,
\]
so $\Ree(\gamma/u)>1$ and hence $a\ge1$.  The upper bound $b\le a$ follows from $\Imm(\gamma/u)\le\Ree(\gamma/u)<a+1$.  Finally,
\[
 -b-1<-\Imm(\gamma/u)\le\Ree(\gamma/u)<a+1,
\]
and integrality gives $b\ge-a-1$.
\end{proof}

\leanref{CCL.CoordinateFloor.digit_cone}

The next lemma replaces a quantifier-elimination certificate by a finite table that translates directly into case splits and polynomial normalization.

\begin{lemma}[three-step contraction]\label{lem:ccf-contraction}
Let $u\ne0$ satisfy $\operatorname{Sq}(u)$.  If $v=F(u)\ne0$ and $w=F(v)$, then
\[
 \Norm(uvw)<1.
\]
\end{lemma}

\begin{proof}
Write
\[
 t=\sqfloor{\frac\gamma u},\qquad
 t'=\sqfloor{\frac\gamma v},
\]
so that
\[
 u=\frac\gamma{t+v},\qquad
 v=\frac\gamma{t'+w}.
\]
If $t=a+bi$, define the integer
\[
 \lambda(t)=
 \begin{cases}
 a^2+b^2,&b\ge0,\\
 a^2,&b=-1,\\
 a^2+(b+1)^2,&b\le-2.
 \end{cases}
\]
For every $q$ satisfying $\operatorname{Sq}(q)$, direct coordinate bounds give
\[
 \Norm(t+q)\ge\lambda(t).
\]
Since $\Norm(w)<2$,
\[
 \Norm(uvw)
 =\frac{4\Norm(w)}{\Norm(t+v)\Norm(t'+w)}
 <\frac8{\lambda(t)\lambda(t')}.
\]
Thus the claim is immediate when $\lambda(t)\lambda(t')\ge8$.

By \autoref{lem:digit-cone}, the inequality $\lambda(t)<8$ leaves only the following finite list:
\[
 \begin{array}{c|cccccccc}
 r&0&1&2&3&4&5&6&7\\ \hline
 \tau_r&1-2i&1-i&1&1+i&2-2i&2-i&2&2+i\\
 \lambda(\tau_r)&2&1&1&2&5&4&4&5.
 \end{array}
\]
Suppose now that $t=\tau_r$, $t'=\tau_c$, and $\lambda(t)\lambda(t')<8$.  Write $w=x+iy$ and $tt'+\gamma=C+Di$.  The identity
\[
 (t+v)(t'+w)=t(t'+w)+\gamma
\]
gives
\[
 \Norm(uvw)=\frac{4\Norm(w)}{\Norm(t(t'+w)+\gamma)}.
\]
It is therefore enough to prove positivity of
\[
 \begin{aligned}
 \Delta_{r,c}(x,y)
 &=\Norm(t(t'+w)+\gamma)-4\Norm(w)\\
 &=(\Norm(t)-4)(x^2+y^2)
   +2(aC+bD)x+2(-bC+aD)y+C^2+D^2,
 \end{aligned}
\]
where $t=a+bi$.  The following finite table records lower bounds valid whenever $0\le x\le1$ and $0\le y\le1$.  A dash means that the preceding coarse estimate already applies.  The cross marks a digit pair that cannot occur.
\[
\renewcommand{\arraystretch}{1.15}
\begin{array}{c|rrrrrrrr}
 r\backslash c&0&1&2&3&4&5&6&7\\ \hline
0&0&1&5&16&-&-&-&-\\
1&\times&0&4&10&4&6&10&16\\
2&0&1&4&8&5&6&9&13\\
3&6&4&6&10&-&-&-&-\\
4&-&6&10&-&-&-&-&-\\
5&-&5&9&-&-&-&-&-\\
6&-&6&10&-&-&-&-&-\\
7&-&9&13&-&-&-&-&-
\end{array}
\tag{5.1}
\]
This table is a finite polynomial certificate, not the output of an unrecorded decision procedure.  To check it, put
\[
 A_r=\Norm(\tau_r)-4.
\]
When $A_r\le0$, the separable quadratic identity
\[
 \begin{aligned}
 \Delta(x,y)
 &= (1-x)(1-y)\Delta(0,0)+x(1-y)\Delta(1,0)\\
 &\quad +(1-x)y\Delta(0,1)+xy\Delta(1,1)\\
 &\quad -A_r\bigl(x(1-x)+y(1-y)\bigr)
 \end{aligned}
\]
shows that the smallest of the four displayed corner values is a lower bound.  When $A_r>0$, only rows $r=0,4,5,7$ occur.  The remaining monotonicity checks are recorded explicitly below; each derivative bound holds on $0\le x,y\le1$.
\[
\begin{array}{cc|cc|c|c}
r&c&\partial_x\Delta_{r,c}&\partial_y\Delta_{r,c}
   &\text{minimizing corner}&\Delta_{r,c}\text{ there}\\ \toprule
0&0&\ge 8&\le-12&(0,1)&0\\
0&1&\ge 8&\le-2 &(0,1)&1\\
0&2&\ge 8&\ge 6 &(0,0)&5\\
0&3&\ge 8&\ge16 &(0,0)&16\\
4&1&\ge16&\le 0 &(0,1)&6\\
4&2&\ge16&\ge 8 &(0,0)&10\\
5&1&\ge12&\le-2 &(0,1)&5\\
5&2&\ge12&\ge 6 &(0,0)&9\\
7&1&\ge16&\le-6 &(0,1)&9\\
7&2&\ge16&\ge 2 &(0,0)&13
\end{array}
\]
Thus every numerical entry in (5.1) follows from an integer corner evaluation and one of the two displayed finite arguments.

The pair marked $\times$ is $(t,t')=(1-i,1-2i)$.  To exclude it, write
\[
 t'+w=A-iB,
 \qquad A=1+x,
 \qquad B=2-y,
\]
so $B>1$.  Since
\[
 v=\frac\gamma{A-iB}=p+iq,
\]
we have, with $R=A^2+B^2$,
\[
 p=\frac{A-B}{R},\qquad q=\frac{A+B}{R},\qquad p^2+q^2=\frac2R.
\]
If $t=1-i$, then
\[
 \Imm\!\left(\frac\gamma{1-i+v}\right)-1
 =\frac{q(1-q)-p(1+p)}{(1+p)^2+(q-1)^2}
 =\frac{2(B-1)}{R\bigl((1+p)^2+(q-1)^2\bigr)}>0.
\]
This contradicts $\operatorname{Sq}(u)$.

The three zero entries in (5.1) can attain equality only at $x=0$, $y=1$.  The actual predicate $\operatorname{Sq}(w)$ has $y<1$, so all admissible cases give $\Delta_{r,c}(x,y)>0$.  Hence $\Norm(uvw)<1$.
\end{proof}

\leanref{CCL.CCFContraction.three_step_contraction}

\begin{theorem}[finite Gaussian-rational CCFs]\label{thm:ccf-finite}
Algorithm~\ref{alg:ccf} terminates on every input $z_0\in\Qi$.  If it emits $t_0,\ldots,t_n$ and $M=A(t_0,\gamma)\cdots A(t_n,\gamma)$, then $M\omega=\Proj{z_0}{1}$.
\end{theorem}

\begin{proof}
The reconstruction statement follows by finite composition of the inverse-branch identities, with the terminal digit satisfying $A(t_n,\gamma)\omega=\Proj{t_n}{1}$.  For termination, carry an uncancelled representation $z_n=P_n/Q_n$ with $P_n,Q_n\in\Zi$ and $Q_n\ne0$.  Then
\[
 P_{n+1}=\gamma Q_n,
 \qquad
 Q_{n+1}=P_n-t_nQ_n,
 \qquad
 \frac{Q_{n+1}}{Q_n}=z_n-t_n.
\]
Execute at most three CCF clauses.  If one is terminal, stop.  If all three are non-terminal, \autoref{lem:ccf-contraction} gives
\[
 \Norm(Q_{n+3})<\Norm(Q_n).
\]
Thus the encoder is defined by well-founded recursion on the positive natural number $\Norm(Q_n)$, with each recursive call producing at most three digits before passing to a strictly smaller measure.
\end{proof}

\leanref{CCL.CCFEncoder.encode_reconstruct}

\section{Complex continued logarithms}

\subsection{The finite sector table}

The ordered finite representative list is
\[
 \mathcal R_8=(1,1+i,i,-1+i,-1,-1-i,-i,1-i).
\]
The three-bit word $uvw$ selects $s=(-1)^u i^v\gamma^w$.  Equivalently, the following finite table defines both the sector predicate and the representative.
\[
\begin{array}{ccl}
uvw&s&\text{rational sector predicate for }z=x+iy\\ \toprule
000&1&x\ge\tfrac12,\quad -\tfrac12<y<\tfrac12\\
001&1+i&x\ge\tfrac12,\quad y\ge\tfrac12\\
010&i&-\tfrac12<x<\tfrac12,\quad y\ge\tfrac12\\
011&-1+i&x\le-\tfrac12,\quad y\ge\tfrac12\\
100&-1&x\le-\tfrac12,\quad -\tfrac12<y<\tfrac12\\
101&-1-i&x\le-\tfrac12,\quad y\le-\tfrac12\\
110&-i&-\tfrac12<x<\tfrac12,\quad y\le-\tfrac12\\
111&1-i&x\ge\tfrac12,\quad y\le-\tfrac12
\end{array}
\tag{6.1}
\]
The boundary convention makes sector selection a deterministic nest of rational comparisons.  For $s$ in the displayed eight-entry list, write $\operatorname{Sect}_s(z)$ for the corresponding row predicate.  This is a finite family of decidable predicates.

Define three norm predicates:
\[
 \begin{aligned}
 \operatorname{Band}(z)&:\Longleftrightarrow 1\le\Norm(z)\le2,\\
 \operatorname{Entry}(z)&:\Longleftrightarrow 1\le\Norm(z)<2,\\
 \operatorname{Loop}(z)&:\Longleftrightarrow 1<\Norm(z)\le2.
 \end{aligned}
\]
These are predicates, not names for collections of points.

\begin{lemma}[sector partition]\label{lem:sector-partition}
If $\Norm(z)\ge1$, exactly one row of (6.1) holds.
\end{lemma}

\begin{proof}
The inequalities $|x|<1/2$ and $|y|<1/2$ would imply $\Norm(z)<1/2$.  Hence at least one coordinate crosses a displayed threshold.  If both do, their signs choose a diagonal row; otherwise the sign of the unique large coordinate chooses an axial row.  The strict transverse inequalities in the axial rows assign every threshold line to exactly one diagonal row.
\end{proof}

\leanref{CCL.Sectors.Sector.select_spec; CCL.Sectors.Sector.holds_unique}

\begin{lemma}[sector contraction]\label{lem:sector-contraction}
Assume $\operatorname{Band}(z)$ and let $s$ be selected by (6.1).  Then
\[
 \Norm(z-s)<\frac12.
\]
\end{lemma}

\begin{proof}
Sign changes and coordinate exchange reduce the proof to $s=1$ and $s=1+i$.

For $s=1$, write $z=x+iy$ with $x\ge1/2$ and $|y|<1/2$.  From $x^2+y^2\le2$ we get $x<3/2$, because $x\ge3/2$ would imply $x^2>2$.  Hence $(x-1)^2\le1/4$, while $y^2<1/4$.

For $s=1+i$, put
\[
 x=\frac12+a,\qquad y=\frac12+b,
 \qquad a,b\ge0,
\]
and write $p=a+b$, $q=a^2+b^2$.  The band inequalities give
\[
 \frac12\le p+q\le\frac32.
\]
We claim $q<p$.  Otherwise $p\le q\le p^2$.  The lower band inequality excludes $p=0$, so $p\ge1$, whence $p+q\ge2$, a contradiction.  Therefore
\[
 \Norm(z-(1+i))
 =\left(a-\frac12\right)^2+\left(b-\frac12\right)^2
 =q-p+\frac12<\frac12.
\]
\end{proof}

\leanref{CCL.Sectors.Sector.select_contraction}

\subsection{Raw bitstream algorithm}

\begin{algorithm}[CCL expansion]\label{alg:ccl}
Let $z_0\in\Qi$ be non-zero.

\emph{Initial normalization.}  Write $\Norm(z_0)=A/B$ with $A,B>0$.
\begin{itemize}
\item If $2B\le A$, use \autoref{lem:finite-normalization}(a) to obtain $k\ge1$ with
\[
 \operatorname{Entry}(z_0/\gamma^k),
\]
put $y_0=z_0/\gamma^k$, and emit $1^k0$.
\item If $A<2B$, use \autoref{lem:finite-normalization}(b) to obtain $m\ge0$ with
\[
 \operatorname{Entry}(\gamma^mz_0),
\]
put $y_0=\gamma^mz_0$, and emit $0\,1^m0$.
\end{itemize}
Select $s_0$ by (6.1), emit its three-bit code, and put $r_0=y_0-s_0$.  If $r_0=0$, stop.  Otherwise put
\[
 z_1=\frac1{\gamma r_0}.
\]

\emph{Subsequent loop.}  By \autoref{lem:sector-contraction}, $\Norm(z_j)>1$.  Use \autoref{lem:finite-normalization}(c) to obtain the unique $k\ge0$ such that
\[
 \operatorname{Loop}(z_j/\gamma^k).
\]
Put $y_j=z_j/\gamma^k$, emit $1^k0$, select $s_j$ by (6.1), emit its three-bit code, and put $r_j=y_j-s_j$.  Stop if $r_j=0$; otherwise put
\[
 z_{j+1}=\frac1{\gamma r_j}
\]
and repeat.
\end{algorithm}

The initial stage uses $[1,2)$ because its leading bit distinguishes division from multiplication.  Every later stage uses $(1,2]$.  This single boundary change removes the endpoint-preservation case from the denominator argument and makes every non-terminal loop step strictly descending.

The zero input may be represented by a separate finite zero tag.  It is not forced into the annular code.

\subsection{Complete-block notation}

The raw bitstream definition comes first because it is the representation.  We now introduce notation that abbreviates only \emph{completed} blocks.

Let $\code(s)\in\bits^3$ be the three-bit code in (6.1).  For $k\in\N$ and one of the eight displayed representatives, define
\[
 [k,s]\;:=\;1^k0\,\code(s).
\]
For the low initial branch, write
\[
 0[m,s]\;:=\;0\,[m,s]\;=\;0\,1^m0\,\code(s).
\]
Thus $0[m,s]$ means one extra initial zero followed by the completed block $[m,s]$.

A complete non-zero CCL word therefore has one of the forms
\[
 [k_0,s_0][k_1,s_1]\cdots[k_r,s_r],\qquad k_0\ge1,\quad k_1,\ldots,k_r\ge0,
\]
or
\[
 0[m,s_0][k_1,s_1]\cdots[k_r,s_r],\qquad m\ge0,\quad k_1,\ldots,k_r\ge0.
\]
The notation is not a replacement for raw finite words.  A prefix can end after only some of the unary ones, immediately after the delimiter zero, or after one or two sector bits.  Such a prefix has no lossless representation as a list of pairs $[k,s]$.  Its unresolved bits nevertheless affect its cylinder predicate and its bit length.  In particular,
\[
 |[k,s]|=k+4,\qquad |0[m,s]|=m+5.
\]

\subsection{Decoder matrices}

For a completed regular block $[k,s]$, the encoder equations are
\[
 y=\frac z{\gamma^k},\qquad
 r=y-s,
 \qquad
 z^+=\frac1{\gamma r}.
\]
Thus
\[
 z=\gamma^k\left(s+\frac1{\gamma z^+}\right),
\]
and a Gaussian-integral decoder matrix is
\[
 B(k,s)=S_{\gamma^k}T_sI_\gamma
 =\mat{\gamma^{k+1}s}{\gamma^k}{\gamma}{0}.
\tag{6.2}
\]
A completed block updates a prefix matrix by $M\mapsto MB(k,s)$.

For the low initial block $0[m,s]$, put
\[
 E^-_m=\mat{1}{0}{0}{\gamma^m}.
\]
Its matrix is
\[
 L(m,s)=E^-_mT_sI_\gamma
 =\mat{\gamma s}{1}{\gamma^{m+1}}{0}.
\tag{6.3}
\]
The high initial branch uses $B(k,s)$ exactly as in (6.2).

\subsection{Incremental angular parsing}\label{subsec:angular-folding}

The existing word $uvw$ already has a useful algebraic structure.  Put
\[
 \epsilon=(-1)^u i^v=i^{2u+v}\in(1,i,-1,-i).
\]
Thus $uv$ is the ordinary two-bit index of a quarter-turn and the last bit $w$ chooses the axial representative $\epsilon$ or the following diagonal representative $\epsilon\gamma$.  Equivalently, the integer with binary digits $uvw$ is the cyclic octant index.  No alternative angular code is needed.

\begin{lemma}[unit covariance of the sector table]\label{lem:sector-unit-covariance}
For $\epsilon\in(1,i,-1,-i)$, $h\in\bits$, and $y\in\Qi$,
\[
 \operatorname{Sect}_{\epsilon\gamma^h}(y)
 \quad\Longleftrightarrow\quad
 \operatorname{Sect}_{\gamma^h}(\epsilon^{-1}y).
\tag{6.4}
\]
\end{lemma}

\begin{proof}
Multiplication by a Gaussian unit rotates the coordinate inequalities by a multiple of ninety degrees.  Direct substitution in the eight rows of (6.1), including their strict and non-strict boundary signs, gives (6.4).  This is a finite check of four units and two canonical rows.
\end{proof}

\leanref{CCL.SectorCovariance.holds_covariant}

Suppose the current decoder relation is represented by
\[
 X=M\Proj{y}{1},\qquad M=\mat{a}{b}{c}{d}.
\]
After reading $u$, put $y_1=(-1)^uy$ and right-multiply by $S_{(-1)^u}$.  In particular,
\[
 u=1:
 \qquad
 MS_{-1}=\mat{-a}{b}{-c}{d}.
\tag{6.5}
\]
After reading $v$, put $y_2=i^{-v}y_1$ and right-multiply by $S_{i^v}$.  The accumulated frame is $\epsilon$, so the temporary matrix is $MS_\epsilon$ and the remaining branch is one of the two canonical choices $\gamma^0$ or $\gamma^1$.

Let
\[
 C_h=T_{\gamma^h}I_\gamma\qquad(h\in\bits).
\]
The full branch satisfies the projective matrix identity
\[
 T_{\epsilon\gamma^h}I_\gamma
 \ \sim\ 
 S_\epsilon C_hS_\epsilon,
\tag{6.6}
\]
where $\sim$ means equality up to multiplication by a non-zero scalar.  Indeed, the right-hand matrix is exactly $\epsilon$ times the left-hand matrix.  Hence a finite parser may process the angular code as two unit folds, one binary canonical branch, and one outgoing unit fold.  It needs only a four-valued orientation frame and a one-bit axial--diagonal choice, rather than an unrelated eight-way matrix table.

The unit factors in (6.5)--(6.6) are exact changes of tail coordinate.  Their accompanying bit predicates remove incompatible sectors, but the matrix fold itself does not yet assert numerical contraction.  Section~\ref{subsec:projective-quadrance-contraction} makes this distinction precise.

\subsection{Finite termination}

\begin{lemma}[inversion-stage descent]\label{lem:ccl-inversion-descent}
Let $y=P/Q$ with $P,Q\in\Zi$, $Q\ne0$, and $\operatorname{Band}(y)$.  Let $s$ be the selected sector representative.  If $P-sQ\ne0$, then
\[
 z^+=\frac{Q}{\gamma(P-sQ)}
\]
is an uncancelled representation of the next post-inversion state and
\[
 \Norm\bigl(\gamma(P-sQ)\bigr)<\Norm(Q).
\]
\end{lemma}

\begin{proof}
By \autoref{lem:sector-contraction},
\[
 \Norm(P-sQ)<\frac12\Norm(Q).
\]
Multiplication by $\gamma$ doubles the norm.
\end{proof}

\leanref{CCL.Decoder.inversion_stage_descent}

\begin{theorem}[finite Gaussian-rational CCLs]\label{thm:ccl-finite}
Algorithm~\ref{alg:ccl} terminates on every non-zero input $z_0\in\Qi$.  If $M$ is the product of (6.2)--(6.3) for the complete blocks it emits, in emission order, then $M\omega=\Proj{z_0}{1}$.
\end{theorem}

\begin{proof}
Finite matrix composition proves the reconstruction statement: at the final exact subtraction, the unresolved post-inversion coordinate is $\omega$, and the last block sends it to the terminal annulus value.  For termination, write each annulus-stage value without cancellation as
\[
 y_j=\frac{P_j}{Q_j},\qquad P_j,Q_j\in\Zi,
 \qquad Q_j\ne0.
\]
The initial value satisfies $\operatorname{Entry}(y_0)$, and every later one satisfies $\operatorname{Loop}(y_j)$.

Assume the $j$th sector step is non-terminal.  Before radial normalization, \autoref{lem:ccl-inversion-descent} gives
\[
 P_j^+=Q_j,
 \qquad
 Q_j^+=\gamma(P_j-s_jQ_j),
 \qquad
 \Norm(Q_j^+)<\Norm(Q_j).
\]
The next loop normalization chooses $k\ge0$ with
\[
 1<\frac{\Norm(P_j^+)}{2^k\Norm(Q_j^+)}\le2.
\]
It then sets
\[
 P_{j+1}=P_j^+,
 \qquad
 Q_{j+1}=\gamma^kQ_j^+.
\]
Therefore
\[
 \Norm(Q_{j+1})<\Norm(P_j^+)=\Norm(Q_j).
\]
Thus one non-terminal sector step followed by the next finite normalization makes a recursive call with a strictly smaller positive natural denominator norm.  This is a well-founded definition of the encoder.
\end{proof}

\leanref{CCL.Encoder.encode_reconstruct}

\section{Finite-prefix cylinder semantics}

\subsection{Raw prefixes and parser state}

A bit word is a finite list over $\bits$.  For a non-zero rational input $z$, let $\operatorname{enc}(z)$ be the finite word produced by \autoref{alg:ccl}.  For a finite bit word $w$, define the cylinder predicate
\[
 \Cyl(w,z)\;:\Longleftrightarrow\; w\pref\operatorname{enc}(z),
\tag{7.1}
\]
where $\pref$ is the ordinary finite-list prefix relation.  This is a predicate, not the name of a set.  It is decidable by running the encoder only until it has emitted $|w|$ bits or has terminated.

At a radial parser phase, let $x$ denote the current rational state before normalization.  The information in an unfinished unary run is already a single rational inequality; no list of its unboundedly many possible final exponents is formed.  The exact equivalences are:
\[
\begin{array}{ccl}
\text{phase}&\text{raw fragment just read}&\text{predicate on }x\\ \toprule
\text{initial high}&1^r,\ r\ge1&2^r\le\Norm(x)\\
\text{initial high}&1^r0,\ r\ge1&2^r\le\Norm(x)<2^{r+1}\\
\text{initial low}&0&\Norm(x)<2\\
\text{initial low}&0\,1^r,\ r\ge1&2^{r-1}\Norm(x)<1\\
\text{initial low}&0\,1^m0,\ m\ge0&1\le2^m\Norm(x)<2\\
\text{subsequent loop}&1^r,\ r\ge1&2^r<\Norm(x)\\
\text{subsequent loop}&1^k0,\ k\ge0&2^k<\Norm(x)\le2^{k+1}
\end{array}
\tag{7.2}
\]
These statements are immediate from the maximal or minimal bounded searches in \autoref{lem:finite-normalization}.  After a delimiter, the corresponding normalized value is $x/\gamma^r$, $\gamma^m x$, or $x/\gamma^k$, and any already-read sector bits further refine its predicate by a finite disjunction over the eight table rows.

A deterministic parser scans the finite word $w$ from left to right.  Its phase records only:
\begin{enumerate}[label=\textup{(\roman*)}]
\item whether the initial high or low header is being read;
\item the finite unary count accumulated in the current scale run;
\item zero, one, or two already-read sector bits; or
\item that a completed block has just updated the matrix.
\end{enumerate}
For a given finite prefix the counter is bounded by $|w|$; no parser state stores future exponents or future continuations.

If the current sector-bit fragment is $q$ of length at most two, obtain the finite candidate list $\Sigma_q$ by filtering the fixed eight-entry list $\mathcal R_8$ for representatives whose code begins with $q$, without changing the order.  This list has respectively eight, four, or two entries.  The cylinder predicate at that point is the finite disjunction of the corresponding sector predicates.  Once the third bit arrives, $\Sigma_q$ has one entry and the block matrix can be applied.

This distinction is essential.  The block notation $[k,s]$ summarizes completed syntax, while (7.1) uses every raw bit.  For example, the ten-bit prefix
\[
 1110101110
\]
of the code in \autoref{ex:ccl-355} below contains the complete block $[3,-1-i]$ followed by the raw fragment $110$.  That fragment says that the next scale exponent is exactly two but contains no sector bit.  It changes the cylinder depth from seven bits to ten even though no second pair $[k,s]$ is complete.

\subsection{Canonical projective truncation}

The parser multiplies (6.2) or (6.3) whenever a block becomes complete.  Let $M_w^\circ$ be the resulting product after the finite scan of $w$.  An incomplete unary scale factor may also be accumulated, but it fixes $\omega$, so it does not change the following coordinate.  An incomplete sector code has no unique translation matrix and is therefore not applied.

Define the canonical projective record
\[
 \Pi(w)=\bigl(w,\,M_w^\circ\omega\bigr).
\tag{7.3}
\]
The first component is a finite tag and is part of the record.  Distinct prefixes may have the same projective coordinate while remaining distinct records.  The coordinate is not asserted to be an affine rational whose encoding begins with $w$; the tag places it at the corresponding formal projective boundary.  Thus (7.3) is the finite tagged replacement for adjoining a projective-closure endpoint, without naming a closure collection.  At a completed block boundary it is the usual matrix truncation.  Before the first completed block, its coordinate is the formal point $\omega$.

\subsection{A finite bit-prefix minimax theorem}

Append an external terminal marker $\term$ to a completed rational code.  This marker is not a CCL bit and is not charged to the bit budget; it merely distinguishes a terminated code from a strict binary extension.

A rational cylindrical gauge is a function $g$ assigning a positive rational $g(w)$ to every finite binary word and satisfying
\[
 g(wb)<g(w)\qquad(b\in\bits).
\]
The bit gauge is
\[
 g_2(w)=\frac1{2^{|w|}}.
\]
For two finite well-formed tags $a,b$ over bits with an optional final terminal marker, define
\[
 d_g(a,b)=
 \begin{cases}
 0,&a=b,\\
 g(\operatorname{lcp}(a,b)),&a\ne b,
 \end{cases}
\tag{7.4}
\]
where $\operatorname{lcp}$ is their longest common binary prefix.  The usual longest-common-prefix argument gives the strong triangle inequality, so this is an ultrametric on finite tags, not on numerical points.

\begin{theorem}[finite two-witness minimax]\label{thm:finite-minimax}
Let $w$ be a finite non-terminal prefix.  Suppose there are rational witnesses $z_0,z_1$ such that
\[
 \Cyl(w0,z_0),\qquad \Cyl(w1,z_1).
\]
Put $x_b=\operatorname{enc}(z_b)\term$ and, for a candidate finite tag $a$, define the finite two-witness radius
\[
 R_w(a;z_0,z_1)=\max\bigl(d_g(a,x_0),d_g(a,x_1)\bigr).
\]
Then:
\begin{enumerate}[label=\textup{(\alph*)}]
\item If $w\pref a$, then $R_w(a;z_0,z_1)=g(w)$.  Moreover, for every rational $z$ satisfying $\Cyl(w,z)$,
\[
 d_g\bigl(a,\operatorname{enc}(z)\term\bigr)\le g(w).
\]
\item If $w\not\pref a$, then $R_w(a;z_0,z_1)>g(w)$.
\end{enumerate}
Consequently the canonical record $\Pi(w)$, whose tag is $w$, is minimax for this finite-prefix gauge.  The optimum is not unique: every candidate tag extending $w$ has the same radius.
\end{theorem}

\begin{proof}
Assume first that $w\pref a$.  Every compatible terminal code shares $w$ with $a$, so repeated strict child monotonicity gives the displayed upper bound.  If $a=w$, both witness distances equal $g(w)$.  If $a$ has a next symbol after $w$, at least one of $x_0,x_1$ has a different next symbol; this also covers the case in which $a$ ends there with the terminal marker.  Its longest common prefix with $a$ is therefore exactly $w$.  Hence the finite maximum is $g(w)$.

If $w\not\pref a$, put $v=\operatorname{lcp}(a,w)$.  Then $v$ is a proper prefix of $w$, and both witnesses have longest common prefix $v$ with $a$.  Therefore their distances equal $g(v)>g(w)$.
\end{proof}

\leanref{CCL.PrefixMinimax.finite_twoWitness_minimax}

\begin{corollary}[pointwise bit budget]\label{cor:bit-budget}
Under the hypotheses of \autoref{thm:finite-minimax}, a decoder that has read exactly the word $w$ cannot have two-witness radius below $1/2^{|w|}$ in the bit gauge.  The canonical record $\Pi(w)$ attains exactly $1/2^{|w|}$.
\end{corollary}

\begin{proof}
Apply Theorem~\ref{thm:finite-minimax} with the bit gauge
$g_2(w)=2^{-|w|}$.
\end{proof}

This theorem deliberately concerns prefix information.  It does not assert
optimality of the numerical coordinate in (7.3), either in Euclidean
distance or in Hermitian projective quadrance.

\subsection{Hermitian projective quadrance and completed-block contraction}\label{subsec:projective-quadrance-contraction}

For non-zero projective coordinates define their \emph{Hermitian projective
quadrance}
\[
 \chi^2\bigl(\Proj{x_0}{x_1},\Proj{y_0}{y_1}\bigr)
 =\frac{\Norm(x_0y_1-y_0x_1)}
 {(\Norm(x_0)+\Norm(x_1))(\Norm(y_0)+\Norm(y_1))}.
\tag{7.5}
\]
This is the Gaussian-norm extension of projective quadrance
\cite{wildberger-projective}; conjugation makes the underlying coordinate
form Hermitian.  Writing $\overline{z}$ for Gaussian conjugation, the
following rational Lagrange identity shows that $0\le\chi^2\le1$:
\[
 \begin{aligned}
 &(\Norm(x_0)+\Norm(x_1))(\Norm(y_0)+\Norm(y_1))
   -\Norm(x_0y_1-y_0x_1)\\
 &\hspace{4em}=
 \Norm\bigl(x_0\overline{y_0}+x_1\overline{y_1}\bigr)\ge0.
 \end{aligned}
\tag{7.6}
\]
The expression is unchanged when either homogeneous coordinate pair is multiplied by a non-zero scalar.

For a matrix $M=\mat{a}{b}{c}{d}$ and an affine tail $t$, put
\[
 D_M(t)=\Norm(at+b)+\Norm(ct+d).
\]

\begin{proposition}[one-tail truncation identity]\label{prop:projective-quadrance-truncation}
Let
\[
 p=M\omega=\Proj{a}{c},\qquad
 z=M\Proj{t}{1}=\Proj{at+b}{ct+d},
\]
where $\det M\ne0$.  Then
\[
 \chi^2(z,p)
 =\frac{\Norm(\det M)}
 {(\Norm(a)+\Norm(c))D_M(t)}.
\tag{7.7}
\]
\end{proposition}

\begin{proof}
The projective numerator is
\[
 \Norm\bigl((at+b)c-a(ct+d)\bigr)=\Norm(bc-ad).
\]
The denominator is the one displayed.
\end{proof}

\leanref{CCL.Chordal.Mat2.chordalSq_truncation}

\begin{proposition}[two-tail homographic identity]\label{prop:two-tail-projective-quadrance}
Let $\det M\ne0$ and let $t,u\in\Qi$.  Then
\[
 \chi^2\left(M\Proj{t}{1},M\Proj{u}{1}\right)
 =\frac{\Norm(\det M)\Norm(t-u)}{D_M(t)D_M(u)}.
\tag{7.8}
\]
Consequently
\[
 \frac{\chi^2\left(M\Proj{t}{1},M\Proj{u}{1}\right)}
 {\chi^2\left(\Proj{t}{1},\Proj{u}{1}\right)}
 =\frac{\Norm(\det M)(1+\Norm(t))(1+\Norm(u))}
 {D_M(t)D_M(u)}
\tag{7.9}
\]
when $t\ne u$.
\end{proposition}

\begin{proof}
The determinant of the two output columns is
\[
 (at+b)(cu+d)-(au+b)(ct+d)=(ad-bc)(t-u).
\]
Taking Gaussian norms gives (7.8), and division by the affine formula for $\chi^2(t,u)$ gives (7.9).
\end{proof}

\leanref{CCL.Chordal.Mat2.chordalSq_act_affine}

\begin{lemma}[scale and unit folds preserve the truncation expression]\label{lem:scale-fold-invariance}
Let $\lambda\in\Qi^\times$, let $t=\lambda t'$, and put $M'=MS_\lambda$.  Then
\[
 M'\Proj{t'}{1}=M\Proj{t}{1},
 \qquad
 M'\omega=M\omega,
\]
and
\[
 \begin{aligned}
 \Norm(\det M')&=\Norm(\lambda)\Norm(\det M),\\
 \Norm(a\lambda)+\Norm(c\lambda)
   &=\Norm(\lambda)(\Norm(a)+\Norm(c)),\\
 D_{M'}(t')&=D_M(t).
 \end{aligned}
\tag{7.10}
\]
Thus the value of (7.7) is exactly unchanged.  If also $u=\lambda u'$, then
\[
 \chi^2\left(M'\Proj{t'}{1},M'\Proj{u'}{1}\right)
 =\chi^2\left(M\Proj{t}{1},M\Proj{u}{1}\right),
\]
so the two-tail expression is unchanged as well.
\end{lemma}

\begin{proof}
All statements follow by multiplying $M$ on the right by $S_\lambda$.  The first-column factor in the denominator of (7.7) acquires the same norm multiplier as the determinant, while $D_M$ and the transported point remain unchanged.  For two tails, $\Norm(t'-u')=\Norm(t-u)/\Norm(\lambda)$ cancels the determinant multiplier in (7.8).
\end{proof}

\leanref{CCL.ScaleFold.chordalSq_twoTail}

A unary $1$ is the case $\lambda=\gamma$: the first decoder column and the determinant norm both double.  The bit $u=1$ in the angular code is the unit case $\lambda=-1$, giving exactly (6.5); the $v$ bit similarly uses $\lambda=i^v$.  These bits strengthen the rational tail predicate, but coefficient growth alone cannot prove contraction because it is cancelled in (7.7).

For $\lambda\in\Qi^\times$ and $s\in\mathcal R_8$, define the completed-block map
\[
 F_{\lambda,s}(t)=\lambda\left(s+\frac1{\gamma t}\right)
\]
by the projective matrix
\[
 A_{\lambda,s}
 =\mat{\lambda\gamma s}{\lambda}{\gamma}{0}.
\tag{7.11}
\]
For $\lambda=\gamma^k$ this is $B(k,s)$ from (6.2); for $\lambda=\gamma^{-m}$, multiplication by the non-zero scalar $\gamma^m$ gives $L(m,s)$ from (6.3).  It sends $\omega$ to $\Proj{\lambda s}{1}$.  Put
\[
 L=\Norm(\lambda)>0,
 \qquad
 q(L)=\frac{2L}{(1+L)^2}.
\tag{7.12}
\]

\begin{theorem}[completed-block contraction of Hermitian projective quadrance]\label{thm:block-projective-quadrance-contraction}
For $j\in\{1,2\}$, let $X_j$ be either $\omega$ or an affine coordinate $\Proj{t_j}{1}$ with $t_j\in\Qi^\times$.  For every affine $X_j$, put
\[
 y_j=s+\frac1{\gamma t_j}
\]
and assume $\operatorname{Band}(y_j)$ and $\operatorname{Sect}_s(y_j)$.  Then
\[
 \chi^2\bigl(A_{\lambda,s}X_1,A_{\lambda,s}X_2\bigr)
 \le q(L)\chi^2(X_1,X_2)
 \le\frac12\chi^2(X_1,X_2).
\tag{7.13}
\]
If $X_1\ne X_2$, the first inequality is strict whenever at least one tail is affine.
\end{theorem}

\begin{proof}
For an affine compatible tail, \autoref{lem:sector-contraction} gives
\[
 \Norm(y_j-s)<\frac12,
 \qquad
 \Norm(t_j)=\frac1{2\Norm(y_j-s)}>1.
\tag{7.14}
\]
Moreover
\[
 \Norm(\det A_{\lambda,s})=2L
\]
and, because $\lambda\gamma t_jy_j=\lambda\gamma s t_j+\lambda$,
\[
 D_{A_{\lambda,s}}(t_j)
 =2\Norm(t_j)\bigl(1+L\Norm(y_j)\bigr).
\tag{7.15}
\]
If both tails are affine and distinct, (7.9) becomes
\[
 \frac{\chi^2(F_{\lambda,s}(t_1),F_{\lambda,s}(t_2))}
 {\chi^2(t_1,t_2)}
 =\frac L2
 \prod_{j=1}^2
 \frac{1+1/\Norm(t_j)}{1+L\Norm(y_j)}.
\]
By (7.14), each numerator factor is strictly below $2$; by the band predicate, each denominator factor is at least $1+L$.  Hence the ratio is strictly below $q(L)$.  Equal affine tails give zero on both sides.

If one tail is $\omega$, symmetry reduces to $X_2=\omega$.  Since
\[
 A_{\lambda,s}\omega=\Proj{\lambda s}{1},
\]
direct use of the affine formula gives
\[
 \frac{\chi^2(F_{\lambda,s}(t_1),F_{\lambda,s}(\omega))}
 {\chi^2(t_1,\omega)}
 =\frac L2\,
 \frac{1+1/\Norm(t_1)}
 {(1+L\Norm(y_1))(1+L\Norm(s))}
 <\frac{L}{(1+L)^2}<q(L),
\]
because $\Norm(s)\ge1$.  If both tails are $\omega$, both sides vanish.  Finally,
\[
 (1+L)^2-4L=(1-L)^2\ge0
\]
gives $q(L)\le1/2$.
\end{proof}

\leanref{CCL.BlockContraction.Tail.contraction}

A regular or high initial block with unary exponent $r$ has $L=2^r$.  A low initial block $0[r,s]$ has $L=2^{-r}$.  In both cases the factor is
\[
 q_r=\frac{2^{r+1}}{(1+2^r)^2},
 \qquad
 q_0=\frac12,
 \qquad
 q_r<\frac12\quad(r>0).
\tag{7.16}
\]
Thus long unary runs do improve the completed block's contraction, but only after the sector branch has been completed.

For a finite word $w$ and a rational $Q\ge0$, define the finite universal
Hermitian-projective-quadrance bound predicate, abbreviated by
\[
 \operatorname{HPQBound}(w,Q)
 :\Longleftrightarrow
 \forall z,z'\in\Qi^\times,
 \quad
 \Cyl(w,z)\land\Cyl(w,z')
 \Longrightarrow
 \chi^2\bigl(\Proj{z}{1},\Proj{z'}{1}\bigr)\le Q.
\tag{7.17}
\]
This is a pairwise certificate, not a supremum over a named cylinder.

\begin{corollary}[shrinking completed-block prefixes]\label{cor:completed-block-shrink}
Let $w$ be the raw concatenation of $n$ well-formed completed CCL blocks from the initial state, with unary exponents $r_0,\ldots,r_{n-1}$.  Then
\[
 \operatorname{HPQBound}\left(
 w,\prod_{j<n}q_{r_j}
 \right),
 \qquad
 \prod_{j<n}q_{r_j}\le2^{-n}.
\tag{7.18}
\]
\end{corollary}

\begin{proof}
Take two rational inputs satisfying the same prefix predicate.  Finite decoder substitution expresses each as the same outer block applied to a residual projective tail.  A residual tail is affine when its rational code continues and is $\omega$ when it terminates at that block.  At every shared block the corresponding two intermediate tails satisfy the hypotheses of \autoref{thm:block-projective-quadrance-contraction}.  Iterating the theorem through the finite block list gives the product in (7.18).  The unresolved final tails have Hermitian projective quadrance at most $1$ by (7.6), and every factor is at most $1/2$.
\end{proof}

\leanref{CCL.Cylinder.completedPrefix_chordBound}

Every raw extension preserves an already proved bound: if $w\pref v$, then (7.1) gives
\[
 \operatorname{HPQBound}(w,Q)
 \Longrightarrow
 \operatorname{HPQBound}(v,Q).
\tag{7.19}
\]
Strict improvement at every raw bit is false.  For example, the initial prefix $1$ means $2\le\Norm(z)$, while $10$ means $2\le\Norm(z)<4$.  If $R=\Norm(z)$ and $S=\Norm(z')$ are at least $2$, then
\[
 \Norm(z-z')\le2(R+S)
\]
from $\Norm(z-z')+\Norm(z+z')=2(R+S)$, and
\[
 \frac{2(R+S)}{(1+R)(1+S)}\le\frac89.
\]
The last inequality follows from
\[
 4RS-5R-5S+4
 =(4S-5)(R-2)+3(S-2)\ge0.
\]
Hence both $1$ and $10$ admit the universal bound $8/9$.  The rational witnesses $\gamma$ and $-\gamma$ satisfy both prefixes and have Hermitian projective quadrance $8/9$, so the delimiter bit produces no strict reduction at all.  Partial bits refine predicates; completed blocks supply the uniform geometric contraction.

\begin{remark}[why the expression is squared]
For affine $z,w\in\Qi$,
\[
 \chi^2\bigl(\Proj{z}{1},\Proj{w}{1}\bigr)
 =\frac{\Norm(z-w)}{(1+\Norm(z))(1+\Norm(w))},
\]
and
\[
 \chi^2\bigl(\Proj{z}{1},\omega\bigr)=\frac1{1+\Norm(z)}.
\]
We work with $\chi^2$ itself because every term is rational on
Gaussian-rational projective coordinates; taking a square root would
generally leave $\Q$.  Hermitian projective quadrance is therefore an exact
algebraic error expression.  Proposition~\ref{prop:projective-quadrance-truncation}
measures one canonical truncation, while
Proposition~\ref{prop:two-tail-projective-quadrance} and
Theorem~\ref{thm:block-projective-quadrance-contraction} control the pairwise
numerical width of a finite prefix.  These geometric statements complement,
but do not replace, prefix minimax optimality.
\end{remark}

\subsection{Finite approximation towers and progressive boxed names}

A finite rational code already contains approximations to its value.  The relevant substrings are its \emph{initial} substrings, namely its prefixes.  An internal substring has discarded the initial decoder state and is generally a name fragment for a different tail, not an approximation to the original input.

\begin{proposition}[prefix tower of a finite name]\label{prop:finite-prefix-tower}
Let $z\in\Qi$ be non-zero, let $v=\operatorname{enc}(z)$, put $m=|v|$, and for $0\le n\le m$ let $w_n$ be the first $n$ bits of $v$.  Then, for $0\le n<m$,
\[
 w_n\pref w_{n+1},
\]
while for every $0\le n\le m$,
\[
 \Cyl(w_n,z).
\]
If $x=v\term$ is the exact terminal tag, then for $n<m$,
\[
 d_g(w_n,x)=g(w_n).
\]
The completed decoder product satisfies
\[
 \Pi(v)=\bigl(v,\Proj{z}{1}\bigr).
\]
Whenever both binary children of $w_n$ admit rational witnesses, \autoref{thm:finite-minimax} says that $\Pi(w_n)$ is minimax among all records based only on that prefix.  If only one child is admissible, the next bit is forced rather than informative; one first extends through the finite run of forced bits and applies the theorem at the next genuine branch.
\end{proposition}

\begin{proof}
The prefix relations and cylinder predicates follow from the definition of $w_n$.  For $n<m$, the longest common binary prefix of $w_n$ and $v\term$ is exactly $w_n$, which gives the distance formula.  The final decoder identity is \autoref{thm:ccl-finite}.  The minimax assertion is \autoref{thm:finite-minimax}.
\end{proof}

\leanref{CCL.PrefixTower.distance_prefixAt}

Thus ``best approximation'' here has a precise information-theoretic meaning.  At every genuine branch, no method that has read only $w_n$ can improve its worst-case prefix radius.  Independently, whenever $w_n$ ends after $j$ completed blocks, \autoref{cor:completed-block-shrink} says that any two rational realizations of that prefix have Hermitian projective quadrance at most $2^{-j}$.  In particular, any rational witness carried with the prefix is a certified numerical approximant to the encoded rational.  The numerical coordinate of $\Pi(w_n)$ need not change at every raw bit and is not asserted to lie inside the affine cylinder; incomplete unary or sector fragments can still refine the tag and its enclosure before enough syntax exists to apply another block matrix.  At the final block of a rational name, the canonical projective coordinate is the rational itself.

We now add explicit Euclidean enclosures without introducing a completed complex point.  A \emph{rational box} is finite data
\[
 B=(x_-,x_+,y_-,y_+)\in\Q^4,
 \qquad x_-\le x_+,
 \qquad y_-\le y_+.
\]
For $z=x+iy\in\Qi$, write
\[
 \operatorname{InBox}(B,z)
 \;:\Longleftrightarrow\;
 x_-\le x\le x_+
 \ \text{and}\ 
 y_-\le y\le y_+.
\]
Its rational squared diagonal is
\[
 \Delta^2(B)=(x_+-x_-)^2+(y_+-y_-)^2.
\]
For boxes $C=(x'_-,x'_+,y'_-,y'_+)$ and $B$, write $C\sqsubseteq B$ for the four finite inequalities
\[
 x_-\le x'_-\le x'_+\le x_+,
 \qquad
 y_-\le y'_-\le y'_+\le y_+.
\]
This is the box-containment relation expressed entirely by rational comparisons.

\begin{definition}[boxed prefix certificate]
A boxed prefix certificate is finite data $(w,B,q)$, where $w$ is a finite bit word, $B$ is a rational box, and $q\in\Qi^\times$, together with proofs of
\[
 \Cyl(w,q)
\]
and
\[
 \forall z\in\Qi^\times,
 \qquad
 \Cyl(w,z)\Longrightarrow \operatorname{InBox}(B,z).
\tag{7.20}
\]
The witness $q$ proves finite satisfiability.  It is not declared to be the value named by the certificate.
\end{definition}

If $w\pref v$, then
\[
 \Cyl(v,z)\Longrightarrow\Cyl(w,z)
\tag{7.21}
\]
by transitivity of finite-list prefixes.  Hence extending a name prefix always refines its cylinder predicate, even though no cylinder is introduced as a named collection.  A certificate $(v,C,q')$ \emph{refines} $(w,B,q)$ when $w\pref v$ and $C\sqsubseteq B$.

A completed initial normalization header already gives a coarse rational box.  For a high header $1^k0$, every compatible input satisfies $\Norm(z)<2^{k+1}$ and is therefore enclosed by
\[
 [-2^{k+1},2^{k+1}]\mathbin{\times}[-2^{k+1},2^{k+1}].
\]
For a low header $0\,1^m0$, every compatible input satisfies $\Norm(z)<2$ and is enclosed by $[-2,2]\times[-2,2]$.  Every later prefix retains these bounds; additional bits permit sharper ones.

\begin{proposition}[completed blocks give explicit rational boxes]\label{prop:projective-quadrance-to-box}
Let $w$ be a completed-block prefix as in \autoref{cor:completed-block-shrink}, let
\[
 Q_w=\prod_{j<n}q_{r_j},
\]
and let $q\in\Qi^\times$ satisfy $\Cyl(w,q)$.  Suppose $K\in\Q_{>0}$ satisfies
\[
 \Cyl(w,z)\Longrightarrow\Norm(z)\le K
\tag{7.22}
\]
for every rational realization $z$.  Then
\[
 \Cyl(w,z)
 \Longrightarrow
 \Norm(z-q)\le(1+K)^2Q_w.
\tag{7.23}
\]
Given a positive rational $\varepsilon$, if
\[
 Q_w\le\frac{\varepsilon^2}{16(1+K)^2},
\tag{7.24}
\]
then the rational box
\[
 B_\varepsilon(q)=
 \left(\Ree q-\frac\varepsilon4,\Ree q+\frac\varepsilon4,
       \Imm q-\frac\varepsilon4,\Imm q+\frac\varepsilon4\right)
\tag{7.25}
\]
satisfies (7.20), and
\[
 \Delta^2(B_\varepsilon(q))=\frac{\varepsilon^2}{2}\le\varepsilon^2.
\]
\end{proposition}

\begin{proof}
For compatible $z$ and $q$, (7.18), the affine formula for $\chi^2$, and (7.22) give
\[
 \Norm(z-q)
 =\chi^2(z,q)(1+\Norm(z))(1+\Norm(q))
 \le(1+K)^2Q_w.
\]
Under (7.24), each coordinate difference has square at most $\varepsilon^2/16$, hence absolute value at most $\varepsilon/4$.  This proves membership in (7.25); its squared diagonal is a direct rational calculation.

Clearing the positive denominators in (7.24) gives the equivalent rational
test $(1+K)^2Q_w\le(\varepsilon/4)^2$.  The calculation is homogeneous in
the projective coordinates and therefore also covers an explicitly rescaled
integral decoder realization.
\end{proof}

\leanref{CCL.BoxTransport.completed_affine_epsilonBox_contains_of_paper_bound}

For a completed high header $1^k0$, one may take $K=2^{k+1}$; for a completed low header $0\,1^m0$, one may take $K=2$.  Thus a deterministic refinement procedure that, on input $n$, returns a longer satisfiable prefix containing at least $n$ completed blocks can be upgraded to a progressive boxed name.  Given $\varepsilon>0$, clear denominators and perform a bounded binary scan to choose $n$ with
\[
 2^{-n}\le\frac{\varepsilon^2}{16(1+K)^2}.
\]
Request that many blocks, use its rational witness as the centre in (7.25), and intersect with the previous certified box to retain nesting.  This proves, using only finite data, that arbitrarily many completed blocks produce arbitrarily small affine enclosures on every fixed bounded initial branch.

The next lemma remains useful for partial blocks, arithmetic transducers, and sharper anisotropic boxes.

\begin{lemma}[finite rational box transport]\label{lem:box-transport}
Let
\[
 M=\mat{a}{b}{c}{d},
 \qquad
 H_M(t)=\frac{at+b}{ct+d}.
\]
Suppose a finite parser state expresses every compatible affine input as $H_M(t)$ for rational tails satisfying a rational predicate $P(t)$.  Suppose finite evidence supplies a rational tail box $T$, a positive rational $\delta$, and the implications
\[
 P(t)\Longrightarrow\operatorname{InBox}(T,t),
 \qquad
 P(t)\Longrightarrow \delta\le\Norm(ct+d).
\]
Then finite rational interval arithmetic computes a rational box $B$ satisfying
\[
 P(t)\Longrightarrow
 \operatorname{InBox}\bigl(B,H_M(t)\bigr).
\]
If $P$ is the finite disjunction produced by an incomplete sector code, the same construction is performed for each entry of the finite candidate list and the coordinatewise rational hull is taken.
\end{lemma}

\begin{proof}
Write $t=x+iy$, $u=at+b$, and $v=ct+d$.  The real and imaginary parts of $u$ and $v$ are affine rational expressions in $x,y$, so ordinary interval addition and multiplication give rational bounds on them over $T$.  The affine value is
\[
 \frac{u}{v}=\frac{u\overline v}{\Norm(v)}.
\]
Interval arithmetic bounds the two coordinates of $u\overline v$.  It also gives a rational upper bound for $\Norm(v)$; replacing that upper bound by its maximum with $\delta$ gives a positive rational $D$ such that, on $P$,
\[
 0<\delta\le\Norm(v)\le D.
\]
Thus $1/\Norm(v)$ lies in the rational interval $[1/D,1/\delta]$, and one more finite interval multiplication bounds both coordinates of $u/v$.  A finite disjunction is handled by finitely many repetitions followed by rational minima and maxima.  No extremum over a named point collection is taken.  If no positive denominator-separation certificate is yet available, one retains the previous coarse box and reads more prefix bits rather than performing an uncertified division.

For completeness, an even simpler square-root-free enclosure is available.
Let $L_T$ be the endpoint-computed upper bound
for $\Norm(t)$ on $T$, and put
\[
 N(M,T)=2\bigl(\Norm(a)L_T+\Norm(b)\bigr),\qquad
 R(M,T,\delta)=\frac{N(M,T)}{\delta}.
\]
Then $\Norm(at+b)\le N(M,T)$ and the denominator certificate gives
$\Norm(H_M(t))\le R(M,T,\delta)$.  Since
$R\le(R+1)^2$, both coordinates lie between $-(R+1)$ and $R+1$.
This explicit symmetric rational box proves the lemma without a square root;
the coordinatewise interval construction above is a sharper interchangeable
construction.  The finite-list hull uses only rational minima and maxima and
covers every member of an incomplete sector's candidate list.
\end{proof}

\leanref{CCL.BoxTransport.homographicBox_contains;
CCL.BoxTransport.hullList_contains_of_mem}

Given an old box $B$ and a newly computed box $C$, one may replace $C$ by the coordinatewise intersection with $B$.  The new rational witness lies in both boxes by (7.20)--(7.21), so the intersection is certified non-empty; it is nested in $B$ and has no larger squared diagonal.

\begin{definition}[progressive boxed CCL name]
A finite boxed refinement trace is a finite list of boxed prefix certificates in which each entry refines the preceding one.  A progressive boxed CCL name is specified operationally by an initial certificate and a deterministic refinement procedure which, given the current finite trace and a positive rational $\varepsilon$, returns a finite extension whose last box $B$ satisfies
\[
 \Delta^2(B)\le\varepsilon^2.
\tag{7.26}
\]
Every invocation must terminate and return its finite word, finite box, rational witness, and proofs.
\end{definition}

\begin{proposition}[boxed names supply squared-error processes]\label{prop:boxed-name-process}
Every progressive boxed CCL name supplies a progressive squared-error
enclosure process of the kind used in CCL-2.  At a certificate $(w,B,q)$
its candidates are precisely the non-zero rational values satisfying the raw
prefix predicate, its centre is the fixed rational anchor of $B$, and its
squared radius is $\Delta^2(B)$.  Raw-word extension proves persistence, and
the initial box gives a uniform rational bound for all later centres.

More explicitly, for a positive squared-error target $\tau$, request the
boxed name at
\[
 r(\tau)=\min\!\left(1,\frac{\tau}{4}\right).
\]
Then $r(\tau)>0$ and $r(\tau)^2<\tau$, so the returned squared radius is
strictly below $\tau$.  Consequently two progressive boxed CCL names may be
passed directly to the constructive addition, subtraction, and multiplication
controllers, and to division when a positive denominator-ball separation
certificate is supplied.
\end{proposition}

\begin{proof}
If $0<r\le1$, then $r^2\le r$.  The definition also gives
$r\le\tau/4$, hence
\[
 r^2\le r<4r\le\tau.
\]
If $(v,C,q')$ refines $(w,B,q)$, (7.21) takes every candidate for $v$
to a candidate for $w$, while $C\sqsubseteq B$ preserves its enclosure.
Every later anchor lies in the initial box, whose explicit endpoint bound
therefore bounds its quadrance.  All statements concern the current finite
certificates only.
\end{proof}

\leanref{CCL.BoxedNames.Refines.witness_distance_le}

The procedure is permitted never to emit a terminal marker.  This permits it to capture a value not supplied by a finite Gaussian-rational code---the intended non-rational case---not by postulating a completed complex point, but by producing progressively longer cylinder predicates and nested rational boxes at every requested precision.  Non-termination alone is not asserted to decide irrationality.  No object consisting of all prefixes, all compatible values, or the intersection of all boxes is formed.

The rational witnesses also give explicit numerical approximants.  If $(w',B',q')$ refines $(w,B,q)$, then both $q$ and $q'$ lie in $B$, because $w\pref w'$ and (7.20)--(7.21).  Therefore
\[
 \Norm(q'-q)\le\Delta^2(B).
\tag{7.27}
\]
Along a trace satisfying increasingly small instances of (7.26), the witnesses are thus a directly certified Cauchy approximation process.  For a finite rational name, \autoref{prop:finite-prefix-tower} is the finite version of the same process and terminates at the exact decoder value.

\begin{remark}[interface for constructive root isolation]
A constructive root-isolation argument can couple each boxed prefix certificate with exact polynomial evidence that its current box remains a valid root enclosure, and can answer a precision request by extending the CCL prefix until (7.26) holds.  The returned root is then the refinement procedure itself.  This is the finite-data interface needed for a constructive fundamental theorem of algebra; the root-isolation and root-existence certificates are separate from the CCL encoding developed here.
\end{remark}

\fi

\ifCCLtwo
\section*{Introduction}

CCL-1 constructs the representation consumed here.  Its finite raw parser
associates each ordinary bit prefix with an exact rational projective matrix
and a finite predicate on admissible Gaussian-rational tails.  A completed
canonical block contracts the Hermitian projective quadrance
\[
 \chi^2\bigl(\Proj{x_0}{x_1},\Proj{y_0}{y_1}\bigr)
 =\frac{\Norm(x_0y_1-y_0x_1)}
 {\bigl(\Norm(x_0)+\Norm(x_1)\bigr)
  \bigl(\Norm(y_0)+\Norm(y_1)\bigr)}
\]
of two compatible tails by at most one half.  This is the Gaussian-norm, or
Hermitian, extension of projective quadrance \cite{wildberger-projective}.
The boxed-name interface answers a positive
rational accuracy request with a longer finite word, a nested rational box,
a displayed rational witness, and proofs of enclosure and prefix
persistence.

Those are the only representation-level facts assumed below.  This paper
adds the redundant-output and quantitative sidecar layer required for
productive arithmetic.  It first applies that layer to the shrinking
zero/pole tracks supplied by the RFAP paper, then develops elementary
arithmetic, exponential and logarithmic transforms, operational laws,
limit-free calculus, integration, and finite root families.  The canonical
encoder and its half-open boundary conventions are never changed.

The same one-way boundary also couples the finite approximation-sharpness
runs of RFAP-2 to ordinary CCL prefixes.  One final finite run is reserved
first; intermediate states then reveal longer initial segments of its
window list and of one rational sample's terminating word.  This gives a
prefix-coherent finite tower without asserting an infinite source or
compatibility between independently computed final runs.

As in CCL-1, every numerical object is finite rational data.  A Gaussian
rational is written $x+iy$ with squared norm
$\Norm(x+iy)=x^2+y^2$, and $\gamma=1+i$ has norm two.  A ``name'' is a
terminating request interface returning a finite word and finite
certificates, not an actual infinite stream or a completed point.  The formal
projective coordinate $\omega=\Proj{1}{0}$ is only a pair of rational
coordinates.  All searches range over displayed finite lists or bounded
natural-number intervals.

\section{Wrapping finite FTA and RFAP box tracks}
\label{sec:rfap-wrapper}

The CCL construction does not replace the winding and Sturm calculation of the companion finite FTA and RFAP theorem.  It consumes its output.  For a displayed polynomial, that theorem first performs finite squarefree preprocessing.  Every linear factor supplies an exact Gaussian rational zero.  Every higher factor supplies a degree-indexed family of pairwise separated winding-one tracks:
\[
 (B_{j,0},B_{j,1},\ldots),
 \qquad
 B_{j,n+1}\subseteq B_{j,n},
 \qquad
 \Delta^2(B_{j,n})\le\varepsilon_n^2
\tag{1.1}
\]
operationally, one requested finite refinement at a time.  An exact non-zero
terminal is passed to the finite CCL encoder of CCL-1; zero retains its
distinguished finite tag.  A higher track itself supplies a progressive
squared-error process: at a finite state its candidate predicate is membership
in the displayed box, its ball is based at the rational box anchor with
squared radius $\Delta^2(B)$, and refinement uses the square-root-free
accuracy request established for boxed names in CCL-1.  Separately, every
finite refinement can encode a displayed non-zero rational sample of that
box.  No root is first extracted from the boxes.

For a rational function $F/G$, the Paper-I algorithm first computes and cancels the displayed gcd.  It then returns the pairwise-coprime squarefree multiplicity bases of the reduced numerator and denominator.  A numerator factor of multiplicity $m$ supplies winding $+m$; a denominator factor of multiplicity $m$ supplies winding $-m$.  Constant factors contribute no finite sites, linear factors supply exact rational terminals, and higher factors supply independently shrinking tracks.  The CCL wrapper applies the same terminal/track conversion to both signs and retains the Paper-I cross-separation certificates.

\begin{definition}[CCL site observation]
A CCL site observation is one of two finite records.
\begin{enumerate}
\item An \emph{exact observation} stores a displayed Gaussian rational,
      its direct polynomial evaluation certificate, and either its terminating
      CCL word or the distinguished exact-zero tag.
\item A \emph{track observation} stores one Paper-I winding-one track, its
      current rational box, and one persistent arithmetic cursor whose
      quantitative candidates are members of that track's initial box.
\end{enumerate}
The site index, multiplicity, sign, and separation witnesses remain fields of
the surrounding Paper-I record; the CCL attachment does not recompute them.
\end{definition}

\begin{lemma}[a winding track is a progressive enclosure]
Let $(B_n)$ be a Paper-I winding-one track.  Taking the lower-left anchor of
$B_n$ as centre and $\Delta^2(B_n)$ as squared radius defines a progressive
enclosure process.  A request at positive target $\tau$ returns a nested box
whose ball radius is strictly below $\tau$, and every later candidate remains
a member of the earlier box.
\end{lemma}

\begin{proof}
Use the track's finite refinement operation with the square-root-free request
$r(\tau)=\min(1,\tau/4)$.  Since $0<r(\tau)\le1$,
$r(\tau)^2\le r(\tau)<\tau$.  The track theorem supplies both box nesting
and the requested diagonal bound.  Box membership is the possibility
predicate, so nesting proves persistence without choosing a point from the
boxes.
\end{proof}

\leanref{CCL.TrackProcess.toProcess}

\begin{theorem}[persistent request for one tracked site]
\label{thm:persistent-site-request}
Initialize a joint state at a higher site's initial box $B_0$ and the empty
arithmetic cursor.  Given a later state with box $B$, a positive rational
$\varepsilon$, and $r\in\N$, one finite request returns:
\begin{enumerate}
\item a box $B'\subseteq B$ with
      $\Delta^2(B')\le\varepsilon^2$;
\item an exposed raw word extending the previous exposed word by exactly
      $r$ bits; and
\item a quantitative cylinder of pairwise Hermitian projective quadrance at most
      $\varepsilon^2$.
\end{enumerate}
More explicitly, let $\widehat w$ be the internally committed word, including
the finite unexposed buffer, let $M_{\widehat w}$ be its ordinary parser
matrix, and let $\tau_q$ be the current projective tail chart.  For any two
admitted sidecar candidates $t,u$,
\[
 \chi^2\!\left(M_{\widehat w}\tau_q(t),
                  M_{\widehat w}\tau_q(u)\right)
 \le\varepsilon^2,
 \qquad \Delta^2(B')\le\varepsilon^2.
\tag{1.2}
\]
Every admitted parser image is still an explicit rational candidate of
$B_0$; hence the box and cylinder retain the same finite site owner.
\end{theorem}

\begin{proof}
First ask the cursor for $r$ further bits.  This changes the exposed prefix
but not its originating process.  Refine only the finite process state in the
resulting quantitative mode.  This second operation changes neither exposed
nor buffered bits, parser phase, nor projective chart.

Write $A_q$ for that chart matrix.  The finite matrix
$M_{\widehat w}A_q$ has non-zero determinant, so its explicitly calculated
transport coefficient $K_q$ is positive.  Request sidecar squared radius
$\varepsilon^2/(4K_q)$.  Two admitted candidates in the resulting ball obey
$\Norm(t-u)\le4\rho$ by the two-leg squared-triangle estimate.  Affine
Hermitian projective quadrance is at most this norm, and transport through
$M_{\widehat w}A_q$ contributes at most the factor $K_q$.  This proves the
first inequality of (1.2).  Independently, the winding-track request gives
$B'\subseteq B$ and the second inequality.

The exact transition invariant is preserved by both finite operations.
Therefore every parser image still reconstructs to a candidate of $B_0$.
The track index never changes, so this is also the claimed ownership
certificate.
\end{proof}

\leanref{CCL.PersistentTrack.State.request_cylinder_accurate;
CCL.PersistentTrack.State.candidate_in_initial_box}

\begin{corollary}[finite CCL-wrapped FTA and RFAP]
\label{thm:ccl-wrapped-rfap}
The finite polynomial FTA returns CCL site observations for exactly its
degree-weighted zero sites.  The finite rational-function argument principle
returns them for every reduced zero and pole site, preserving signed
multiplicity, pairwise separation, and
\[
 \operatorname{wind}_{\partial B}(F/G)
   =\deg F_{\rm red}-\deg G_{\rm red}.
\tag{1.3}
\]
Exact sites retain terminating rational CCL data.  Higher sites satisfy
Theorem~\ref{thm:persistent-site-request}.  At any request they may also
return the complete canonical encoding of a displayed non-zero sample in
$B'$; consecutive samples move by squared norm at most the earlier box
diagonal.
\end{corollary}

\begin{proof}
Run Paper I once.  Mapping its finite factor-observation lists is bounded
recursion.  A linear observation already includes direct evaluation at its
displayed rational root, so attach the finite encoder output or exact-zero
tag.  A higher observation retains its existing track, applies the preceding
lemma, and initializes the persistent state of
Theorem~\ref{thm:persistent-site-request}.  The executable box anchor supplies
the optional canonical sample, while box membership bounds consecutive
anchor movement.

No CCL operation changes the Paper-I indices or their multiplicity,
separation, and winding fields.  Applying the same attachment to the reduced
numerator and denominator lists therefore preserves the outer identity
(1.3).  Every construction is a finite map over those displayed lists.
\end{proof}

\leanref{CCL.Paper.polynomialRootStreams;
CCL.PolynomialRootStreams.NonzeroResult.mem_rootIndices;
CCL.PolynomialRootStreams.NonzeroResult.rootRequest_contract;
CCL.Paper.rationalFunctionArgumentPrincipleWithCCL}

\subsection{Finite CCL prefixes for robust sharpness reserves}

RFAP-2 proves a different finite box theorem.  Given a displayed rational
$\tau$ satisfying
\[
 \frac34\le\tau<1,
 \qquad \tau^2-12\tau+9<0,
\tag{1.4}
\]
a scale $T$, and a run length $N$, it computes one proper rational box $B$
and $N$ strictly scale-increasing consecutive-best windows.  Every displayed
rational member of $B$ violates the proposed threshold in every one of those
windows.  The following adapter adds finite CCL observations without entering
the RFAP dependency closure.

\begin{theorem}[prefix-coherent finite sharpness reserve]
\label{thm:ccl-sharpness-reserve}
From $\tau,T,N$ satisfying (1.4), a terminating calculation returns the RFAP
run and box $B$, a displayed $q\in B\setminus\{0\}$, and its complete finite
CCL word $v=\operatorname{enc}(q)$.  For any supplied
\[
       s\le T,\qquad n\le N,\qquad b\le |v|,
\]
the corresponding stage exposes the same box $B$, the first $n$ certified
windows, and the word $w_b=v\mathbin{\upharpoonright}b$.  It satisfies
\[
 \Cyl(w_b,q),\qquad |w_b|=b,
\]
and every exposed window lies beyond $s$, is consecutive best for every
displayed member of $B$, and violates $\tau$ there.

If $s\le s'\le T$, $n\le n'\le N$, and $b\le b'\le|v|$, then the later stage
extends the earlier one:
\[
 B\subseteq B,\qquad w_b\pref w_{b'},\qquad
 (W_0,\ldots,W_{n-1})\pref(W_0,\ldots,W_{n'-1}).
\tag{1.5}
\]
Accordingly a stateful request for a new scale, additional windows, and
additional bits computes the componentwise extension (1.5).  It adds exactly
the requested numbers while they fit in the declared finite reserve, and
saturates at $T$, $N$, and the terminal length $|v|$ otherwise.
\end{theorem}

\begin{proof}
Run the robust RFAP construction once at the final values $T,N$.  Properness
of $B$ gives a nonzero displayed sample by checking its lower-left corner and,
only if that corner is zero, using its upper-right corner.  Convert this
sample coefficientwise to a Gaussian-integer fraction and run the terminating
CCL encoder once.  Define every word and window list by finite list
\texttt{take}.  The cylinder claim is the defining prefix predicate for the
encoder output.  Monotonicity in (1.5) is the elementary fact that a shorter
\texttt{take} is a prefix of a longer one; box refinement is reflexivity.
The uniform consecutive-best and violation claims are inherited directly
from the RFAP box certificate, and its final scale bound implies every smaller
exposed scale bound.  Clipping each request by a natural-number minimum proves
the saturation clauses.
\end{proof}

\leanref{CCL.ApproximationSharpness.reserve;
CCL.ApproximationSharpness.ReservedTower.Stage.cylinder;
CCL.ApproximationSharpness.ReservedTower.Stage.uniformConsecutive;
CCL.ApproximationSharpness.ReservedTower.Stage.uniformViolation;
CCL.ApproximationSharpness.ReservedTower.stage_monotone;
CCL.ApproximationSharpness.ReservedTower.Stage.extend_extends;
CCL.Paper.finiteApproximationSharpnessTower}

Here $\Cyl(w_b,q)$ is a \emph{CCL prefix cylinder}: it says that the finite
encoder word of the displayed rational $q$ begins with $w_b$.  The
error--denominator cylinders used inside the best-approximation proof are
finite inequalities on approximation pairs and are unrelated objects.  The
reserve is also deliberately finite.  Rebuilding it with a larger final
scale or run length may select a different rational source and a different
CCL word; the theorem makes no completed-source claim across those separate
calculations.

\section{Gosper-style arithmetic from finite predicates}
\label{sec:arithmetic}

Every completed input digit or CCL block is a homographic substitution.  A unary operation is represented by
\[
 H(x)=\frac{\alpha x+\beta}{\gamma_0x+\delta},
 \qquad
 M_H=\mat{\alpha}{\beta}{\gamma_0}{\delta}.
\]
When an input block has matrix $A$, consumption is the finite update
\[
 M_H\longmapsto M_HA.
\]

For output, let $\operatorname{Tail}_w(t)$ be the rational predicate produced by the finite parser after input prefix $w$, and let $\operatorname{Branch}_b(y)$ be the rational predicate for output bit $b$ in the current output phase.  The tail predicate also records that every denominator used by $H$ is non-zero.  Emitting $b$ is justified by proving the implication
\[
 \operatorname{Tail}_w(t)\Longrightarrow
 \operatorname{Branch}_b\bigl(H(t)\bigr)
\tag{2.1}
\]
for rational $t$.  At an incomplete CCL sector code, $\operatorname{Tail}_w$ is a finite disjunction indexed by the finite list $\Sigma_q$.  After conjoining denominator non-vanishing, ordered inequalities can be multiplied by the positive squared norms of those denominators.  Hence (2.1) reduces to finitely many rational polynomial implications.  No named collection of all tails is needed.

Binary arithmetic uses the bihomographic state
\[
H(x,y)=\frac{axy+bx+cy+d}{exy+fx+gy+h}.
\]
Consuming a block from one input precomposes in that variable.  Addition,
subtraction, multiplication, and division arise from the corresponding
initial coefficient tuples.  The remainder of this section supplies the
redundant branch semantics, squared-error transport, terminal padding, and
finite margin needed for a productivity theorem; none of those additions
changes the canonical encoder or prefix semantics of CCL-1.

\subsection{Why the canonical branch test is insufficient}

The finite rational encoder may decide the half-open comparison $\Norm(z)<2$ or $2\le\Norm(z)$ exactly.  A transducer receiving only progressive input names cannot do so uniformly.  At a point with $\Norm(z)=2$, every finite enclosure may also contain rational points on both sides.  Any algorithm which emits the unique canonical branch after inspecting only that finite information gives the same answer on two compatible refinements for which opposite answers are required.

The same obstruction occurs when a generic arithmetic operation happens to have an exact rational result.  Opaque names for one chosen solution of $X^2-2=0$ do not themselves supply the algebraic certificate $x^2=2$.  Their product can therefore be certified to arbitrary accuracy as $2$ without the arithmetic transducer ever being able to decide that a terminal marker is justified.  Exact termination is available when the extra equality certificate is supplied; it is not demanded from the numerical names alone.

\subsection{Redundant safe sectors}

For $s\in\mathcal R_8$, put
\[
 D_s(z)\;:\Longleftrightarrow\;\Norm(z-s)<\frac12.
\tag{2.2}
\]
These predicates overlap.  They are used only by the arithmetic output
controller; the disjoint predicate $\operatorname{Sect}_s$ and the canonical
finite encoder remain exactly as in CCL-1.

\begin{lemma}[annulus-free safe inversion]\label{lem:redundant-safe-inversion}
If $D_s(z)$ and $z\ne s$, then
\[
 \Norm\!\left(\frac1{\gamma(z-s)}\right)>1.
\tag{2.3}
\]
If a rational enclosure lies strictly inside the axis-aligned unit cell centred at $s$, then every realization satisfies $D_s$.
\end{lemma}

\begin{proof}
The strict cell inequalities give
\[
 |\Ree(z-s)|<\frac12,
 \qquad |\Imm(z-s)|<\frac12,
\]
and hence $\Norm(z-s)<1/4+1/4=1/2$.  Since $\Norm(\gamma)=2$,
\[
 \Norm\!\left(\frac1{\gamma(z-s)}\right)
 =\frac1{2\Norm(z-s)}>1.
\]
Every comparison is a cross-multiplied rational inequality.  Notice that the non-overlapping unit cells alone would merely move the old ambiguity to their common edges; the overlapping predicates (2.2), rather than ownership of a grid tile, are what remove that ambiguity.
\end{proof}

\leanref{CCL.RedundantSectors.safe_of_openCell;
CCL.RedundantSectors.postInversion_norm_gt_one}

The completed-block contraction proof in CCL-1 uses its upper annulus bound
only through the finite sector-contraction calculation.  Replacing that
derived fact by (2.2) gives the following stronger statement.

\begin{theorem}[redundant completed-block contraction]\label{thm:redundant-block-contraction}
Suppose two displayed values $z,z'$ use the same representative $s$, satisfy
\[
 1\le\Norm(z),\quad 1\le\Norm(z'),
 \qquad D_s(z),\quad D_s(z'),
\]
and use the same displayed radial block.
For each displayed value use the affine inverse tail when its residual is
non-zero, and use the formal projective tail $\omega$ when its residual is
zero.  Then the corresponding completed block of radial norm $L$ obeys
\[
 \chi^2(A_{\lambda,s}t,A_{\lambda,s}t')
 \le \frac{2L}{(1+L)^2}\chi^2(t,t')
 \le\frac12\chi^2(t,t').
\tag{2.4}
\]
with the evident projective reading when either tail is $\omega$.  Hence any
two finite projective realizations of the same $n$ redundant completed blocks
satisfy
\[
 \chi^2(P,Q)\le 2^{-n}.
\tag{2.4a}
\]
Thus crossing the old upper annulus boundary incurs no approximation error,
and an exact residual incurs no exceptional failure of contraction.  Both are
exact redundant representation choices.
\end{theorem}

\begin{proof}
If the residual is non-zero, Equation (2.3) supplies the strict tail lower
bound used in (7.14), while $1\le\Norm(z)$ supplies the output-weight lower
bound used in (7.15).  If it is zero, the block acts on $\omega=[1:0]$ and
reconstructs the displayed representative exactly.  The same weight estimate
proves the affine--$\omega$ case; the $\omega$--$\omega$ numerator is zero.
The one-block half bounds then iterate by structural induction on the supplied
finite list of block realizations, proving (2.4a).  The determinant calculation
and all remaining inequalities in the affine--affine case are unchanged from
the completed-block contraction theorem of CCL-1.
\end{proof}

\leanref{CCL.RedundantBlockContraction.Tail.contraction_half}

\begin{proposition}[hypothesis-free redundant-block semantics]
\label{prop:redundant-block-exact}
For every rational normalized value $y$, every displayed radial factor
$\lambda$, and every representative $s$, with no annulus or safety
hypothesis, the homogeneous decoder identity is
\[
 A_{\lambda,s}
   \Proj{1}{\gamma(y-s)}\sim\Proj{\lambda y}{1},
\tag{2.4b}
\]
where $\lambda y$ is the pre-block value.
\end{proposition}

\begin{proof}
Multiply the displayed matrix by the column
$[1:\gamma(y-s)]$.  Its two coordinates have the same non-zero rational
scalar multiple as $[\lambda y:1]$.  No inequality is used.
\end{proof}

The lower-norm hypotheses in
Theorem~\ref{thm:redundant-block-contraction} are therefore contraction
hypotheses, not exactness hypotheses.  An initial or low-chart digit crossing
the old lower annulus boundary is still exact by
Proposition~\ref{prop:redundant-block-exact}; its uncertainty is simply left
to the sidecar.  A non-zero safe residual has inverse-tail norm greater than
one by (2.3), while a zero residual uses $\omega$, so the following loop is
normalized in either case.

\begin{lemma}[finite projective transport]
\label{lem:finite-projective-transport}
For
$M=\mat{a}{b}{c}{d}$ put
\[
 C(M)=2\bigl(\Norm(a)+\Norm(b)+\Norm(c)+\Norm(d)+1\bigr),
 \qquad
 K(M)=\frac{C(\operatorname{adj}M)^2}{\Norm(\det M)}.
\]
For every displayed $M$ with non-zero determinant and all valid projective
points $P,Q$, direct determinant and weight estimates give
\[
 \chi^2(MP,MQ)\le K(M)\chi^2(P,Q).
\tag{2.4c}
\]
Consequently, if $P,Q$ realize the same $n$ later certified redundant
blocks, then
\[
 \chi^2(MP,MQ)\le K(M)2^{-n}.
\tag{2.4d}
\]
\end{lemma}

\begin{proof}
The alternating numerator in Hermitian projective quadrance is multiplied by
$\det M$.  Applying $\operatorname{adj}M$ to each image and expanding the
four Gaussian-rational terms gives
\[
 \Norm(\det M)W(P)\le C(\operatorname{adj}M)W(MP)
\]
for $W([p:q])=\Norm(p)+\Norm(q)$.  Multiplying the two weight inequalities
and clearing only positive rational denominators proves (2.4c).
Apply the $n$-block estimate (2.4a) and then (2.4c) to obtain (2.4d).
\end{proof}

Here $M$ may contain every already committed unary digit and the one
annulus-free initial or low-chart sector.  Since $M$ is a finite ordinary
decoder prefix, $K(M)$ is an explicitly computed positive rational.  Thus a
fixed early loss of contraction changes only the finite bit budget: later
normalized blocks drive its transported uncertainty below every positive
rational request.  No error symbols are added to the CCL word.

This suffix is not merely hypothetical.  After the one-time initial/low chart
prefix, every finite run of whole controller chunks has a finite proof-sidecar
trace.  Unary normalizations and residual chart changes are neutral trace
steps; every live normalized sector decision contributes its actual ordinary
decoder block.  Reading a final candidate backwards through the trace gives a
projectively coherent realization of precisely that finite block list.
Equation~(2.4d) therefore applies to every pair of its admitted candidates.
This trace is separate from the executable one-bit cursor and does not change
its output type.

\begin{lemma}[redundant cover of the normalized loop]
The eight sector-contraction cases proved in CCL-1 give
\[
 \{z:1\le\Norm(z)\le2\}\subseteq\bigcup_{s\in\mathcal R_8}D_s.
\tag{2.5}
\]
Consequently the loop region has the redundant finite cover
\[
 \{z:1<\Norm(z)\}
 \subseteq
 \{z:2<\Norm(z)\}\cup\bigcup_{s\in\mathcal R_8}D_s.
\tag{2.6}
\]
\end{lemma}

\begin{proof}
The eight finite sector cases give (2.5).  Split a point of norm greater than
one by the rational comparison with two: the high branch is the first member
of (2.6), and the other branch lies in the annulus covered by (2.5).
\end{proof}

\begin{lemma}[explicit redundant-sector margin]
The eight rational cases strengthen the sector estimate to
\[
 1\le\Norm(c)\le\frac{13}{6}
 \quad\Longrightarrow\quad
 \Norm(c-s(c))\le\frac7{16},
\tag{2.7}
\]
where $s(c)$ is the row selected at the rational centre.  The weighted square-root-free inequality
\[
 \Norm(u+v)\le16\Norm(u)+\frac{16}{15}\Norm(v)
\tag{2.8}
\]
holds for all $u,v\in\Qi$.
\end{lemma}

\begin{proof}
Check (2.7) in the eight rational rows of the sector table.  Equation~(2.8)
is the expansion of the non-negative quadrance $\Norm(15u-v)$.
\end{proof}

\begin{lemma}[bounded redundant loop decision]
\label{lem:bounded-loop-decision}
A squared-error ball $\mathcal B(c,\rho)$ with $\rho<1/480$ and
$1\le\Norm(c)$ admits a bounded decision.  If $\Norm(c)\le13/6$, every
realization lies in $D_{s(c)}$.  Otherwise every realization has norm greater
than $2$.  The tests use the exact overlap
\[
 \frac{13}{6}-\frac{32}{15}=\frac1{30},
 \qquad \frac1{480}=\frac{1/30}{16}.
\tag{2.9}
\]
\end{lemma}

\begin{proof}
If $\Norm(c)\le13/6$, apply (2.7) at the centre and (2.8)
to the centre error to obtain one common $D_{s(c)}$ certificate.  Otherwise
(2.9) gives $(16/15)2+16\rho<\Norm(c)$, so (2.8) in the
reverse direction gives $2<\Norm(z)$.
All eight representative checks and both centre comparisons are bounded
finite searches.  Equality at $\Norm(z)=2$ therefore selects a safe sector
rather than blocking.
\end{proof}

\subsection{Terminal unary padding and zero}

\begin{definition}[non-return terminal]
If a residual is exactly zero, its inverse branch uses the projective tail
$\omega$.  A \emph{non-return terminal} is the finite cursor state which
answers later requests by extending the unfinished unary run rather than by
announcing a terminal marker.  The already committed sector matrix is
retained, and unary scaling fixes $\omega$.
\end{definition}

\begin{lemma}[terminal-padding contraction]
\label{lem:terminal-padding}
For an affine tail satisfying $2^r<\Norm(t)$,
\[
 \chi^2(t,\omega)=\frac1{1+\Norm(t)}
 <\frac1{1+2^r}.
\tag{2.10}
\]
Let $M$ be the already emitted finite decoder product.  Its explicit adjugate weight estimate gives a positive rational coefficient $C_M$ such that
\[
 \chi^2(Mt,M\omega)\le\frac{C_M}{1+2^r}.
\tag{2.11}
\]
The exact tail $\omega$ makes the left side zero.  Thus every finite padding request has a certified dyadically shrinking cylinder, although the arithmetic procedure never announces termination.
\end{lemma}

\begin{proof}
The first equality is the definition of Hermitian projective quadrance at
$\omega$.  Apply the finite transport estimate
Lemma~\ref{lem:finite-projective-transport} to obtain (2.11).  The exact
projective tail gives zero directly.
\end{proof}

\begin{corollary}[pairwise terminal-padding cylinder]
If two residual candidates $r,r'$ are both admitted by the small
ball which justifies padding length $q$, their homogeneous tails are each
projectively equivalent to the corresponding affine-or-$\omega$ witness.
The square-root-free two-leg Hermitian-projective-quadrance estimate and
(2.11) give the pairwise
bound
\[
 \chi^2\!\left(M\Proj{1}{\gamma r},
                 M\Proj{1}{\gamma r'}\right)
 \le \frac{4C_M}{1+2^q}.
\tag{2.11a}
\]
\end{corollary}

\begin{proof}
Each residual candidate is within the displayed small quadrance of zero.
Apply terminal-padding contraction to both affine-or-$\omega$ tails, then use
the square-root-free two-leg Hermitian-projective-quadrance inequality.  It
contributes the factor
four.
\end{proof}

\begin{proposition}[initial-low zero padding]
Zero has an analogous initial-low no-return name.  After the initial low bit
and $r+1$ unary ones its finite predicate is
\[
 2^r\Norm(z)<1,
\]
and its affine Hermitian projective quadrance from zero is below $2^{-r}$.
\end{proposition}

\begin{proof}
The unary predicate itself is $2^r\Norm(z)<1$, which is exactly the claimed
affine Hermitian-projective-quadrance bound after dropping the positive denominator
$1+\Norm(z)$.  No later sector block is needed.
\end{proof}

These conventions permit subtraction of two opaque equal names to remain
productive without deciding equality.  They extend only progressive
arithmetic semantics; the canonical finite encoder and its terminal marker
remain unchanged.

\begin{lemma}[finite residual small-or-apart decision]
\label{lem:residual-small-or-apart}
Let a residual ball have centre $c$, squared error $\rho$, and let
$\theta>0$.  If $\rho<\theta/6$, a rational comparison returns one of
\[
 2(\rho+\Norm(c))<\theta,
 \qquad\text{or}\qquad
 0<\frac{\Norm(c)}2-\rho\le\Norm(z)
 \quad\text{for every realization }z.
\]
The first alternative justifies more unary padding; the second justifies
affine inversion.  With $\theta=2^{-(r+1)}$, the returned certificate also
supports $r$ padding bits.  Neither branch decides whether the residual is
zero.
\end{lemma}

\begin{proof}
Test the first strict rational inequality.  If it holds, the squared-triangle
estimate gives $\Norm(z)<\theta$ for every realization.  If it fails and
$\Norm(c)\le2\rho$, then
$\theta\le2(\rho+\Norm(c))\le6\rho<\theta$, an impossibility.
Therefore $0<\Norm(c)/2-\rho$, and the reverse squared-triangle estimate
gives the displayed lower bound for every realization.  With
$\theta=2^{-(r+1)}$, multiplicativity of $\Norm$ gives the inverse-tail
bound.  Exact zero uses $\omega$, so neither branch decides equality.
\end{proof}

\subsection{Squared-error sidecars}

\begin{definition}[squared-error ball]
For $c\in\Qi$ and $\rho\in\Q_{\ge0}$, the squared-error ball
$\mathcal B(c,\rho)$ is the finite predicate
\[
 \Norm(z-c)\le\rho,
 \qquad z\in\Qi.
\tag{2.12}
\]
Here $\rho$ is already a squared error.  A rational box becomes such a ball by taking its lower-left anchor as $c$ and its squared diagonal as $\rho$.
\end{definition}

\begin{lemma}[addition and subtraction of squared-error balls]
If $x\in\mathcal B(a,\rho)$ and $y\in\mathcal B(b,\sigma)$, then
\[
\begin{array}{rcl}
 x+y&\in&\mathcal B(a+b,\;2(\rho+\sigma)),\\
 x-y&\in&\mathcal B(a-b,\;2(\rho+\sigma)).
\end{array}
\tag{2.13}
\]
\end{lemma}

\begin{proof}
Apply $\Norm(u+v)\le2(\Norm(u)+\Norm(v))$ to
$(x-a)+(y-b)$ and to $(x-a)-(y-b)$.
\end{proof}

\begin{lemma}[multiplication of squared-error balls]
If $x\in\mathcal B(a,\rho)$ and $y\in\mathcal B(b,\sigma)$, then
\[
 xy\in\mathcal B\!\left(ab,\;
 4(\rho\Norm(b)+\Norm(a)\sigma)+2\rho\sigma\right).
\tag{2.13a}
\]
\end{lemma}

\begin{proof}
Expand $xy-ab=(x-a)b+a(y-b)+(x-a)(y-b)$.  Two applications of the
squared triangle inequality give the displayed bound, using
multiplicativity of $\Norm$ for each product.
\end{proof}

\begin{lemma}[apartness-certified division]
Suppose $x\in\mathcal B(a,\rho)$, $y\in\mathcal B(b,\sigma)$, and finite
evidence supplies $\delta>0$ with
\[
 y\in\mathcal B(b,\sigma)\Longrightarrow
 \delta\le\Norm(y).
\tag{2.14}
\]
Then $b\ne0$ and
\[
 \frac{x}{y}\in
 \mathcal B\!\left(\frac ab,
 \frac{2(\rho\Norm(b)+\Norm(a)\sigma)}
      {\delta\Norm(b)}\right).
\tag{2.15}
\]
\end{lemma}

\begin{proof}
The certificate at $y=b$ gives $\delta\le\Norm(b)$, hence $b\ne0$.  Now
\[
 \frac{x}{y}-\frac ab
 =\frac{(x-a)b-a(y-b)}{yb};
\]
the numerator is bounded by the addition/subtraction lemma and the
denominator quadrance is at least $\delta\Norm(b)$.  These are finite
rational calculations.  Division is deliberately partial without (2.14):
a zero or pole outcome must remain explicit rather than being hidden by an
invalid inverse.
\end{proof}

\begin{definition}[progressive enclosure process]
A progressive enclosure process consists of finite states $q$, a ball
$\mathcal B_q$, a finite-state possibility predicate $P_q(z)$, and a request
operation
\[
 R(q,\eta),\qquad \eta\in\Q_{>0},
\]
such that
\[
 P_q(z)\Longrightarrow z\in\mathcal B_q,
 \qquad
 P_{R(q,\eta)}(z)\Longrightarrow P_q(z),
 \qquad
 \rho_{R(q,\eta)}<\eta.
\tag{2.16}
\]
Each state contains one displayed rational witness satisfying its current
possibility predicate.  This is data returned by the finite procedure, not an
existence assertion subsequently eliminated by choice.
\end{definition}

\begin{lemma}[candidate apartness becomes ball apartness]
Let $w$ be the displayed witness in a process state with ball
$\mathcal B(c,\rho)$.  If $\delta\le\Norm(w)$ and
$\rho<\delta/12$, then
\[
 \frac{\delta}{8}\le\Norm(z)
 \quad\hbox{for every rational }z\in\mathcal B(c,\rho).
\tag{2.16a}
\]
\end{lemma}

\begin{proof}
Both $w$ and $z$ are in $\mathcal B(c,\rho)$.  The squared triangle
inequality first bounds $\Norm(w-z)$ by $4\rho$.  Its reverse form then gives
$\Norm(w)\le2\Norm(z)+2\Norm(w-z)$, so
$\delta<2\Norm(z)+2\delta/3$.  Weakening the resulting
$\delta/6<\Norm(z)$ to the stated $\delta/8$ keeps the certificate rational
and non-strict.
\end{proof}

Consequently one finite refinement turns candidate-level non-zero evidence
into the whole-ball apartness required for affine division.  Persistence of
the possibility predicate keeps that construction available after every
later refinement.

\begin{lemma}[weighted quadrance inequality]
For every rational $m>0$ and $u,v\in\Qi$,
\[
 \Norm(u+v)\le(m+1)\Norm(u)+\frac{m+1}{m}\Norm(v),
\tag{2.16b}
\]
\end{lemma}

\begin{proof}
Expand the non-negative quadrance $\Norm(mu-v)$ and rearrange.
\end{proof}

\begin{lemma}[finite preparation of the loop lower bound]
Suppose every current candidate has norm at least $L>1$.  Put
$m=2/(L-1)$ and $\kappa=(m+1)/m$.  A refinement below
\[
 \frac{L-\kappa}{m+1}
\]
has a rational centre of norm at least one.  Hence a common target below this
number and $1/480$ supplies the hypotheses of the bounded redundant loop
decision in one finite request.
\end{lemma}

\begin{proof}
The chosen $m$ has $\kappa<L$.  If the refined centre had norm below one, apply
(2.16b) to the displayed witness written as centre plus error.  The requested
error bound would make the witness norm strictly smaller than $L$, contrary
to its candidate certificate.
\end{proof}

\begin{proposition}[exact residual rebasing]
After a safe sector has been emitted, rebase the finite process at the
state on which that decision was made.  Every later candidate is therefore
still admitted by the safe state.  Exact translation by the representative
produces the residual process $r=z-s$ and preserves the explicit strict upper
bound already present in the safe-ball inequality.  If residual apartness is
returned, one further calculated refinement makes the whole residual ball
apart and division constructs
\[
 t=\frac{\gamma^{-1}}r=(\gamma r)^{-1}.
\]
The retained residual upper bound supplies a uniform rational lower bound
$L>1$ for every candidate $t$.
\end{proposition}

\begin{proof}
Persistence gives the first assertion.  Translation is exact and preserves
the proved upper bound.  Candidate-to-ball apartness and
apartness-certified division construct $t$.  Multiplicativity of quadrance
turns the residual upper bound into the asserted lower bound for $t$, after
which the loop-preparation lemma applies.  Thus this is an exact chart
change, not accumulated numerical error.
\end{proof}

Each process also carries a rational $M\ge0$ bounding the squared norm of every
displayed centre.  This is finite persistence data, not a selected completed
complex value.  Binary operations use pairs of operand states; for example,
the new finite predicate for addition says that the displayed rational
candidate is $x+y$ for two candidates satisfying the two current operand
predicates.  The implication in (2.16) is then preserved componentwise.

\begin{proposition}[closed operation requests and centre bounds]
For a desired output squared error $\tau>0$, let $M_x,M_y$ be the two centre
bounds.  The following common operand requests suffice:
\[
\begin{array}{c|c}
\text{operation}&\text{operand squared-error target }\eta\\ \midrule
x\mathbin{\pm}y&\tau/4\\[2pt]
xy&\displaystyle
 \frac{\tau}{\tau+4(M_x+M_y)+2}\\[7pt]
x/y&\displaystyle
 \frac{\tau\delta^2}{4(M_x+M_y+1)}
 \quad\text{when }\delta\le\Norm(y).
\end{array}
\tag{2.17}
\]
The transported centre bounds are respectively
\[
 2(M_x+M_y),\qquad M_xM_y,\qquad M_x/\delta.
\tag{2.18}
\]
Thus the output of each operation is itself a progressive finite enclosure
process and may be used as an operand of another operation.
\end{proposition}

\begin{proof}
Addition and subtraction follow immediately from (2.13).  For
multiplication put $K=4(M_x+M_y)+2$.  The displayed request satisfies
$0<\eta<1$ and $\eta K<\tau$, so (2.13a) is strictly below $\tau$.
For division, the numerator in (2.15) is at most
$2\eta(M_x+M_y)$ and its denominator is at least $\delta^2$; the strict
inequality $2(M_x+M_y)<4(M_x+M_y+1)$ gives the result.  Applying the same
estimates to the rational centres gives (2.18).  Componentwise refinement
proves persistence, so no search is hidden in any request.
\end{proof}

\begin{lemma}[affine inverse-chart sensitivity]
For $z,w\ne s$,
\[
 \Norm\!\left(
 \frac1{\gamma(z-s)}-\frac1{\gamma(w-s)}\right)
 =\frac{\Norm(z-w)}
 {2\Norm(z-s)\Norm(w-s)}.
\tag{2.19}
\]
\end{lemma}

\begin{proof}
Subtract the two fractions, use $\Norm(\gamma)=2$, and use
multiplicativity of quadrance.
\end{proof}

The lemma explains why a scalar $\varepsilon^2$ alone is insufficient: an
affine inversion additionally needs residual separation.  When separation
has not yet been proved, the controller retains a projective enclosure and
emits unary padding rather than forming an uncertified affine quotient.

\subsection{Bihomographic operations and productivity}

\begin{definition}[initial bihomographic arithmetic states]
The four arithmetic operations start from the following bihomographic
coefficient tuples:
\[
\begin{array}{c|rrrr|rrrr}
 &a&b&c&d&e&f&g&h\\ \toprule
 x+y&0&1&1&0&0&0&0&1\\
 x-y&0&1&-1&0&0&0&0&1\\
 xy&1&0&0&0&0&0&0&1\\
 x/y&0&1&0&0&0&0&1&0.
\end{array}
\tag{2.20}
\]
\end{definition}

Input blocks are consumed by finite precomposition in the relevant variable.  Between block consumptions, (2.13)--(2.18) compute an output enclosure.  The controller refines the input, using the closed targets (2.17), until one redundant output predicate holds for the entire enclosure.  It then emits that exact output bit or sector code, substitutes the corresponding inverse output matrix into the bihomographic state, and continues.

The executable realization may transport the two input balls and the current
output chart directly instead of retaining all eight coefficients at every
step.  The two descriptions are propositionally linked.  Every candidate of
the progressive addition, subtraction, multiplication, or
apartness-certified division process is projectively equal to the value of
the corresponding initial row of (2.20).  Moreover, for an arbitrary
bihomographic state $T$, input decoder matrices $L,R$, and output matrix $O$,
the finite coefficient updates satisfy the following identity.

\begin{lemma}[bihomographic transition identity]
For every bihomographic state $T$, input decoder matrices $L,R$, and output
matrix $O$,
\[
 \bigl(O\circ(T\circ(L,R))\bigr)(x,y)
   =O\bigl(T(Lx,Ry)\bigr).
\tag{2.21}
\]
For an emitted decoder block $A$, take $O=\operatorname{adj}(A)$.
\end{lemma}

\begin{proof}
Expand both sides as homogeneous matrix products and use associativity of
finite matrix multiplication.  No division or normalization is involved.
\end{proof}

Thus direct enclosure transport and Gosper's matrix bookkeeping are certified
views of the same arithmetic transition; productivity is proved for the
former while (2.21) proves the exact substitution interpretation.

\begin{definition}[exact cursor invariant]
The exact cursor invariant uses a projective interpretation of each
quantitative mode.  If $v$ is a candidate of the mode, put
\[
\begin{array}{c|c}
q&\tau_q(v)\\ \midrule
\text{initial}&\Proj{v}{1}\\[2pt]
\text{initial-low with count }m&\Proj{(\gamma^m)^{-1}v}{1}\\[2pt]
\text{regular with count }k&\Proj{\gamma^k v}{1}\\[2pt]
\text{residual}&\Proj{1}{\gamma v}.
\end{array}
\tag{2.22}
\]
Let $\widehat w$ be the exposed raw prefix followed by the unexposed
remainder of its one current finite output chunk.  If $q_0$ is the original
finite process state, the invariant is
\[
 P_q(v)\Longrightarrow
 \exists z\in\Qi\quad
 P_{q_0}(z)\quad\hbox{and}\quad
 M_{\widehat w}\tau_q(v)\sim\Proj{z}{1}.
\tag{2.23}
\]
The existential witness here is returned by the finite candidate proof.  It
is not a root or a selected completed value.
\end{definition}

\begin{theorem}[closure of progressive enclosure arithmetic]
Given two progressive enclosure processes, addition, subtraction, and
multiplication construct progressive enclosure processes.  Division does so
from a positive denominator-apartness certificate.  Every positive rational
squared-error request is answered by one finite request to each operand.
\end{theorem}

\begin{proof}
Use the product finite state and possibility predicate described in (2.16).
The three ball-transport lemmas and the closed-operation request proposition
prove strict output accuracy and the new centre bounds.  Componentwise refinement proves
persistence, while each product state carries the displayed pair of operand
witnesses.
\end{proof}

\begin{theorem}[finite totality of the redundant controller]
At a normalized loop state, one refinement below $1/480$ returns either a
sound unary branch or a sound redundant sector.  At a residual state, one
refinement below $\theta/6$ returns either a sound residual bound below
$\theta$ or positive apartness sufficient for affine inversion.  The initial
chart is likewise prepared by a fixed finite refinement.  None of these
decisions tests equality of an opaque result.
\end{theorem}

\begin{proof}
At a loop state whose displayed centre has norm at least one, request radius
below $1/480$ and apply the bounded loop-decision lemma; its ball-level
certificate is sound for every candidate admitted by the refined finite
state.  For a requested residual target $\theta$, request radius below
$\theta/6$ and apply the residual small-or-apart lemma.  In the
small branch, taking $\theta=2^{-(r+1)}$ supplies the projective loop-tail
witness for $r$ unary padding bits; in the apart branch the positive lower
bound supplies the exact hypothesis of (2.15).  Each construction is a
finite function call on rational data, and no completed stream or completed
complex value is formed.

For the initial chart, first refine below $1/480000$.  If the displayed centre
has norm below one, emit the initial-low selector and one low-unary bit; the
weighted estimate (2.8) gives the rescaled candidate bound
$\Norm(\gamma z)<43/20<13/6$.  If the centre has norm at least one, the
overlapping decision following (2.9) emits either an exponent-zero low block
or the first high-unary bit.  In a subsequent low state, the retained
$43/20$ bound and (2.16b), with $m=1000$, keep the next sufficiently refined
centre strictly below $13/6$.  A centre of norm at least one therefore emits
a safe sector, while a centre below one emits another low-unary bit and
preserves the same bound.  Exact zero follows this latter branch forever.
The ordinary loop and residual modes then use the two decisions above.
\end{proof}

\begin{theorem}[productive raw cursor and unchanged grammar]
The controller decisions assemble into a finite-state raw-bit cursor.  From
every state and for every $n\in\N$, bounded recursion returns an extension
containing exactly $n$ further bits without retracting the old prefix.
Parsing every committed whole chunk with the ordinary CCL-1 parser gives the
syntax phase recorded by the controller.  Thus the accuracy sidecar adds no
symbol to the raw CCL grammar.
\end{theorem}

\begin{proof}
Every emitted chunk is non-empty: initial-low and sector chunks may be
buffered, but the public cursor exposes them one bit at a time.  Structural
recursion on $n$ consequently proves both prefix extension and the exact
length increase $n$.  A finite case analysis of the initial-high,
initial-low, loop, and residual modes shows that each chunk takes the
syntax-only raw parser from its stored phase to the phase stored by the next
mode.  Associativity of finite-list append lifts this local fact through the
one-bit buffer and the same structural recursion, proving the committed-word
parser invariant.
\end{proof}

\begin{theorem}[exact semantics of the productive arithmetic cursor]
\label{thm:ccl-arithmetic-productive}
Equation~(2.23) holds after every finite request of the productive cursor for
addition, subtraction, multiplication, or apartness-certified division.
Consequently every internally committed extension has the exact semantics of
the ordinary raw parser for the original arithmetic process.

After a redundant completed block, a non-zero residual $r$ has the exact
affine tail $(\gamma r)^{-1}$, while $r=0$ has the exact projective tail
$\omega$.  More generally, $\Norm(r)<2^{-(q+1)}$ gives either the exact
$\omega$ tail or an exact affine tail of norm greater than $2^q$ and
Hermitian projective quadrance at most $(1+2^q)^{-1}$ from $\omega$.
\end{theorem}

\begin{proof}
For exact semantics, the unary branches use the two cancellation identities
\[
 \gamma^{k+1}(\gamma^{-1}v)=\gamma^kv,
 \qquad
 (\gamma^{m+1})^{-1}(\gamma v)=(\gamma^m)^{-1}v.
\]
A sector branch uses the hypothesis-free homogeneous identity (2.4b).
Translation makes its next residual candidate exactly $v-s$.  Small-residual
padding leaves $\Proj{1}{\gamma(v-s)}$ unchanged.  In the apart branch,
normalizing the affine inverse by the already emitted unary count gives
\[
 \Proj{\gamma^k(\gamma^k)^{-1}(\gamma r)^{-1}}{1}
 \sim\Proj{1}{\gamma r},
\]
so the silent chart change is exact before the next ordinary loop step.
For a normalized-loop sector outcome, the loop-domain lower certificate and
the selected whole-ball safe-sector certificate apply simultaneously to
every pair of actual post-emission residual candidates.  Replacing their
homogeneous tails by the equivalent affine-or-$\omega$ tails therefore
specializes (2.4) directly to that live controller transition.
Non-zero determinants preserve projective validity, allowing these local
equivalences to compose with the previous parser product.  Associativity of
finite-list append then lifts the result through the buffer and induction on
the requested bit count, proving (2.23).  Finally, (2.10) proves the displayed
no-return bound.  Thus productivity and exactness do not depend on recognizing
an exact rational result.
\end{proof}

\leanref{CCL.ProgressiveTransducer.Named.add_productive;
CCL.ProgressiveTransducer.Named.mul_productive;
CCL.ProgressiveTransducer.Named.div_productive}

\section{Computing exponential and logarithm as CCL names}

The objective of this section is an operation in the CCL representation,
not merely a separate routine producing rational Taylor approximants.  An
arithmetic CCL name consists of the productive raw-word cursor of
Section~\ref{sec:arithmetic}
together with its rational enclosure sidechannel.  We construct a total
operator
\[
 \mathsf{ExpName}:\mathsf{CCLName}\longrightarrow\mathsf{CCLName}
\]
and certified logarithm operators on their stated domains.  Thus a user can
continue requesting ordinary CCL bits of the exponential result, while each
request also carries enough finite rational accuracy to justify those bits.
The sidechannel is essential: the theorem does not claim that a bare raw word,
with no quantitative enclosure information, suffices to schedule all input
refinements.

The arithmetic interface permits these operations to be constructed without
adding completed real or complex points.  We follow the finite recurrences
used by Collins~\cite{collins}, but change their logical role.
An approximation is a finite program with a natural-number request argument.
On a displayed input $n$ it computes finite data
\[
 (a_n,\rho_n)
 \tag{3.1}
\]
consisting of the Gaussian rational $a_n$ and rational squared allowance
$\rho_n$; for displayed $n\le m$ the same program proves
\[
 \Norm(a_m-a_n)\le\rho_n;
 \tag{3.2}
\]
and for every displayed rational $\tau>0$ it computes an index
$r(n,\tau)\ge n$ with $\rho_{r(n,\tau)}<\tau$.  Subscripts here are a request
schema for running that finite program, not a stored infinite stream.  The
possibility predicate at index $n$ is not a named set or an intersection: it
is the constructive proposition that returns a finite witness $m\ge n$ and
an equality between the supplied rational value and $a_m$.  Thus (3.2) puts
every such witnessed later sample in the displayed ball
$\mathcal B(a_n,\rho_n)$.

This approximation process is the internal accuracy sidechannel from which
the output CCL controller makes its finite decisions.  This is the meaning of
\emph{error} in this section.  It is a checked finite
tail-error certificate between rational samples.  We neither consume an
infinite sequence as input nor measure distance to a silently postulated
analytic value.  The words ``exponential'' and ``logarithm'' first name the
following recurrence interfaces.  Their productivity proofs do not require a
completion.  The analytic-identification layer is instead phrased as finite
operational equality: two programs are equal when every positive rational
accuracy request computes later displayed tails which are pairwise that
close.  This retains the intended function laws without silently adjoining
completed non-rational values.

\subsection{The local exponential recurrence}

For $x\in\Qi$ satisfying $\Norm(x)\le1$, define the finite Horner terms
\[
 E_{0,j}(x)=1,
 \qquad
 E_{n+1,j}(x)=1+\frac{x}{j}E_{n,j+1}(x),
 \qquad
 e_n(x)=E_{n+1,1}(x).
 \tag{3.3}
\]
Only natural recursion on the displayed depth $n$ occurs.  In particular,
(3.3) is not notation for an infinite continued expression.  Direct
quadrance estimates prove
\[
 \Norm(e_m(x)-e_n(x))\le \frac{10}{4^n}
 \qquad(n\le m),
 \tag{3.4}
\]
and, for two unit-disk inputs,
\[
 \Norm(e_n(x)-e_n(y))\le16\Norm(x-y).
 \tag{3.5}
\]
The proof bounds every finite inner Horner tail and uses
$j^{-2}\le1/4$ for $j\ge2$; induction then propagates the displayed
quarter-contraction.  Consequently $e_n(x)$, with allowance $10/4^n$, is a
finite approximation process in the sense of (3.1)--(3.2).

For a progressive input process, take at stage $n$ the rational witness of
its $n$th calculated refinement and evaluate (3.3) to depth $n$.  Two such
input witnesses at stages $n\le m$ have quadrance at most $4/4^n$.  Combining
this fact with (3.4)--(3.5) and the square-root-free two-leg inequality
gives the deliberately factored allowance
\[
 \rho_n^{\exp}=\frac{2(10+16\cdot4)}{4^n}.
 \tag{3.6}
\]
No large decimal numeral is required by the executable proof.

\begin{theorem}[the local exponential is CCL-computable]
For every arithmetic CCL input name whose candidates at the displayed
starting state satisfy $\Norm(x)\le1$, equations (3.3)--(3.6) construct an
arithmetic CCL output name.  Every positive rational squared-error request
terminates, and the output's ordinary CCL cursor returns every requested
finite number of further bits without retracting its previous prefix.  Its
exact-cursor certificate identifies every admitted parser tail with a
rational candidate of the calculated exponential output process at that same
finite state.
\end{theorem}

\begin{proof}
Use (3.6) as the output name's process radius.  A bounded scan of successive
powers of $1/4$ finds the first displayed index below the requested target.
The scan fuel is the binary length of the computed rational ceiling of the
scale-to-target quotient plus one.  Equations (3.4)--(3.6) prove
persistence and ball containment, so this finite process is a valid accuracy
sidechannel for a CCL name.  Feed that process to the unchanged redundant CCL
controller of Section~\ref{sec:arithmetic}.
Theorem~\ref{thm:ccl-arithmetic-productive} returns
every requested output bit, and its exact-cursor induction proves that the
committed ordinary raw prefix and the exponential enclosure describe the
same finite output state.  This last step is what turns the Horner recurrence
into computation \emph{with the CCL}, rather than leaving it as an unrelated
rational approximation algorithm.
\end{proof}

\leanref{CCL.ProgressiveExp.productive}

\subsection{The local positive-real logarithm recurrence}

For a rational $x$ in the overlapping chart
\[
 \frac12\le x\le2,
 \qquad
 z=\frac{x-1}{x+1},
 \qquad w=z^2,
 \tag{3.7}
\]
direct cross-multiplication gives $-1/3\le z\le1/3$ and
$0\le w\le1/9$.  Put
\[
 L_{0,j}(w)=1,
 \qquad
 L_{n+1,j}(w)=1+\frac{2j-1}{2j+1}wL_{n,j+1}(w),
 \qquad
 \ell_n(x)=2zL_{n,1}(w).
 \tag{3.8}
\]
Again every term is a finite rational Horner calculation.  The coefficient
and chart bounds imply
\[
 \bigl|\tfrac{2j-1}{2j+1}w\bigr|^2\le\frac14,
 \qquad
 \Norm(\ell_m(x)-\ell_n(x))\le\frac1{4^n}
 \quad(n\le m).
 \tag{3.9}
\]
A denominator-cleared comparison of two displayed inputs in the same chart
also gives the finite sensitivity estimate
\[
 \Norm(\ell_n(x)-\ell_n(y))\le64(x-y)^2.
 \tag{3.10}
\]
For progressive chart inputs, the witness movement allowance is again
$4/4^n$.  Combining (3.9)--(3.10) therefore gives
\[
 \rho_n^{\log}=\frac{2(1+64\cdot4)}{4^n}.
 \tag{3.11}
\]

\begin{theorem}[the local positive-real logarithm is CCL-computable]
Suppose every candidate of the displayed arithmetic CCL input state is real
and has real coordinate in $[1/2,2]$.  Equations (3.7)--(3.11) construct an
arithmetic CCL output name for the local logarithm recurrence.  Every
positive rational accuracy request terminates, and its ordinary CCL cursor
returns every requested finite number of further bits with exact same-state
parser semantics.
\end{theorem}

\begin{proof}
The chart hypothesis applies to every explicitly selected later witness.
Equations (3.9)--(3.11) prove the approximation fields, and the same
bounded quarter-power scan answers every rational accuracy request.  Use
that process as the output CCL name's accuracy sidechannel and start the
unchanged controller of Section~\ref{sec:arithmetic}.  Its productivity
theorem supplies every
requested ordinary CCL bit, while exact-cursor semantics tie all admitted
parser tails to candidates of the logarithm process at the same state.  Thus
the result is a local CCL-to-CCL operation, not merely a list of rational
Horner values.
\end{proof}

\leanref{CCL.ProgressiveLog.productive}

\subsection{Automatic finite range and branch preparation}

\begin{lemma}[finite exponential range reduction]
Exponential needs no branch information.  At the current input state let
$c$ and $r$ be the rational centre and squared radius, and compute
\[
 B=2\bigl(r+\Norm(c)\bigr).
 \tag{3.12}
\]
Every current candidate $x$ satisfies $\Norm(x)\le B$.  A bounded
quarter-power scan computes $k\in\N$ with
\[
 \frac{B}{4^k}<1.
 \tag{3.13}
\]
Hence every candidate of the exactly scaled process $x/2^k$ lies in the unit
disk.  Apply the local CCL construction there and then apply CCL
multiplication as exact repeated squaring exactly $k$ times.  Thus the global
finite CCL name is defined by
\[
 \mathsf{ExpName}(X)=
 \operatorname{Sq}^{k}\!\left(\mathsf{ExpLocal}(X/2^k)\right);
\tag{3.14}
\]
then every input to $\mathsf{ExpLocal}$ lies in its unit-disk chart.
\end{lemma}

\begin{proof}
The squared triangle inequality gives (3.12).  A bounded scan of the finite
list $B,B/4,\ldots,B/4^k$ finds (3.13).  Since division by $2^k$ divides
quadrance by $4^k$, the scaled process lies in the unit disk.  Local
exponentiation followed by exactly $k$ CCL squarings gives (3.14).
This is an operational range reduction, not an equality between
independently named completed values.
\end{proof}

\begin{definition}[positive-real logarithm certificate]
A positive-real logarithm certificate consists of a rational $L>0$ and the
proof
\[
 P_q(x)\Longrightarrow \Imm(x)=0\ \hbox{ and }\ L\le\Ree(x)
\tag{3.15}
\]
for every candidate of the displayed state.
\end{definition}

\begin{lemma}[finite dyadic logarithm chart preparation]
Given a positive-real logarithm certificate, refine the state once to
squared radius below
\[
 \left(\frac{L}{16}\right)^2.
 \tag{3.16}
\]
If $c$ is the new rational ball centre, a displayed candidate witness and
the square comparison in (3.16) prove
\[
 \frac L2\le\Ree(c),
 \tag{3.17}
\]
so $\Ree(c)>0$ without extracting any common point from the process.

A bounded binary scan now computes a direction $d\in\{\mathsf{up},
\mathsf{down}\}$ and exponent $k\in\N$ whose rational dyadic factor $s$ puts
the centre safely inside the overlap:
\[
 \frac34\le s\Ree(c)\le\frac32,
 \qquad
 s=2^k\ \hbox{or}\ s=2^{-k}.
 \tag{3.18}
\]
For values above one, the scan fuel is $\lceil\Ree(c)\rceil$; below one it is
$\lceil1/\Ree(c)\rceil$.  The first adjacent power certificate is adjusted
by at most one exponent to obtain (3.18).  Equations (3.16)--(3.18) give
\[
 s\frac{L}{16}\le\frac14,
 \tag{3.19}
\]
and therefore every candidate in the refined ball satisfies
\[
 \Imm(sx)=0,
 \qquad
 \frac12\le\Ree(sx)\le2.
\tag{3.20}
\]
This proves the candidate-wise local chart required by (3.7), including
inputs on dyadic boundaries.  The rational ball supplies the deviation
bound; the separate certificate (3.15), rather than the disk geometry,
proves that all process candidates are real.  No claim is made that every
irrelevant geometric point of the two-dimensional disk is real.
\end{lemma}

\begin{proof}
The displayed candidate witness and (3.16) imply (3.17) by the reverse
squared triangle estimate.  The bounded upward or downward power-of-two scan
gives adjacent dyadic bounds, adjusted by at most one exponent to obtain
(3.18).  Applying the same estimate to every candidate gives (3.19) and
(3.20).  All comparisons are rational and the scan has the displayed finite
fuel.
\end{proof}

\begin{definition}[global positive-real logarithm name]
Generate the local CCL name
$\ell_n(2)$ once and combine it by CCL arithmetic.  Writing
$\mathsf{LogTwoName}=\mathsf{LogLocal}(2)$, define
\[
 \mathsf{LogName}(X)=
 \begin{cases}
   \mathsf{LogLocal}(X/2^k)+k\mathsf{LogTwoName},
      &d=\mathsf{down},\\
   \mathsf{LogLocal}(2^kX)-k\mathsf{LogTwoName},
      &d=\mathsf{up}.
 \end{cases}
\tag{3.21}
\]
Here each occurrence of ``$\log$'' denotes its displayed finite recurrence
inside a productive CCL name.  Multiplication by the displayed integer $k$
and addition are the already proved finite CCL operations.  Hence automatic
positive-real charting does not leave the representation.  No exact
non-rational constant $\log2$ is introduced.
\end{definition}

\subsection{A certified complex branch chart}

\begin{definition}[local complex logarithm recurrence]
The same finite recurrence works over Gaussian rationals once its reduced
argument is supplied directly.  For $z\in\Qi$ with
\[
 \Norm(z)\le\frac19,
 \qquad w=z^2,
 \tag{3.22}
\]
define $L_{n,j}(w)$ exactly as in (3.8) and put
$\ell_n^{\mathrm C}(z)=2zL_{n,1}(w)$.
\end{definition}

\begin{lemma}[local complex logarithm bounds]
For $n\le m$ and two reduced arguments $z,v$ in the disk (3.22),
\[
 \Norm\bigl(\ell_m^{\mathrm C}(z)-\ell_n^{\mathrm C}(z)\bigr)
     \le \frac5{4^n},
 \qquad
 \Norm\bigl(\ell_n^{\mathrm C}(z)-\ell_n^{\mathrm C}(v)\bigr)
     \le64\Norm(z-v).
 \tag{3.23}
\]
The deliberately loose factor $5$ avoids square roots and keeps every bound
in the Gaussian-rational field.  Synchronizing finite input refinement and
Horner depth therefore gives the progressive allowance
\[
 \rho_n^{\mathrm C\log}=\frac{2(5+64\cdot4)}{4^n}.
\tag{3.24}
\]
\end{lemma}

\begin{proof}
The bounds in (3.23) follow by the same finite Horner induction as
(3.9)--(3.10), now using multiplicativity of Gaussian quadrance.  The loose
factor $5$ avoids square roots and keeps every allowance rational.
Combining tail and input movement gives (3.24) by the two-leg inequality.
\end{proof}

\begin{definition}[certified complex logarithm chart]
A one-chart logarithm for an arithmetic CCL input $Y$ consists of the
following displayed finite data.  Choose $a\in\Qi$; for every finite input state prove a
fixed rational lower quadrance $\lambda>0$ for every candidate of $aY+1$;
and at the starting state prove, conservatively for every independently
admitted pair of candidates $y_1,y_2$, that
\[
 \Norm\!\left(\frac{ay_1-1}{ay_2+1}\right)\le\frac19.
 \tag{3.25}
\]
The independent pair is required because finite enclosure arithmetic may
refine the two occurrences of $Y$ separately.  One calculated refinement
below a fixed rational fraction of $\lambda$ converts candidate-wise
separation into apartness of the entire displayed denominator ball.  Hence
progressive division can construct the reduced-argument name
\[
 Z=\frac{aY-1}{aY+1}.
\]
Finally supply a progressive correction name $C$.
The correction is the branch choice.  It may itself be non-rational and is
therefore a finite-request name, not a postulated completed constant.  The
safe-chart interface does not guess a cut, an anchor, or an integer branch
index.  Safety and productivity alone permit an arbitrary correction.  To
certify that it is genuinely logarithmic, we additionally require the finite
operational branch condition
\[
 \mathsf{ExpName}(C)\simeq \mathsf{Exact}(a^{-1}),
 \qquad a\ne0,
\tag{3.26a}
\]
where $\simeq$ is the finite-request equivalence defined below.  This is
insensitive to the expected $2\pi i$ branch shifts but excludes an unrelated
correction.
\end{definition}

\begin{definition}[charted complex logarithm name]
For a certified chart, prepare the division
$Z=(aY-1)/(aY+1)$ and put
\[
 \mathsf{ComplexLogName}(Y;a,C)
   =\mathsf{ComplexLogLocal}(Z)+C.
\tag{3.26}
\]
Prepared division, the local Gaussian recurrence, and correction addition
are all CCL operations, so this is an arithmetic CCL output name.
\end{definition}

\begin{definition}[CCL-computable operation]
An operation on arithmetic CCL names is \emph{CCL-computable} when, for every
rational $\tau>0$ and $n\in\N$, a finite request returns an output state of
squared radius below $\tau$ and exactly $n$ further ordinary CCL bits.  The
same request must return an exact-cursor certificate: every admitted tail,
acted on by the ordinary parser for the committed word and pending finite
chunk, is projectively equivalent to an explicitly witnessed candidate of
that output state.
\end{definition}

\begin{theorem}[global exponential is CCL-computable]
For every arithmetic CCL name $X$, the name $\mathsf{ExpName}(X)$ of (3.14)
is CCL-computable.
\end{theorem}

\begin{proof}
Equations (3.12)--(3.14) reduce exponential to the local process and a
finite recursion of CCL multiplication nodes.  Consequently range reduction
does not leave the CCL representation: scaling, local exponential, and every
successive squaring produce arithmetic CCL names.  Refine the resulting
process at $\tau$, then invoke Theorem~\ref{thm:ccl-arithmetic-productive} at
that same state for the $n$ bits and its exact-parser certificate.
\end{proof}

\begin{theorem}[positive-real logarithm is CCL-computable]
If $X$ carries the positive-real certificate (3.15), then
$\mathsf{LogName}(X)$ of (3.21) is CCL-computable.  Its preparatory
refinement, dyadic exponent, local chart proof, and correction coefficient
are displayed finite data.
\end{theorem}

\begin{proof}
The dyadic chart-preparation lemma reduces the input to (3.7).  The local
logarithm theorem, the independently generated $\mathsf{LogTwoName}$, and
finite CCL arithmetic construct (3.21).  Refine its output process at
$\tau$ and invoke the arithmetic cursor theorem at that state for the
requested bits and exact semantics.
\end{proof}

\begin{theorem}[charted complex logarithm is CCL-computable]
If $X$ carries the certified complex chart data (3.25)--(3.26a), then
$\mathsf{ComplexLogName}(X;a,C)$ is CCL-computable.
\end{theorem}

\begin{proof}
Candidate apartness becomes whole-ball apartness after one calculated
refinement, so progressive division constructs $Z$.  Equations
(3.22)--(3.24) construct its local Gaussian logarithm, and CCL addition adds
the supplied correction.  The generic output refinement and exact cursor
theorem answer the joint accuracy-and-bit request.  Denominator preparation,
Horner evaluation, correction addition, and cursor execution are all finite.
\end{proof}

\begin{corollary}[certified exponential-minus-logarithm composition]
\label{cor:g-productive}
Let $X$ and $Y$ be progressive Gaussian-rational input processes, and give
$Y$ the positive-real certificate (3.15).  Define, operationally,
\[
 \mathsf{GName}(X,Y)=\mathsf{ExpName}(X)-\mathsf{LogName}(Y).
 \tag{3.27}
\]
For every rational $\tau>0$ and $n\in\N$, a finite request returns a rational
output ball of squared radius below $\tau$, exactly $n$ exposed ordinary CCL
bits, and the same-state exact-cursor certificate of the preceding theorem.
The same conclusion holds when $Y$ instead carries the complex chart data
(3.25)--(3.26) and $\mathsf{LogName}(Y)$ is replaced by
$\mathsf{ComplexLogName}(Y;a,C)$.
\end{corollary}

\begin{proof}
Apply the already constructed exponential and certified logarithm processes,
their finite progressive subtraction, and then the generic combined request
at target $\tau$ and count $n$.
The same proof, replacing $\mathsf{LogName}$ by
$\mathsf{ComplexLogName}$ and using (3.25)--(3.26), gives the charted
complex version.
\end{proof}

\leanref{CCL.TranscendentalOutput.exp_ccl_computable}
\leanref{CCL.TranscendentalOutput.log_ccl_computable}
\leanref{CCL.TranscendentalOutput.complex_log_ccl_computable}

Positive-real logarithm is not a claim about a global complex branch.  The
complex extension above likewise does not manufacture branch information:
it consumes the explicit scale, denominator clearance, reduced-disk proof,
and correction.  Subject to such chart data at every logarithm node, finite
expression trees built from arithmetic, exponential, and logarithm remain
productive by composition.
Corollary~\ref{cor:g-productive} provides one such $g$ node.  This is the
precise domain condition needed before using
repeated $g$ constructions to specify further transcendental operations.

\subsection{Finite operational equality and function laws}

\begin{definition}[finite operational equality]
For approximation programs $A=(a_n,\rho_n)$ and $B=(b_n,\sigma_n)$, write
$A\simeq B$ when, from any displayed starting indices $p,q$ and any rational
$\tau>0$, a terminating comparison computes $r\ge p$ and $s\ge q$ such that
\[
 \Norm(a_m-b_n)<\tau
 \qquad\text{for every displayed }m\ge r\text{ and }n\ge s.
 \tag{3.28}
\]
All quantifiers here range over natural indices and Gaussian-rational
samples.  The comparison is data: it returns the synchronizing indices.
\end{definition}

\begin{proposition}[operational equality is an equivalence relation]
Finite operational equality is reflexive, symmetric, and transitive.
\end{proposition}

\begin{proof}
Reflexivity uses identical indices, and symmetry swaps the returned pair.
For transitivity, request quarter of the target from each comparison,
synchronize the two intermediate indices, and apply the two-leg quadrance
inequality.  Every index adjustment is a maximum of finitely displayed
natural numbers.
\end{proof}

Applying the canonical nested finite sampling path of an arithmetic CCL name
therefore gives a finite operational equality for finite-request CCL
programs.

The raw process accuracy field alone is not enough for this conclusion:
independently requested refinements may follow unrelated branches.  We
therefore attach the following finite certificate.

\begin{definition}[coherent arithmetic CCL name]
A coherence certificate assigns each finite
state $s$ a non-negative rational semantic allowance $\eta(s)$ and a positive
rational request modulus $R(\tau)$ such that every two explicitly possible
rational readings $u$ at $s$ and $v$ at $t$ satisfy
\[
 \Norm(u-v)\le 4\bigl(\eta(s)+\eta(t)\bigr),
 \tag{3.28a}
\]
and, whenever a positive process request $r$ obeys $r<R(\tau)$,
\[
 \eta(\mathsf{refine}(s,r))<\tau.
\tag{3.28b}
\]
Both $R$ and the refined state are computed finite data.  Exact rational
names have allowance zero.
\end{definition}

\begin{proposition}[coherence yields operational equality]
Any two starting states of one coherent process give operationally equal CCL
names.  Coherence is preserved by addition, negation, subtraction, and
multiplication.
\end{proposition}

\begin{proof}
Use (3.28b) twice with a sufficiently small rational fraction of the target,
then apply (3.28a).  Addition and negation use the squared triangle estimate.
Multiplication uses the squared-error product polynomial in the two semantic
allowances and their rational centre bounds, followed by the closed request
of (2.17).
\end{proof}

Every displayed-sample approximation inherits a
certificate from its later-sample estimate.  The certificates are preserved
by the preceding proposition.  The local and range-reduced exponential,
repeated squaring, the positive logarithm and its dyadic correction, and a
complex logarithm with a coherent explicit branch correction all carry this
certificate.  No common completed value is selected.

\begin{lemma}[local exponential at zero]
For every displayed finite depth $n$,
\[
 E_{n,j}(0)=1,\qquad e_n(0)=1
\tag{3.29}
\]
and therefore $\mathsf{ExpLocal}(0)\simeq\operatorname{exact}(1)$.
\end{lemma}

\begin{proof}
Direct induction through (3.3) gives $E_{n,j}(0)=1$, hence $e_n(0)=1$.
Identical displayed samples satisfy every operational tolerance.
\end{proof}

\begin{theorem}[local exponential addition law]
For
$x,y\in\Qi$ with $\Norm(x)\le1$, $\Norm(y)\le1$, and
$\Norm(x+y)\le1$, let the left approximation display
$e_n(x)e_n(y)$ and the right approximation display $e_n(x+y)$.  Then
\[
 \mathsf{ExpLocal}(x+y)\simeq
 \mathsf{ExpLocal}(x)\mathsf{ExpLocal}(y).
\tag{3.30}
\]
The product on the right is the actual stateful CCL multiplication process.
\end{theorem}

\begin{proof}
The constructive proof does not appeal to uniqueness of an infinite
series.  Bounded induction identifies (3.3) exactly with the finite
polynomial $\sum_{k\le n}x^k/k!$: expose each nested Horner contribution,
rewrite its consecutive denominator product as the executable finite
factorial, and add the final monomial.  A finite Cauchy-product calculation
and the binomial theorem identify every coefficient through the displayed
truncation degree.  The remaining finite triangular tail satisfies, at the
synchronized odd depth $2k+1$,
\[
 \Norm\bigl(e_{2k+1}(x)e_{2k+1}(y)-e_{2k+1}(x+y)\bigr)
 \le {16\over4^k}.
 \tag{3.30a}
\]
The independent tail allowances on both sides are also explicit, so a
bounded quarter-power scan computes the synchronization indices required by
(3.28) for every displayed $\tau>0$.  Thus (3.30) is an operational law,
not merely a pointwise estimate at one depth.
Moreover, the formal statement is now connected to the actual arithmetic
expression tree.  A separate finite bridge follows the stateful
\texttt{Named.mul} process to one displayed multiplication state, bounds its
two component radii, and compares every later product reading with the
pointwise product approximation.  Hence the proved left side of (3.30) is
literally the CCL multiplication of the two progressive local exponential
names, rather than a product stream merely wrapped in a fresh name.  The
local exponential also respects operational equality of arbitrary
unit-disk input names, by applying the $16\delta$ finite Horner sensitivity
bound at one synchronized depth and absorbing the two finite truncation
tails.
\end{proof}

\begin{lemma}[congruence of finite repeated squaring]
For every displayed $k\in\N$, coherent names $X\simeq Y$ satisfy
\[
 \mathsf{SquareIterate}(k,X)\simeq
 \mathsf{SquareIterate}(k,Y).
\tag{3.30d}
\]
\end{lemma}

\begin{proof}
Canonically square
the displayed samples of any approximation and transport its later-sample
allowance through the two-factor product estimate.  For a coherent CCL name,
a computed request makes both operand readings of the actual stateful
self-multiplication close to one canonical sampled reading; this compares
\texttt{Named.mul $X$ $X$} with that canonical square approximation.  Finite
induction on $k$ proves the result through the actual repeated
self-multiplication tree.
\end{proof}

\begin{lemma}[finite local logarithm reconstruction bound]
For $\Norm(z)\le1/9$, put
\[
 \ell_n^{\mathrm C}(z)=2\sum_{j=0}^n{z^{2j+1}\over2j+1},
 \qquad
 P_n(z)=\sum_{r=0}^{2n+2}{(\ell_n^{\mathrm C}(z))^r\over r!}.
\]
Then the executable estimate
\[
 \Norm\left(P_n(z)-{1+z\over1-z}\right)
 \le {118\over4^n}.
\tag{3.30b}
\]
\end{lemma}

\begin{proof}
Finite coefficient differentiation and substitution prove that the
coefficients of $P_n$ agree through degree $2n+2$ with those of
$1+2z+2z^2+\cdots$.  A non-negative finite coefficient majorant bounds the
remaining terms.  Substitution of $\Norm(z)\le1/9$ gives (3.30b).
\end{proof}

\begin{theorem}[displayed rational exponential--logarithm inverse]
For a displayed rational chart input $y\in\Qi$, put
$z=(ay-1)/(ay+1)$.  If the chart certificate gives $\Norm(z)\le1/9$ and
the anchor (3.26a) supplies $q_n$ with
$\Norm(q_n-a^{-1})<4^{-n}$, then
\[
 \Norm\bigl(P_n(z)q_n-y\bigr)
 \le {2\bigl(50+118\Norm(a^{-1})\bigr)\over4^n}.
 \tag{3.30c}
\]
where $50$ is a uniform quadrance bound for $P_n(z)$.  Consequently
\[
 \mathsf{Exp}(\mathsf{ComplexLog}(y;a,C))\simeq y
\tag{3.31}
\]
for the displayed rational input.
\end{theorem}

\begin{proof}
The chart certificate proves the denominator non-zero, and direct
Gaussian-rational cancellation gives $(1+z)/(1-z)=ay$.  Combine the local
reconstruction bound with the anchor error and the uniform bound $50$ to get
(3.30c).  A bounded quarter-power scan supplies the indices for (3.31).
Every sample is rational; the finite round-trip theorem compares controller
resampling with the originating approximation, so the emitted CCL name is
operationally equivalent to exact $y$.  No completed logarithm value is
asserted.
\end{proof}

\begin{theorem}[arithmetic congruence]
Operational equality of coherent names is a congruence for addition and
multiplication.  Multiplication is operationally commutative and associative.
For every displayed $k$,
\[
 \mathsf{SquareIterate}(k,X)\mathsf{SquareIterate}(k,Y)
 \simeq \mathsf{SquareIterate}(k,XY).
\tag{3.30e}
\]
\end{theorem}

\begin{proof}
Equal-index sampled sums
and products give common rational approximations for coherent names.
Comparing actual stateful arithmetic with these approximations proves
two-input congruence for addition and multiplication, and proves
commutativity and associativity of multiplication, without selecting
completed operand values.  Reassociation and induction on $k$ give (3.30e).
\end{proof}

\begin{theorem}[automatic exponential addition for rational inputs]
For all displayed $x,y\in\Qi$,
\[
 \mathsf{ExpName}(\operatorname{exact}x)
 \mathsf{ExpName}(\operatorname{exact}y)
 \simeq
 \mathsf{ExpName}(\operatorname{exact}(x+y)).
\tag{3.30g}
\]
The three automatic range exponents need not agree.
\end{theorem}

\begin{proof}
Dyadic exponent reconciliation is finite for exact Gaussian-rational data.
Let $D_0(z)=z$ and $D_{k+1}(z)=D_k(z)/2$, and put
$E_k(z)=\mathsf{SquareIterate}(k,\mathsf{LocalExp}(D_k(z)))$ whenever the
displayed local unit-disk certificate holds.  Applying (3.30) to
$D_{k+1}(z)+D_{k+1}(z)=D_k(z)$ proves $E_{k+1}(z)\simeq E_k(z)$; finite
induction permits any larger admissible exponent.  Taking the finite maximum
of three independently supplied exponents and using (3.30e) therefore gives
\[
 E_{k_x}(x)E_{k_y}(y)\simeq E_{k_s}(x+y)
 \tag{3.30f}
\]
for exact rational $x,y$ and any three admissible $k_x,k_y,k_s$.
For the public automatic entry point, actual multiplication of the exact
dyadic scale and exact input is operationally equivalent to their displayed
rational product.  Local exponential congruence and (3.30d) transport this
bridge through the whole automatically selected squaring tree, proving
(3.30g).
\end{proof}

\begin{theorem}[local exponential addition for coherent names]
If the possible finite readings of coherent names $X$, $Y$, and $X+Y$ lie in
the unit disk, then
\[
 \mathsf{LocalExp}(X)\mathsf{LocalExp}(Y)
 \simeq \mathsf{LocalExp}(X+Y).
\tag{3.30h}
\]
\end{theorem}

\begin{proof}
For arbitrary coherent progressive inputs, the local addition law is now
complete.  The actual sum is transported to an equal-index sampled wrapper;
its possible finite readings inherit the actual sum's unit-disk certificate.
At one computed odd depth the exact polynomial bound (3.30a) is instantiated
with the two displayed operand readings.  Explicit dynamic Horner tails
compare actual multiplication of the local exponential outputs with the
dynamic sum-Horner name.  A second finite scan synchronizes that name with
local exponential evaluation on the sampled-sum wrapper, after which local
exponential congruence returns to the actual sum process.  No completed
operand values are selected.
\end{proof}

\begin{theorem}[operational exponential laws]
\label{thm:operational-exponential-laws}
For coherent progressive CCL inputs,
\[
 \mathsf{ExpName}(X)\mathsf{ExpName}(Y)
 \simeq \mathsf{ExpName}(X+Y),
\tag{3.30i}
\]
and $\mathsf{ExpName}(\operatorname{exact}0)\simeq
\operatorname{exact}1$.  The total exponential also respects operational
equality.  Every synchronization index is computed from finite certificates
and the requested positive rational tolerance.
\end{theorem}

\begin{proof}
The same law now holds for the total automatic transform.  The three range
reductions for $X$, $Y$, and $X+Y$ may select different exponents.  Their
finite maximum determines three explicit halving chains.  Each chain is
proved operationally equivalent to the canonical sampler multiplied by the
same displayed inverse dyadic.  Distributivity at those rational samples,
the common-scale unit-disk bound, (3.30h), and finite repeated-squaring
congruence give (3.30i).  This is the actual automatic
process on all three sides, not an idealized stream and not a theorem about
selected completed values.  The exact-zero specialization is also connected
to the finite identity sample:
\[
 \mathsf{ExpName}(\operatorname{exact}0)
 \simeq \operatorname{exact}1.
 \tag{3.30j}
\]
The common-refinement argument, local Horner sensitivity, and repeated
squaring congruence also transport operational equality through the total
exponential.
\end{proof}

\begin{corollary}[operational inverse and non-vanishing]
For every coherent name $X$,
\[
 \mathsf{ExpName}(-X)\mathsf{ExpName}(X)\simeq1,
 \qquad
 \mathsf{ExpName}(X)\mathsf{ExpName}(-X)\simeq1.
\tag{3.30k}
\]
In particular $\mathsf{ExpName}(X)\not\simeq0$.
\end{corollary}

\begin{proof}
Addition has operational commutativity, associativity, zero, and additive
inverse laws; multiplication has commutativity, associativity, unit, zero,
and both distributivity laws.  The exponential addition law and additive
inverse law give (3.30k).  Exact zero
and exact one fail the comparison (3.28) already at the displayed squared
tolerance $1/2$.  If $\mathsf{ExpName}(X)\simeq0$, multiplication congruence
with the independently constructed $\mathsf{ExpName}(-X)$ would make that
same product equivalent both to exact zero and to exact one, contradicting
this finite separation.  No completed exponential value or global
apartness oracle is asserted.
\end{proof}

\leanref{CCL.TranscendentalCoherence.automaticExp_add}
\leanref{CCL.ProgressiveCalculusCongruence.exponential_addition}

\begin{lemma}[M\"obius reconstruction sensitivity]
If $\Norm(r),\Norm(s)\le1/9$, then
\[
 \Norm\left({1+r\over1-r}-{1+s\over1-s}\right)
 \le {81\over4}\Norm(r-s).
\tag{3.31c}
\]
\end{lemma}

\begin{proof}
Direct cancellation gives
$(1+r)/(1-r)-(1+s)/(1-s)=2(r-s)/((1-r)(1-s))$.
The square-root-free estimates $\Norm(1-r),\Norm(1-s)\ge4/9$, followed by
multiplicativity of quadrance, give the result.
\end{proof}

\begin{lemma}[charted-input sensitivity]
Put $T_a(x)=(ax-1)/(ax+1)$.  If the displayed denominators for $x$ and $y$
have quadrance at least $\delta>0$, then
\[
 \Norm(T_a(x)-T_a(y))
 \le {4\Norm(a)\over\delta^2}\Norm(x-y).
\tag{3.31d}
\]
\end{lemma}

\begin{proof}
Put the two fractions over their common denominator.  The numerator reduces
to $2a(x-y)$; use the two denominator lower bounds and multiplicativity of
quadrance.
\end{proof}

\begin{lemma}[exact synchronized chart reconstruction]
Let $y_n$ be one rational candidate read from a prepared input state and set
$r_n=T_a(y_n)$.  Then
\[
 {1+r_n\over1-r_n}a^{-1}=y_n.
\tag{3.31e}
\]
The samples $r_n$ form a progressive reduced-input name with an explicit
quarter-power allowance and inherit the chart's $1/9$ disk certificate.
\end{lemma}

\begin{proof}
The charted-input sensitivity lemma transports the input coherence allowance
to the $r_n$.  The disk certificate is candidate-wise and hence applies to
each displayed $y_n$.  Equation~(3.31e) is direct rational cancellation.
\end{proof}

\begin{lemma}[asynchronous quotient comparison]
Write $F_a(x,y)=(ax-1)/(ay+1)$.  Suppose $\Norm(z)\le M$ and every displayed
denominator has quadrance at least $\delta>0$.  Then
\[
 \Norm(F_a(x,y)-T_a(z))
 \le {4(\Norm(a)^2M+\Norm(a))\over\delta^2}
       \bigl(\Norm(x-z)+\Norm(y-z)\bigr).
\tag{3.31f}
\]
Consequently, for a coherent name $Y$, the stateful prepared quotient is
operationally equivalent to the synchronized name:
\[
 R_a^{\rm state}(Y)
 \simeq \mathsf{Name}\bigl((T_a(y_n))_{n\ge0}\bigr).
\tag{3.31g}
\]
\end{lemma}

\begin{proof}
Common-denominator cancellation and the squared triangle estimate give
(3.31f).  For a requested tolerance, compute a smaller coherence budget.
One bounded scan drives the prepared-division request below its corresponding
input request, and a second reaches a diagonal source index with the same
allowance.  Apply (3.31f) to the two source witnesses and the diagonal
witness.  This returns the finite synchronization indices required by
(3.31g), without selecting a common value.
\end{proof}

\begin{theorem}[anchored exponential--logarithm inverse]
\label{thm:anchored-exp-log-inverse}
Let $Y$ be a coherent arithmetic CCL name and let $(a,C)$ be an anchored
logarithm chart for it: the chart includes the displayed denominator and
reduced-disk certificates, a productive correction name $C$, and the finite
exponential anchor for $a^{-1}$.  Then the actual CCL processes satisfy
\[
 \mathsf{ExpName}(\mathsf{ComplexLog}(Y;a,C))\simeq Y.
\]
The comparison algorithm works for every requested positive rational squared
tolerance and does not select a completed logarithm value.
\end{theorem}

\begin{proof}
The total automatic exponential respects operational equality between
arbitrary coherent inputs.  Applying (3.30i) in reverse to the actual charted-logarithm
sum, then applying the branch certificate, gives
\[
 \mathsf{ExpName}(\mathsf{ComplexLog}(Y;a,C))
 \simeq
 \mathsf{ExpName}(\mathsf{LocalLog}(R_a(Y)))\,\operatorname{exact}(a^{-1}).
 \tag{3.31a}
\]
At this stage the arbitrary-input inverse obligation reduces precisely to the
local charted reconstruction on the right of (3.31a); exponent splitting,
branch correction, and total-transform congruence are no longer involved.
The finite reconstruction below discharges that obligation.  Branch changes
remain explicit: they are transported only through checked overlapping
finite enclosures, and a zero-discrepancy logarithm law requires a displayed
alignment certificate.  Nothing here defines a global single-valued
logarithm.

The sharper finite Horner estimate puts every displayed local-log sample in
the exponential unit disk.  At reduced-sample depth $n$, use the finite
program $P_n(r_n)$ with exponential depth $2n+1$.  The reconstruction bound
(3.30b), together with the explicit truncation and input-sensitivity
allowances, gives
\[
 \mathsf{ExpName}(\mathsf{LocalLog}(R_a(Y)))
       \operatorname{exact}(a^{-1})
 \simeq
 \mathsf{Name}\left(
   \left({1+r_n\over1-r_n}a^{-1}\right)_{n\ge0}
 \right).
 \tag{3.31b}
\]
The name on the right has a closed quarter-power error schedule; it denotes a
finite approximation program, not a selected limit.

The M\"obius reconstruction and charted-input sensitivity lemmas transport
operational equality through the last chart.  One finite
denominator-preparation state is projected back to its input component; at
depth $n$ the reference reads one rational candidate $y_n$ there and emits
$T_a(y_n)$.  The exact synchronized reconstruction lemma gives (3.31e), and
the asynchronous quotient comparison gives (3.31g).  Exact reconstruction
turns the scaled diagonal chart into a finite
subsampling of the prepared input sampler.  The explicit input-coherence
certificate rebases that prepared state to the original input state.  Combining
this arithmetic comparison with (3.31a)--(3.31e) proves the full anchored
round trip for arbitrary coherent progressive inputs,

\[
 \mathsf{ExpName}(\mathsf{ComplexLog}(Y;a,C))\simeq Y.
 \tag{3.31h}
\]

This completes the finite comparison algorithm and hence the proof.
\end{proof}

\leanref{CCL.ComplexLogCorrectness.actualExp_chartedLog_named_equivalent}

\subsection{Finite logarithm branches and certified compositions}

\begin{definition}[anchored logarithm chart and transition]
An anchored logarithm chart is finite data: the rational M\"obius chart,
its denominator and reduced-disk certificates, a productive correction CCL,
and the operational anchor needed for the exponential round trip.  Given two
such charts $A,B$ over the same input, define their transition by the actual
CCL subtraction tree
\[
 \Delta(A,B)=\mathsf{LogName}(Y;B)-\mathsf{LogName}(Y;A).
 \tag{3.31i}
\]
\end{definition}

\begin{theorem}[finite logarithm branch calculus]
\label{thm:finite-log-branches}
For anchored charts $A,B,C$ over the same coherent input,
$\mathsf{ExpName}(\Delta(A,B))\simeq1$ and
\[
 \Delta(A,A)\simeq0,\qquad
 \Delta(A,B)+\Delta(B,A)\simeq0,\qquad
 \Delta(A,B)+\Delta(B,C)\simeq\Delta(A,C).
\]
\end{theorem}

\begin{proof}
First, $\mathsf{ExpName}(\Delta(A,B))\simeq1$.  Transitions to the same chart
are operationally zero, opposite transitions add to zero, and three supplied
charts satisfy
\[
 \Delta(A,B)+\Delta(B,C)\simeq\Delta(A,C).
 \tag{3.31j}
\]
Finite induction gives the corresponding statement for any displayed finite
chain of charts.
\end{proof}

\begin{theorem}[branch-safe logarithm product law]
For supplied branches of $X$, $Y$, and a process
operationally equivalent to the actual product $XY$, the discrepancy between
$\log(XY)$ and $\log X+\log Y$ exponentiates to exact one.
\end{theorem}

\begin{proof}
Expand the discrepancy using the total exponential addition and inverse
laws.  The three anchored inverse theorems and the supplied product-input
equivalence reduce its exponential to exact one.  The discrepancy remains an
explicit CCL name rather than a silently chosen completed multiple of
$2\pi i$.
\end{proof}

\begin{corollary}[aligned logarithm product law]
When a separate
finite alignment certificate proves that this actual CCL discrepancy is
zero, finite addition and cancellation give the stronger guarded law
\[
 \mathsf{LogName}(XY;C)\simeq
   \mathsf{LogName}(X;A)+\mathsf{LogName}(Y;B).
 \tag{3.31j'}
\]
\end{corollary}

\begin{proof}
The product-input equivalence and the zero discrepancy are both fields of the
certificate; neither is inferred from the fact that the discrepancy has unit
exponential.
\end{proof}

\begin{theorem}[productivity of finite certified expressions]
These operations are closed under a finite certified expression package.
Each node stores its constructed coherent CCL output and an explicit natural
node count.  Leaves and exact rationals need no side condition; division
stores denominator apartness and a logarithm node stores its anchored chart.
Addition, negation, multiplication, exponential, and those guarded nodes
compute their output processes recursively.  For every rational $\tau>0$
and bit count $n$, the expression computes one refined finite state whose
ball has squared radius below $\tau$, emits exactly $n$ further raw CCL bits,
and returns the same-state exact parser certificate.  Its coherence sidecar
also computes a positive rational request modulus.  Thus arbitrary finite
composition preserves the CCL-with-error computational contract; no semantic
value is assigned to the expression outside its finite requests.
\end{theorem}

\begin{proof}
Proceed by induction on the finite expression tree.  Leaves use exact or
supplied coherent names.  Arithmetic nodes use progressive closure;
exponential and logarithm use their CCL-computability theorems; division and
logarithm consume the stored apartness or chart certificate.  Compose the
rational request moduli and invoke the exact cursor at the resulting state.
\end{proof}

\begin{lemma}[quarter-power scan cost]
Every quarter-power refinement scan used by these constructions has a closed
finite cost bound.  Starting at index $j$ with non-negative scale $S$ and
positive target $\tau$, the returned index is at most
\[
 j+\left(\log_2\!\left\lceil S/\tau\right\rceil+1\right),
 \tag{3.31k}
\]
where the logarithm and ceiling are executable operations on displayed
natural and rational data.
\end{lemma}

\begin{proof}
Run the scan with exactly the displayed fuel.  Repeatedly multiplying the
denominator by four puts the scale below $\tau$ before the fuel is exhausted.
This is bounded natural recursion, not an unbounded convergence argument.
Expression nodes compose their displayed request moduli; no uniform cost
independent of those finite node certificates is claimed.
\end{proof}

\leanref{CCL.LogBranchGeometry.logarithm_product_discrepancy_unit}
\leanref{CCL.LogBranchGeometry.logarithm_product_of_alignment}
\leanref{CCL.CertifiedExpression.Expression.bitstream_and_error}

\subsection{Projective rational-function evaluation}

\begin{proposition}[finite projective rational-function evaluation]
\label{prop:projective-rational-evaluation}
Let $f,g\in\Qi[T]$ and let $X$ be a certified expression.  Finite Horner
recursion constructs coherent numerator and denominator expressions and
retains the result projectively as
\[
 [\mathsf{Eval}(f,X):\mathsf{Eval}(g,X)].
 \tag{3.31l}
\]
This construction respects operational equality and specializes to the
ordinary affine or reciprocal-affine chart whenever the corresponding
displayed denominator apartness certificate is supplied.
\end{proposition}

\begin{proof}
Horner evaluation respects operational equality.  At a displayed rational
input it is operationally equivalent to the exact displayed polynomial
value, by finite induction through multiplication and addition.  No
denominator decision is required for the projective pair.  The ordinary
affine expression $f(X)/g(X)$ is exposed only from supplied finite apartness
of the denominator.  Around a pole, supplied numerator apartness instead
exposes the reciprocal chart $g(X)/f(X)$; denominator equivalence to exact
zero together with that numerator certificate is the finite pole record.
Thus rational functions compose with CCL computations without adjoining an
affine infinity or completing either coordinate.
\end{proof}

\leanref{CCL.ProjectiveRationalEvaluation.evaluate_exact}

\subsection{Finite divided differences and differential laws}

The differential assertions used here do not quantify over an infinitesimal,
take a limit, or first identify a function with a completed set of ordered
pairs.  They begin with finite syntax and an exact two-input factor.

\begin{definition}[power divided-difference kernel]
For a natural exponent define
\[
 K_n(x,y)=
 \begin{cases}
  0,&n=0,\\
 x^{n-1}+yK_{n-1}(x,y),&n>0.
 \end{cases}
\]
\end{definition}

\begin{lemma}[power divided difference]
For every displayed natural $n$ and every $x,y\in\Qi$,
\[
 x^n-y^n=(x-y)K_n(x,y),
 \qquad K_n(x,x)=n x^{n-1}.
 \tag{3.31m}
\]
\end{lemma}

\begin{proof}
Finite induction in $\Qi$ proves the first identity in (3.31m).  The second
equality is obtained only after the first polynomial factor has
been constructed; no division by $x-y$ occurs, including on the diagonal.
\end{proof}

\begin{corollary}[finite polynomial divided differences]
\label{lem:finite-divided-differences}
Every displayed finite polynomial $p$ has an executable kernel $\Delta p$
satisfying
\[
 p(x)-p(y)=(x-y)\Delta p(x,y),
 \qquad \Delta p(x,x)=p'(x),
 \tag{3.31n}
\]
where $p'$ is the executable finite formal derivative.
\end{corollary}

\begin{proof}
Take the coefficient-weighted sum of the power kernels.  Distribute the exact
factor $(x-y)$ before evaluating on the diagonal; the diagonal identity is
then the finite coefficient formula for $p'$.
\end{proof}

\begin{proposition}[finite rational-expression kernels]
Signed powers, finite Laurent term lists, and a rational expression $f/g$
with the required displayed non-zero denominator certificates have finite
divided-difference kernels.  Product, composition, reciprocal, and quotient
kernels obey their usual exact rules.
\end{proposition}

\begin{proof}
Negative powers use explicit Gaussian-rational inverse cancellation.
Substitute the already constructed kernels into the finite product,
composition, reciprocal, and quotient identities and clear the certified
non-zero denominators.  These identities define the derivative data; they do
not appeal to an ambient analytic derivative.
\end{proof}

\leanref{CCL.DividedDifference.polynomialKernel_diagonal}

\begin{theorem}[finite differential law for exponential]
The finite exponential program has an exact two-input secant factorization,
valid on the diagonal, whose first coefficient is the exponential program.
After subtraction of the linear term, the remainder has an explicit rational
bound proportional to $\Norm(h)^2$.  The factor and remainder are productive
CCL names, giving the operational law $D\exp(x)=\exp(x)$.
\end{theorem}

\begin{proof}
The local exponential has an especially economical factor.  Put
\[
 R_n(z)=E_{n,2}(z),
 \qquad S_n(z)={1\over2}E_{n,3}(z),
\]
using the finite Horner programs of (3.3).  At every displayed depth,
including $z=0$,
\[
 e_n(z)-1=zR_n(z),
 \qquad R_{n+1}(z)-1=zS_n(z),
 \qquad e_{n+1}(z)-1-z=z^2S_n(z).
 \tag{3.31o}
\]
Thus $R_n$ is the executable relative exponential usually denoted
$\operatorname{exprel}$, but it is defined by the middle Horner program rather
than by the unsafe quotient $(e_n(z)-1)/z$.  Likewise $S_n$ is a finite second
relative remainder.  On $\Norm(z)\le1$ both carry explicit quarter-power
squared-error schedules and uniform rational centre bounds, so each feeds the
ordinary productive CCL controller and emits every requested finite prefix.
Direct induction gives $R_n(0)=1$.

For a displayed base $x$ and increment $h$, the product computation
$e_n(x)e_n(h)$ therefore has exact finite secant factor
\[
 e_n(x)e_n(h)-e_n(x)
 =h\,e_n(x)R_n(h),
\]
and, one depth later, exact quadratic remainder
\[
 e_{n+1}(x)e_{n+1}(h)-e_{n+1}(x)-h e_{n+1}(x)
 =h^2e_{n+1}(x)S_n(h).
 \tag{3.31p}
\]
The displayed unit-disk bounds make the word ``quadratic'' quantitative in
the paper's quadrance convention:
\[
 \Norm\bigl(e_{n+1}(h)-1-h\bigr)\le \Norm(h)^2,
 \qquad
 \Norm(\text{translated remainder})\le10\Norm(h)^2.
\]
Since $\Norm(h)$ is already squared ordinary magnitude, these are the usual
$O(|h|^2)$ remainder bounds expressed without square roots.
On the diagonal $h=0$ the displayed secant factor is exactly $e_n(x)$.
The already proved finite operational addition law identifies the product
computation with the direct computation of $e(x+h)$ at arbitrary requested
rational accuracy.  This is the constructive differential law
$D\exp(x)=\exp(x)$ for the CCL-computable exponential: an exact finite
factorization, a shrinking quadratic sidecar, and an operational bridge to
the public exponential process.
\end{proof}

\begin{theorem}[finite differential law for an anchored logarithm]
On a supplied anchored chart, the local logarithm program has an exact
two-input secant factorization valid on the diagonal.  Its first coefficient
is operationally the reciprocal program, its quadratic remainder has an
explicit rational bound, and all three programs are productive CCL names.
Thus $D\log(x)=x^{-1}$ on that chart.
\end{theorem}

\begin{proof}
For
$\Norm(u)\le1/4$ define the displayed finite programs
\[
 L_n(u)=\sum_{k=0}^n{(-u)^k\over k+1},
 \qquad
 Q_n(u)=-\sum_{k=0}^n{(-u)^k\over k+2},
 \qquad
 D_n(u)=\sum_{k=0}^n(-u)^k.
\]
All three sums are natural-number recursions over Gaussian rationals.  The
local logarithm sample is $uL_n(u)$, and exact finite rearrangement gives
\[
 L_{n+1}(u)-1=uQ_n(u),
 \qquad uL_{n+1}(u)-u=u^2Q_n(u).
 \tag{3.31q}
\]
The proved uniform bound $\Norm(Q_n(u))\le4$ therefore gives the explicit
finite estimate
\[
 \Norm\bigl(uL_{n+1}(u)-u\bigr)\le4\Norm(u)^2.
\]
More fundamentally, let $F_n(u)=uL_n(u)$.  The finite-polynomial divided
difference constructs a two-input expression $\Delta_n(x,y)$ before any
diagonal substitution, and proves
\[
 F_n(x)-F_n(y)=(x-y)\Delta_n(x,y),
 \qquad \Delta_n(u,u)=D_n(u).
\]
Thus $D_n$ is not merely a separately chosen approximation to the expected
answer: it is exactly the diagonal value of the secant factor for the same
finite local-log program.
The derivative program has the division-free geometric certificate
\[
 (1+u)D_n(u)=1-(-u)^{n+1}.
 \tag{3.31r}
\]
The disk hypothesis proves directly that
$\Norm(1+u)\ge1/4$ and hence
$\Norm((1+u)^{-1})\le4$.  Consequently
\[
 \Norm\bigl(D_n(u)-(1+u)^{-1}\bigr)\le4^{-n}.
 \tag{3.31s}
\]
The proof constructs the comparison index for every positive rational
squared tolerance, so the derivative approximation is operationally
equivalent to the exact displayed reciprocal; it does not appeal to an
infinite geometric series.  The programs $L_n$ and $Q_n$ separately carry
the same kind of productive CCL sidecar.

For an anchored branch write $x=a(1+u)$ with $a\ne0$.  The correction that
selects the logarithm branch is constant with respect to $x$, and therefore
has zero first coefficient.  Explicit inverse cancellation gives
\[
 (1+u)^{-1}a^{-1}=\bigl(a(1+u)\bigr)^{-1}=x^{-1}.
 \tag{3.31t}
\]
This transports the local result to the branch-aware law
$D\log(x)=x^{-1}$ without choosing a global logarithm.
\end{proof}

\begin{corollary}[finite differential laws for exponential and logarithm]
\label{thm:finite-differential-laws}
The exponential and every supplied anchored logarithm chart satisfy their
respective operational differential laws, with exact secant factorizations
and explicit quadratic-error certificates.
\end{corollary}

\begin{proof}
Combine the preceding two theorems.
\end{proof}

\leanref{CCL.ExponentialCalculus.expDifferential_equivalent_direct}
\leanref{CCL.LogarithmicCalculus.anchoredLogDerivative_inverse}

The preceding differential laws extend through finite syntax, coefficient
tables, and path refinements.  We state each layer separately so that its
data, hypotheses, and computational conclusion remain visible.

\begin{definition}[finite first jet]
A first jet is a pair $(v,d)\in\Qi^2$.  Its product, reciprocal, and
composition operations are
\[
 (v,d)(w,e)=(vw,dw+ve),
 \qquad
 (v,d)^{-1}=\bigl(v^{-1},-d(v^{-1})^2\bigr),
 \qquad
 (w,e)\circ(v,d)=(w,de).
 \tag{3.31u}
\]
\end{definition}

\begin{theorem}[finite first-jet evaluation]
Structural recursion evaluates finite expression syntax containing constants,
one input, arithmetic, reciprocal, local exponential, and local logarithm in
one structural pass, returning its displayed value and first coefficient.
Jet evaluation leaves the ordinary finite-depth value unchanged and proves
the coefficient rule at every node.  Required local-domain conditions are
decidable predicates on displayed rational subexpression values.
\end{theorem}

\begin{proof}
Induct on the finite expression syntax.  Constants and the input are
immediate.  Equation~(3.31u) handles products, reciprocals, and composition;
the two differential theorems supply the exponential and guarded logarithm
nodes.  Each induction branch returns the ordinary value equation together
with its first-coefficient equation.
\end{proof}

\begin{definition}[CCL sine and cosine]
For every coherent arithmetic CCL input $Z$, set
in the already constructed total exponential and field operations
\[
 \sin(Z)={\operatorname{Exp}(iZ)-\operatorname{Exp}(-iZ)\over2i},\qquad
 \cos(Z)={\operatorname{Exp}(iZ)+\operatorname{Exp}(-iZ)\over2}.
 \tag{3.31v}
\]
\end{definition}

\begin{theorem}[sine and cosine are CCL-computable]
Both domains are exactly the type of coherent CCL names; no bounded chart is
required.  Each output answers a requested positive rational squared error
and emits exactly the requested finite bit count with the same-state parser
certificate.
\end{theorem}

\begin{proof}
Scalar multiplication, negation, addition, subtraction, and division by the
displayed non-zero rational constants are CCL-computable.  Compose them with
the total exponential theorem in the two finite trees of (3.31v).
\end{proof}

\begin{theorem}[synchronized progressive first jets]
A synchronized progressive first jet has value and slope
two coherent CCL names accompanied by finite indexed provenance: constants,
the independent input, addition, negation, product, reciprocal, quotient,
composition, total exponential, anchored logarithm, sine, and cosine are the
only constructors, and each constructor computes both result channels by its
displayed rule.  In particular, an arbitrary pair of unrelated coherent names
cannot be inserted as a derivative.  Given a target $\tau>0$, the program
forms one positive rational harmonic request strictly below $\tau$ and below
both channels' coherence moduli.  The resulting two enclosure radii and two
semantic allowances are all below $\tau$, and both cursors return the same
requested finite number of bits.  This supplies a finite, bit-producing
forward-differentiation program over arbitrary coherent progressive inputs;
the quantitative change law is stated separately rather than inferred from
the two channels alone.
\end{theorem}

\begin{proof}
Induct through the permitted provenance constructors and apply the first-jet
rule at each node.  Fold a positive harmonic request through the two finite
coherence moduli.  The resulting radii and allowances are below the target;
the persistent cursors supply the exact bit counts and parser certificates.
\end{proof}

\begin{definition}[quadratic change certificate]
A finite certificate at increment $h$ stores displayed
$b,d,o,r,F$ and proves
\[
 o-b-hd=r,\qquad \Norm(r)\le F\Norm(h)^2.
 \tag{3.31w}
\]
\end{definition}

\begin{proposition}[finite quadratic-change calculus]
Affine, sum, negation, scalar, product, apartness-guarded reciprocal, and
quotient rules compute new tuples directly.  For product and reciprocal the
factor is the exact displayed rational ratio
$\Norm(r)/\Norm(h)^2$ when $h\ne0$, with the zero case discharged by the
algebraic expansion before any division.  On $\Norm(h)\le1$, one certificate
has the explicit displacement sensitivity
$S=2(\Norm(d)+F)$.  If an outer certificate is supplied at the inner displayed
displacement, finite substitution gives the composite factor
\[
 F_{\rm comp}=2\bigl(\Norm(d_{\rm out})F_{\rm in}
              +F_{\rm out}S_{\rm in}^{2}\bigr).
 \tag{3.31x}
\]
The exponential and local-log bounds $10$ and $4$ proved above instantiate
this same record.  Consequently a finite composite retains an explicit
rational quadratic bound, rather than only an informal big-O annotation at
its leaves.
\end{proposition}

\begin{proof}
Expand (3.31w) for each arithmetic constructor and collect its exact linear
and residual parts.  When $h\ne0$, the displayed rational ratio
$\Norm(r)/\Norm(h)^2$ supplies the product or reciprocal factor; at $h=0$ the
algebraic identity discharges the case before division.  For composition,
substitute the inner identity into the outer one and apply the squared
triangle estimate, obtaining (3.31x).  The exponential and logarithm
differential theorems supply factors $10$ and $4$.
\end{proof}

\begin{definition}[progressive quadratic trace]
A progressive quadratic trace supplies a certificate
$C_n=(h_n,b_n,o_n,d_n,r_n,F_n)$ at each index, one displayed geometric scale,
finite later-closeness bounds for all five varying fields, and one rational
bound for their centres.
\end{definition}

\begin{theorem}[persistent bitstreams for a quadratic trace]
The trace data construct coherent channels for
$h,b,o,d,r$ and a computable joint request index.  Each channel is then rooted
at that exact process state and retained as an actual persistent CCL cursor.
For every rational target $\tau>0$ and finite bit count, one synchronized
request returns five enclosure-and-cursor records: every root radius is below
$\tau$, every cursor extends its existing raw word by exactly the requested
number of bits, its projective cylinder has the proved pairwise
Hermitian-projective-quadrance bound,
and all five cursors retain exact ordinary-parser semantics.  The five root
centres are literally the fields of $C_n$, and hence prove
\[
 o_n-b_n-h_nd_n=r_n,\qquad \Norm(r_n)\le F_n\Norm(h_n)^2.
 \tag{3.31w'}
\]
This is not a tuple of rational facts adjacent to the CCL calculation: its
same-index finite observations are the roots of the live bit-emitting cursors.
The constant-trace constructor recovers every exact terminating rational
certificate.  A second executable witness uses the changing finite
exponential approximation at the displayed input $1$ as $h_n$, proves
$h_0\ne h_1$ by finite rational calculation, and nevertheless satisfies the
same productive bitstream-and-quadratic contract.  It therefore exercises the
nonterminating-input route rather than silently replacing it by an exact
rational value.
\end{theorem}

\begin{proof}
Fold one request through the five finite later-closeness moduli and root every
cursor at the resulting common state.  The trace certificate at that state
is exactly (3.31w').  The persistent cursor theorem supplies prefix extension,
length, Hermitian-projective-quadrance, and parser certificates for every
channel.  The constant
trace is immediate.  The changing exponential witness is checked by direct
rational evaluation at its first two depths and then uses the same generic
request construction.
\end{proof}

\begin{proposition}[finite channel transport algebra]
If channels $a_n,b_n$ have quarter-power scales $E_a,E_b$ and centre bounds
$B_a,B_b$, their product receives the displayed scale and bound
\[
 2(E_aB_b+B_aE_b),\qquad B_aB_b.
 \tag{3.31w-tr}
\]
Addition, negation, scalar multiplication, multiplication, and semantic
composition therefore construct the complete common trace without a
caller-supplied continuity record.  A reciprocal does require information
which pointwise nonzeroness cannot constructively provide: displayed
$\lambda>0$ with $\lambda\le\Norm(a_n)$ at every sampled stage.  Its channel
then has scale $E_a\lambda^{-2}$ and bound $\lambda^{-1}$, by the exact
difference-of-inverses identity.  This gives automatic reciprocal and
division traces from explicit base and output lower quadrances.
\end{proposition}

\begin{proof}
Apply the squared triangle estimate to the exact product difference and use
the two centre bounds, giving (3.31w-tr).  Addition, negation, and scaling are
special cases.  For reciprocals use
$a^{-1}-b^{-1}=(b-a)/(ab)$ and the displayed lower-quadrance certificates.
Division is multiplication by the resulting inverse channel.
\end{proof}

\begin{theorem}[analytic primitives preserve quadratic traces]
Simultaneous finite
Horner-depth and input refinement gives scale $2(10+16E)$ for the local
exponential channel.  The reduced complex-log chart
$2z(1+z^2/3+\cdots)$---where the dots describe only the displayed finite
recurrence---uses its proved finite input sensitivity and obtains scale
$2(5+64E)$ on $\Norm(z)\le1/9$.  Each becomes an outer quadratic trace and is
then composed with the inner trace.  Local sine and cosine are assembled from
the two transported exponential traces and exact Gaussian-rational scaling.
Thus arithmetic, composition, exponential, reduced-chart logarithm, sine,
and cosine preserve the semantic certificate as productive CCL programs.
The global range-reduced exponential and an anchored log correction remain
the already proved bit-producing outer layers; their local differential
certificate is the trace just described.
\end{theorem}

\begin{proof}
Use the finite Horner tail and input-sensitivity bounds for exponential and
the reduced logarithm, then insert the resulting outer certificate into
(3.31x).  Sine and cosine follow from their exponential definitions and the
finite channel transport algebra.
\end{proof}

\begin{proposition}[finite failure certificates]
Partial constructors return a finite decision.  Rejection names the
missing positive denominator or output lower quadrance, input disk,
safe chart, overlap, or branch-alignment certificate.  It no longer asks the
caller to invent a generic quadratic-transport proof for these automatic
constructors.
\end{proposition}

\begin{proof}
Each partial constructor compares its finite rational domain certificate.
Success invokes the relevant transport theorem; failure returns the
particular missing positive bound or chart field.
\end{proof}

\begin{theorem}[productive finite trace programs]
These nodes assemble into one finite program indexed by its complete
semantic trace.  A source node carries a quadratic trace; every
other node computes its result index from the traces of its children.
Composition is genuinely dependent: the outer program is indexed by the
displayed inner displacement channel, not by an unrelated input sequence.
The program therefore cannot attach arbitrary value, slope, and remainder
names after the calculation.  Its public request theorem simultaneously
returns the quadratic law, extension and exact new length of all five live
raw prefixes, exact ordinary-parser semantics, and a pairwise
Hermitian-projective-quadrance cylinder bound below the requested rational
target.  Exact
terminating certificates specialize literally, while the changing
exponential-sample witness remains changing after nontrivial arithmetic and
continues to emit every requested finite number of bits.
\end{theorem}

\begin{proof}
Induct over the indexed program.  Every constructor computes its result trace
from its children, and dependent composition uses the actual inner
displacement channel.  Apply the persistent quadratic-trace theorem at the
root.  Exact inputs reduce to terminating traces; the changing witness is
preserved by the same constructor equations.
\end{proof}

\begin{definition}[finite Taylor and multivariable jets]
A Taylor jet is a
low-to-high rational coefficient list, with finite truncation, Cauchy product,
composition, and natural-number iteration of the formal derivative.  A
multivariable first jet
is one Gaussian-rational value and one displayed finite list of partial
coefficients; its dimension is that list's length.
\end{definition}

\begin{proposition}[finite jet table algebra]
Taylor-jet truncation, Cauchy product, composition, and formal differentiation
are finite list recursions.  The product rule holds coefficient by
coefficient.  Directional coefficients of a multivariable jet are finite dot
products and obey the addition, product, and chain rules.
\end{proposition}

\begin{proof}
Induct over the displayed coefficient lists.  Cauchy-product coefficients
give the product rule, Horner recursion gives composition, and dot-product
linearity gives the directional laws.  Every recurrence terminates at a
displayed list length.
\end{proof}

\begin{theorem}[productive progressive coefficient tables]
The progressive counterpart replaces each displayed coefficient by a
coherent CCL channel.  Addition, negation, coherent scaling, finite Cauchy
product, finite Horner composition, truncation, and formal differentiation
all recurse over those displayed lists.  A single rational request is folded
through the finitely many channel moduli.  The table now retains the actual
cursor belonging to every entry, rather than reconstructing a fresh cursor
for each query.  For every retained coefficient one joint request proves a
radius and semantic allowance below the visible target, extension of the
previous prefix by the exact requested bit count, and exact parser semantics.
A finite schedule of differently sized requests preserves those origins and
produces compatible names even when redundant-sector choices make the raw
words textually different.  Truncation transports precisely this contract for
every retained entry; it does not leave an implicit obligation for an
unmentioned tail.  Exact rational coefficient tables specialize literally:
after any positive joint request, their returned centres are the original
displayed coefficient list.
\end{theorem}

\begin{proof}
Fold one positive rational request through the finite list of channel moduli.
Apply the persistent cursor theorem to every retained entry at the common
state.  Induction over the requested schedule preserves each cursor origin;
truncation simply restricts the finite family.  For exact tables every
channel has allowance zero, so the returned centres are definitionally the
input coefficients.
\end{proof}

\begin{theorem}[scheduled compatibility of equivalent tables]
Compatibility also holds when the two origins are merely operationally
equivalent and their finite schedules are different.  After each schedule,
one common final request produces two extending prefixes of the exactly
accounted lengths.  Both retain exact parser semantics and each denotes a
cylinder of pairwise Hermitian projective quadrance at most $\tau$; the operational
equivalence program additionally computes later rational sample indices whose
separation is strictly below $\tau$.  The raw words need not coincide.  This
contract is applied to congruence of addition, subtraction, multiplication,
and automatic exponential, and to the field laws,
$\operatorname{Exp}(x+y)\simeq
  \operatorname{Exp}(x)\operatorname{Exp}(y)$,
$\operatorname{Exp}(\operatorname{Log}(y;A))\simeq y$, and the explicitly
zero-aligned local reverse law.  Hence these laws now govern actual emitted
CCL prefixes rather than stopping at rational sample identities.
\end{theorem}

\begin{proof}
After the two finite schedules, compute a common request using operational
equivalence and the two table moduli.  Each persistent cursor supplies its
own extension, length, diameter, and parser certificate, while operational
equality supplies later rational samples closer than the target.  Apply this
construction entrywise and use the already proved arithmetic, exponential,
and logarithm congruence laws.
\end{proof}

\begin{lemma}[finite coefficient evaluation budget]
If
$a_j$ and $\widehat a_j$ are the actual and displayed entries with
$\Norm(a_j-\widehat a_j)\le\eta$, and $\Norm(x)\le X$, a finite Horner
recurrence computes a rational number $E(m,\eta,X)$ and proves
\[
 \Norm\!\left(\sum_{j<m}a_jx^j-
                    \sum_{j<m}\widehat a_jx^j\right)
 \le E(m,\eta,X).
 \tag{3.31w''}
\]
The corresponding progressive theorem obtains the displayed centres from one
actual table request.  Appending a finite discarded coefficient list gives an
exact algebraic retained-plus-tail identity and charges a separately displayed
finite tail budget; there is no implicit infinite remainder.  The same bound
is exposed for finite Cauchy multiplication and Horner composition.
\end{lemma}

\begin{proof}
Induct through the finite Horner recurrence, applying the product and sum
error bounds at each coefficient.  The progressive statement takes all
centres from one table request.  Appending the explicitly discarded finite
list proves the retained-plus-tail identity; the same induction handles
Cauchy multiplication and composition.
\end{proof}

\begin{theorem}[productive multivariable directional channels]
The same construction handles a natural-number-sized multivariable first
jet.  Its partial channels form one finite table, and a directional channel is
a finite dot-product recursion against a displayed Gaussian-rational list
whose length must equal that table's displayed dimension.  This directional
output again has the complete finite CCL request contract.  A second finite
recurrence transports the displayed coefficient errors through the supplied
direction list and proves the resulting rational directional-error budget.
Thus ``all
coefficients'' here always means membership in one explicit list, never an
unbounded coordinate family.
\end{theorem}

\begin{proof}
Check equality between the direction-list length and the displayed jet
dimension, then recurse through the finite dot product.  A second recursion
propagates the coefficient-error allowances through the same direction list.
Progressive arithmetic gives the complete CCL request contract for the
resulting channel.
\end{proof}

\begin{definition}[finite polynomial antiderivative and integral]
For $p(T)=\sum_{k=0}^n a_kT^k$, define
\[
 A_p(T)=c+\sum_{k=0}^n {a_k\over k+1}T^{k+1},\qquad
 I_p(a,b)=A_p(b)-A_p(a).
\]
\end{definition}

\begin{theorem}[finite polynomial fundamental theorem]
Coefficient arithmetic proves literally that $A_p'=p$.  Applying the already
constructed divided difference before diagonal substitution gives the
infinitesimal-free fundamental theorem
\[
 I_p(a,b)=(b-a)\Delta A_p(b,a),\qquad
 \Delta A_p(a,a)=p(a).
 \tag{3.31y}
\]
This immediately constructs the exact finite initial-value solution of
$y'=p(x)$ for displayed rational initial time and value.  For $y'=y$, the
degree-$(n+1)$ exponential polynomial has derivative equal to the degree-$n$
polynomial, and its exact pointwise defect is the single last displayed
Taylor monomial.  These are deliberately finite ODE certificates; improving
or subdividing them means requesting a larger natural degree, not completing
an infinite trajectory.
\end{theorem}

\begin{proof}
Differentiate the displayed coefficient list to obtain $A_p'=p$.  Apply the
polynomial divided-difference identity before diagonal specialization,
giving (3.31y).  The ODE claims follow by the same finite coefficient
calculation; the exponential defect is its single last displayed monomial.
\end{proof}

\begin{definition}[finite path refinement]
At refinement $n$ the caller
returns one finite list of rational path cells.  Cell $j$ has displayed
endpoints $s_j,t_j$, displacement $d_j=t_j-s_j$, a finite rational
polynomial $p_j$, and a rational local squared-error allowance
$\eta_j\geq0$.  Its displayed integral sample is
\[
 \widehat I_j=d_j\int_0^1p_j(u)\,du,
 \qquad
 \widehat I_n=\sum_j\widehat I_j,
 \tag{3.31y'}
\]
where every integral on the right is the coefficient calculation above.  A
finite observation with error $e_j$ satisfies
$N(e_j)\leq\eta_j$, so multiplicativity of Gaussian quadrance gives the
explicit transported bound
\[
 N(d_je_j)\leq N(d_j)\eta_j.
 \tag{3.31y''}
\]
\end{definition}

\begin{lemma}[finite weighted path-error bound]
Repeated use of $N(u+v)\leq2(N(u)+N(v))$ computes a deliberately coarse
rational allowance $B_n$ by a fold over this particular finite cell list.
Induction over that finite list, using (3.31y'') at the head and the
two-term squared-triangle estimate at each cons step, bounds the quadrance of
the corresponding weighted finite error sum by $B_n$.  There is no
quantification over every point of a completed interval.
\end{lemma}

\begin{proof}
Induct over the displayed cell list.  At the head use (3.31y''); combine it
with the induction hypothesis using the two-term squared triangle estimate.
The resulting recurrence is the computable fold defining $B_n$.
\end{proof}

\begin{definition}[certified integral refinement program]
A refinement program supplies a rational scale $\sigma\geq0$
and checks, for each finite call,
\[
 B_n\leq \sigma4^{-(n+1)},\qquad
 N(\widehat I_m-\widehat I_n)\leq\sigma4^{-(n+1)}\quad(n\leq m).
 \tag{3.31y'''}
\]
\end{definition}

\begin{theorem}[a certified path integral is CCL-computable]
The already proved finite search for the first geometric radius below a
positive rational target is therefore a computable refinement modulus.
It turns the samples $\widehat I_n$ into a coherent approximation and then
into an ordinary persistent CCL cursor.  For every displayed positive target
and natural bit count, one finite request returns an enclosure below that
target, a prefix of exactly that length, prefix persistence, the exact bit
count, and exact parser semantics.  Thus the integral is a CCL computation
with a rational accuracy side channel, not a postulated completed value.
\end{theorem}

\begin{proof}
The quarter-power scan computes a refinement index from the requested target.
Equation~(3.31y''') gives both later-sample and local-error bounds, so the
samples form a coherent progressive enclosure process.  Feed that process to
the persistent CCL controller and request the displayed bit count at the same
refined state.
\end{proof}

\begin{proposition}[finite path-refinement algebra]
All expected finite algebra is recorded at this operational boundary.
Reversing a cell uses an explicit polynomial reparameterization certificate
and negates its sample; a structural path reversal reverses both every cell
and their finite list order.  Adjacent path lists concatenate exactly;
polynomial addition and rational scaling are linear.  These operations now
lift to refinement programs themselves: scaling computes the squared scale
transport, addition computes the two-source geometric modulus, concatenation
uses the literal appended path, and reversal reuses the original modulus
after the proved sign identity.  Where reversing order or reconciling two
local partitions changes the deliberately coarse allowance fold, the caller
supplies that finite rational inequality explicitly.  Each resulting program
therefore inherits the same productive CCL request theorem.
\end{proposition}

\begin{proof}
Cell reversal is the displayed polynomial reparameterization and sign
identity; list reversal and concatenation are structural recursions.
Linearity holds coefficientwise.  Compose the corresponding finite
refinement moduli, using a supplied coarse-fold comparison when partitions
change.
\end{proof}

\begin{theorem}[refinement-program fundamental theorem]
When the local models are exact,
the refinement output is operationally equivalent to the terminating result
of the coefficient integrator.  Finally, a supplied finite antiderivative
process records at every refinement both a literal coefficient derivative
equality and the endpoint-difference equality.  Summing those displayed
differences gives the same sample at every index, hence an operationally
equivalent productive CCL output.  This is the finite fundamental theorem for
the general refinement interface: antiderivatives are certified programs,
not functions represented as infinite sets, and no infinitesimal is used.
\end{theorem}

\begin{proof}
For an antiderivative process, sum the literal endpoint differences;
intermediate endpoints cancel in the finite list, so every sample equals the
coefficient integrator's endpoint difference.
\end{proof}

\begin{definition}[zero-aligned logarithm selection data]
For a candidate logarithm output $L$ over $\operatorname{Exp}(X)$, the
selection data consist of: a rational point admitted by the source and target
enclosures; positive denominator-quadrance and reduced-disk certificates; the
source chart factor and correction; a displayed cell $(c,\rho)$ with
$0<\rho\le1$ containing both $X$ and $L$; and a finite local
exponential-injectivity certificate for that pair in the cell.
\end{definition}

\begin{theorem}[selected logarithm--exponential inverse]
Given zero-aligned selection data, the constructed chart over the actual CCL
$\operatorname{Exp}(X)$ satisfies $L-X\simeq0$ and
\[
 \operatorname{Log}(\operatorname{Exp}(X);\text{zero-aligned chart})\simeq X.
 \tag{3.31z}
\]
Merely knowing $\operatorname{Exp}(L-X)\simeq1$ is intentionally
insufficient.
\end{theorem}

\begin{proof}
The safe-overlap record constructs the anchored chart over the actual
$\operatorname{Exp}(X)$ process.  Its anchored inverse gives
$\operatorname{Exp}(L)\simeq\operatorname{Exp}(X)$.  Apply the supplied local
injectivity certificate to the two displayed cell members, obtaining
$L\simeq X$.  Arithmetic congruence gives $L-X\simeq0$, and additive
cancellation yields (3.31z).
\end{proof}

\begin{corollary}[productive finite calculus]
\label{thm:productive-finite-calculus}
Finite first jets, coefficient tables, multivariable directional jets,
quadratic traces, and certified path integrals are closed under their stated
finite operations.  Exponential, guarded logarithm, sine, and cosine preserve
the corresponding CCL contracts; the two finite fundamental theorems and the
selected reverse logarithm law hold under their displayed certificates.
\end{corollary}

\begin{proof}
Combine the focused results of this subsection.
\end{proof}

\leanref{CCL.ProgressiveTraceProgram.Calculation.persistent_bitstream_and_quadratic_law}
\leanref{CCL.ProgressiveIntegration.RefinementProgram.AntiderivativeProcess.ftc_equivalent}
\leanref{CCL.LogBranchGeometry.SelectedLogarithm.logarithm_exponential_inverse}

\subsection{Fractional powers are finite root families}

Complex $q$th root is a finite-valued operation, not a single-valued one.  We
separate its RFAP family, chart-selected member, exact rational candidates,
and cost account.

\begin{definition}[fractional-power polynomial]
For displayed $y\in\Qi$ and positive $q\in\N$, put
\[
 P_{y,q}(T)=T^q-y.
 \tag{3.32}
\]
\end{definition}

\begin{theorem}[RFAP fractional-power family]
The polynomial $P_{y,q}$ has degree exactly $q$, satisfies
\[
 P_{y,q}(z)=z^q-y
 \tag{3.33}
\]
for every displayed $z\in\Qi$, and its RFAP root family has
multiplicity-weighted occurrence count exactly $q$.
\end{theorem}

\begin{proof}
The coefficient list is explicit, so its leading coefficient and degree are
computed directly; finite Horner induction proves (3.33).  The
squarefree/multiplicity
preprocessor followed by the RFAP box construction therefore returns a list
of exact linear terminals and independently shrinking higher-degree CCL
families whose multiplicity-weighted occurrence count is exactly $q$.
For $y=0$ this correctly permits the single distinct site $0$ with
multiplicity $q$; no artificial list of $q$ unequal zeroes is manufactured.
\end{proof}

\begin{definition}[chart-selected fractional power]
For anchored chart $A$, define
\[
 \mathsf{SelectedRoot}(Y;A,q)
   =\mathsf{ExpName}\bigl(q^{-1}\mathsf{LogName}(Y;A)\bigr),
 \tag{3.34}
\]
\end{definition}

\begin{theorem}[selected root reconstruction]
The selected root satisfies
\[
 \mathsf{SelectedRoot}(Y;A,q)^q\simeq Y.
 \tag{3.35}
\]
\end{theorem}

\begin{proof}
Finite induction applies the exponential addition law through the displayed
$q$-fold CCL product.  The rational exponent cancels, and the anchored
exponential--logarithm inverse returns $Y$.
\end{proof}

The polynomial family, rather than $\exp(\log(y)/q)$, is the multi-valued
operation $y^{1/q}$.  The selected root is only one chart-selected member of
(3.32).  Likewise, the branch-selected rational power $Y^{p/q}$ is the
displayed $p$-fold product of (3.34).

\begin{definition}[sparse rational root candidate system]
Define the radius polynomial
\[
 R_{y,q}(S)=S^q-N(y).
 \tag{3.36}
\]
For a displayed rational $s$, define the finite Lucas polynomials
\[
 A_0(X;s)=2,\qquad A_1(X;s)=2X,\qquad
 A_{n+2}(X;s)=2X A_{n+1}(X;s)-sA_n(X;s).
 \tag{3.37}
\]
For each surviving rational radius $s$, use the coordinate polynomials
$A_q(X;s)-2\Re(y)$ and $A_q(X;s)-2\Re((-i)^qy)$.
\end{definition}

\begin{lemma}[root projection into the sparse system]
If $z^q=y$ and $s=N(z)$, then $R_{y,q}(s)=0$ and
\[
 A_q(\Re z;N(z))=2\Re(z^q).
 \tag{3.38}
\]
and
\[
 A_q(\Im z;s)=2\Re((-i)^qy).
 \tag{3.39}
\]
\end{lemma}

\begin{proof}
Multiplicativity gives $s^q=N(y)$ and hence the radius equation.  Adding the
$n$th powers of $z$ and $\bar z$ proves (3.38) by finite induction through
(3.37).  Apply the same identity to $-iz$, whose real coordinate is $\Im z$,
quadrance is $s$, and $q$th power is $(-i)^qy$.
\end{proof}

Each searched polynomial has $q+1$ dense coefficient slots and leading
coefficient respectively $1,2^q,2^q$.  Thus this path constructs no
Sylvester matrix, determinant, scaled inverse, or resultant coefficient
tower.

\begin{lemma}[finite rational-root divisor pool]
For any one of these displayed rational
polynomials, trailing zero coefficients are removed, its finite canonical
denominators are cleared, and its nonzero leading and constant integer
coefficients are denoted $L$ and $C$.  If a nonzero root is in canonical
reduced form $a/b$, homogeneous evaluation and $\gcd(|a|,b)=1$ give directly
\[
 |a|\mid |C|,\qquad b\mid |L|.
 \tag{3.40}
\]
The executable pool tests zero separately and otherwise crosses the signed
positive divisors of $|C|$ with the positive divisors of $|L|$.  In
particular it inspects no fraction merely because its numerator and
denominator happen to lie below a common magnitude bound.  If
\[
 D(P)=1+2|C(P)|\,|L(P)|,
 \tag{3.41}
\]
then the exact filtered candidate list for $P$ has length at most $D(P)$.
\end{lemma}

\begin{proof}
Clear the finite canonical denominators and homogenize at the reduced
fraction $a/b$.  Coprimality makes $a$ divide the constant endpoint and $b$
divide the leading endpoint, proving (3.40).  The pool has at most one zero
and $2|C||L|$ signed divisor pairs, proving (3.41).
\end{proof}

\begin{theorem}[complete automatic rational-root candidate search]
Let $\mathcal R$ be the
finite filtered candidate list for $R_{y,q}$, and write $X_s,Y_s$ for the two
coordinate polynomials at $s$.  At most
\[
 N(y,q)=\sum_{s\in\mathcal R}D(X_s)D(Y_s)
 \tag{3.42}
\]
Gaussian candidates reach duplicate removal and the exact power pass.  The
sum here is a fold over the displayed finite list, not a sum over a named
point collection.  The final pass performs at most $N(y,q)$ power tests and
the simple deduplicator at most $N(y,q)^2$ rational-pair comparisons.  Every
displayed Gaussian-rational solution of $z^q=y$ occurs in the promoted list
and is returned as a terminating CCL name.  If that distinct list has $q$
entries, the degree-sharp check certifies it as complete.
\end{theorem}

\begin{proof}
For a supplied rational root $z$, the sparse-projection lemma puts its norm
and coordinates into the three polynomial root lists.  The divisor-pool
lemma ensures that each reduced rational occurs in its finite pool, and exact
Horner evaluation retains it.  Thus the Cartesian product contains $z$,
deduplication retains one copy, and the exact power test promotes it.  The
product of the coordinate bounds and the finite fold over $\mathcal R$ give
(3.42); nested-loop counting gives the Horner, power, and deduplication
bounds.  A successful test supplies the terminating rational CCL certificate,
and the degree-sharp conclusion follows by finite synthetic division.
\end{proof}

For example, the checked batch $[2,-2]$ for $y=4,q=2$ is returned without a
caller-supplied candidate list.  The search quantifies over no set of all
Gaussian rationals and never tries to recognize an exact rational from an
opaque non-terminating prefix.

\begin{definition}[defensive bit-work ledger]
Put
\[
 \beta(n)=1+\lfloor\log_2 n\rfloor,\qquad
 \beta(a/b)=1+\beta(|a|)+\beta(b),\qquad
 \beta(u+iv)=\beta(u)+\beta(v),
 \tag{3.43}
\]
with the displayed one-bit conventions at zero.  In the defensive
schoolbook model, a rational arithmetic operation with inputs of widths
$r_1,r_2$ and normalized output width $r_3$ is charged
\[
 (r_1+r_2+r_3+1)^2
 \quad\hbox{core bit work},\qquad
 (r_3+1)^3
 \quad\hbox{canonical-normalization work}.
 \tag{3.44}
\]
The cubic term is a deliberately coarse allowance for the finite Euclidean
normalization; it is kept separate rather than hidden inside a unit-cost
field operation.  The formal ledger follows the actual intermediate values
of every Horner multiplication/addition and every repeated Gaussian power,
the actual recursive denominator-clearance lists, every finite divisor scan,
and the complete defensive duplicate-comparison pass.  The sparse
construction materializes exactly
\[
 (1+2|\mathcal R|)(q+1)
 \tag{3.45}
\]
coefficient slots.  Finally
\[
 B(y,q)=B_{\rm clear}+B_{\rm sparse}+B_{\rm divisors}
       +B_{\rm normalize}+B_{\rm Horner}+B_{\rm dedup}+B_{\rm power}
 \tag{3.46}
\]
\end{definition}

\begin{theorem}[defensive exact-root search bound]
The complete automatic rational-root search performs at most $B(y,q)$ bit
work in the schoolbook model (3.44), and materializes exactly the number of
coefficient slots in (3.45).
\end{theorem}

\begin{proof}
The ledger is a fully executable natural number computed from the displayed input and
all finite intermediate lists.  Summing the displayed charges for denominator
clearance, sparse coefficient construction, divisor enumeration, Horner
evaluation, duplicate comparisons, and exact power tests gives (3.46),
together with the Horner, comparison, power-test, and final-output bounds.  It is an upper
account in the stated schoolbook model, not a claim about wall-clock time or
an asymptotically optimal program.  Faster integer
multiplication and gcd, factored divisor generation, and a balanced finite
deduplication table are now explicit places to sharpen (3.44)--(3.46).
\end{proof}

\begin{corollary}[finite family of fractional powers]
\label{thm:finite-root-family}
For displayed $y\in\Qi$ and positive $q$, $T^q-y$ has an RFAP family of
multiplicity-weighted size $q$; every Gaussian-rational root is found and
returned as a terminating CCL name within the candidate and bit-work bounds;
and every anchored logarithm chart supplies one productive selected root
whose $q$th power is operationally $y$.
\end{corollary}

\begin{proof}
Combine the RFAP-family, selected-root, candidate-search, and work-ledger
theorems.
\end{proof}

\leanref{CCL.RootPowers.RootFamily.occurrenceCount_eq_degree}
\leanref{CCL.RootPowers.exactRoot_mem_automaticExactRoots}
\leanref{CCL.RootPowerBitComplexity.automaticExactRoots_bit_complexity}
\leanref{CCL.RootPowers.selectedRoot_power_reconstructs}

\fi

\ifCCLone
\section{Exact finite examples}

\begin{example}[real examples]
The rational $355/113$ has classical continued fraction
\[
 \frac{355}{113}=[3;7,16]
\]
and continued-logarithm word
\[
 1000100011100111001100101.
\]
The rational $13/5$ has continued-logarithm word
\[
 101001.
\]
\end{example}

\begin{example}[CCF examples]
Algorithm~\ref{alg:ccf} gives
\[
 \frac{3+4i}{5}\rightsquigarrow(0,1-i,1-i,2),
\]
\[
 \frac{13}{5}\rightsquigarrow(2,1+i,1,2+2i),
\]
and
\[
 \frac{355}{113}\rightsquigarrow(3,7+7i,16).
\]
Each list is finite and may be checked by exact Gaussian-rational substitution.
\end{example}

\begin{example}[a short CCL code]
For $z=(3+4i)/5$, Algorithm~\ref{alg:ccl} emits
\[
 00001100010011.
\]
In completed-block notation this is
\[
 0[0,1+i]\,[1,1+i]\,[0,-1+i].
\]
The annulus-stage values are
\[
 \frac{3+4i}{5},\qquad \frac12+i,\qquad -1+i,
\]
so the final sector subtraction is exact.
\end{example}

\begin{example}[a real-axis CCL code]
For $13/5$, the revised boundary convention gives
\[
 110110110111111000110001
\]
and the block expansion
\[
 [2,-i]\,[2,1-i]\,[3,1+i]\,[1,1+i].
\]
The last normalized value is $1+i$, so the stream terminates after its sector code.
\end{example}

\begin{example}[the code for $355/113$]\label{ex:ccl-355}
The exact rational $355/113$ has CCL word
\[
 11101011101001111011111001011111111010111111110111.
\]
Its completed-block parsing is
\[
 [3,-1-i]\,[2,-1]\,[4,1-i]\,[2,i]\,[8,-1-i]\,[7,1-i].
\]
The ten-bit prefix discussed in Section~7 is not a list of two pairs: only the first pair is complete, and the final raw fragment is $110$.
\end{example}

\leanref{CCL.Examples.threePlusFourIOverFive_raw;
CCL.Examples.thirteenOverFive_raw;
CCL.Examples.threeFiftyFiveOverOneThirteen_raw}

\section{Conclusion: Cartesian and polar structure}

The CCF digit $t=\sqfloor z$ records a translate of the coordinate unit square.  Its two coordinates are quantized independently, so the representation is Cartesian and its digit alphabet is unbounded.

The CCL first removes a power of $\gamma=1+i$, whose norm is two, and then chooses one entry from the finite eight-direction table (6.1).  A completed regular block is therefore summarized by
\[
 [k,s]=\text{unary radial exponent }k\;\text{ followed by finite angular representative }s.
\]
This is the precise finite sense in which CCL is polar.  The raw representation remains binary, and incomplete raw fragments remain semantically significant even when no pair $[k,s]$ is complete.

The two algorithms now have fully finite termination arguments.  CCF uses a three-step macro-recursion certified by the finite table (5.1).  CCL uses a one-block macro-recursion with strict Gaussian denominator descent.  Its angular code is the ordinary binary octant index: two bits maintain a finite quarter-turn frame and one bit chooses axial versus diagonal.  Unary and unit folds are exact coordinate substitutions, while every completed block contracts compatible projective tails by the explicit factor $q_r\le1/2$.

The CCL is therefore especially well suited to finite exact computation because normalization, sector selection, incremental bit parsing, matrix decoding, pairwise Hermitian-projective-quadrance bounds, box transport, and the termination measure are all explicit calculations on rational data.  Extending a raw prefix always refines its rational predicate; accumulating completed blocks gives an effective $2^{-n}$ Hermitian-projective-quadrance certificate, and on a bounded initial branch this produces rational boxes of any requested squared diagonal.  The resulting operational name supplies the enclosure layer on which the arithmetic and finite-calculus constructions of CCL-2 build.

\fi

\ifCCLtwo
\section{Conclusion: arithmetic as finite requests}

The canonical CCL of CCL-1 is left unchanged.  Productive arithmetic instead
uses overlapping safe output predicates and carries a rational accuracy
sidechannel beside the ordinary raw prefix.  This extra finite state is what
allows an output to keep producing bits when an opaque computation approaches
a canonical boundary or has an exact result that cannot be recognized from
its input names alone.

Addition, subtraction, multiplication, and apartness-certified division are
closed under this interface.  Bounded range preparation and finite Horner
recurrences extend it to exponential and logarithm, while explicit branch
charts record exactly the information required for complex logarithms.
Operational equivalence then expresses the field and transcendental laws at
arbitrary requested accuracy without replacing the finite processes by
completed values.

The same discipline supports divided differences, jets, Taylor bounds,
coefficient-defined integration, progressive rational-path integration,
trigonometric compositions, rational-function evaluation, and finite root
families.  Every conclusion is witnessed at a displayed finite request.
Sharper complexity constants and broader automatically selected logarithm
charts remain possible improvements, but they do not alter the finite
computational foundation proved here.
\fi

\appendix
\section{Independent machine verification}

The references headed ``Formal verification'' point to the corresponding
declarations in the companion Lean development.  They are not premises of the
paper proofs: each mathematical argument appears in the body above.  The
\texttt{CCL.Paper} and \texttt{CCL.Bridge.Paper} facades collect the combined
CCL-1 and CCL-2 verification surface.  A warning-fatal Lean build, the
generated repository-wide \texttt{check-axioms.sh} audit, and an independent
nanoda check verify that surface.  The reported logical foundation is only
\texttt{propext}; no declaration uses \texttt{Classical.choice},
\texttt{Quot.sound}, a completed scalar field, or a carrier-wide collection
of values.

\begin{thebibliography}{99}

\bibitem{wildberger-projective}
N. J. Wildberger,
\emph{One dimensional metrical geometry},
arXiv:math/0701338, 2007.

\bibitem{hakmem}
M. Beeler, R. W. Gosper, and R. Schroeppel,
\emph{HAKMEM}, MIT Artificial Intelligence Laboratory Memo AIM-239, 1972.

\bibitem{borwein-cl}
J. M. Borwein, N. J. Calkin, S. B. Lindstrom, and A. Mattingly,
Continued logarithms and associated continued fractions,
\emph{Experimental Mathematics} \textbf{26} (2017), no.~4, 412--429.

\bibitem{borwein-generalized}
J. M. Borwein, K. G. Hare, and J. G. Lynch,
Generalized continued logarithms and related continued fractions,
\emph{Journal of Integer Sequences} \textbf{20} (2017), Article 17.5.7.

\bibitem{collins}
M. J. Collins,
Computing with continued logarithms,
\emph{arXiv:2602.07245}, 2026.

\bibitem{hockman}
M. Hockman,
The geometry of Gaussian integer continued fractions,
\emph{Journal of Number Theory} \textbf{197} (2019), 145--167.

\bibitem{hurwitz}
A. Hurwitz,
\"Uber die Entwicklung complexer Gr\"ossen in Kettenbr\"uche,
\emph{Acta Mathematica} \textbf{11} (1887), 187--200.

\bibitem{kirsebom}
M. S. Kirsebom,
Extreme value theory for Hurwitz complex continued fractions,
\emph{Journal of Theoretical Probability} \textbf{34} (2021), 2166--2197.

\bibitem{gonzalez}
G. Gonz\'alez Robert,
\emph{Complex Continued Fractions}, Ph.D. thesis, Aarhus University, 2018.

\bibitem{shallit}
J. Shallit,
Length of the continued logarithm algorithm on rational inputs,
\emph{Integers} \textbf{20} (2020), Paper A69.

\end{thebibliography}

\end{document}
